% Lesson 49 — Listening: Texts about subscribing to service packages (telephone/internet services) — Anglais Tle SH — Cours signé OnBuch+
% Programme officiel MINESEC. Compiler avec : tectonic EN49-listening-texts-about-subscribing-to-service-packages-t.tex
\newcommand{\DOCMATIERE}{Anglais}
\newcommand{\DOCNIVEAU}{Tle SH}
\newcommand{\DOCMODULE}{Chapter 5 : Media and Communication}
\newcommand{\DOCLECON}{EN49}
\newcommand{\DOCTITRE}{Lesson 49 — Listening: Texts about subscribing to service packages (telephone/internet services)}
% OnBuch+ — préambule « cours signés » ENRICHI (compilé avec tectonic / XeLaTeX)
% Système visuel « Pop » d'OnBuch+ : fond crème (jamais blanc), bordures encre
% épaisses, ombres « dures » décalées, titres Archivo Black en capitales.
% Version MATHÉMATIQUES : graphes (pgfplots/tikz), boîtes pédagogiques
% (définition, propriété, méthode, attention, le savais-tu, exercice résolu,
% exercice, TP, bilan), page de garde et signature OnBuch+ ancrée en bas.
%
% Macros attendues dans main.tex AVANT \input{preamble} :
%   \DOCMATIERE  \DOCNIVEAU  \DOCMODULE  \DOCLECON (numéro)  \DOCTITRE
\documentclass[11pt]{article}

\usepackage{fontspec}
\usepackage[english,french]{babel} % français = langue principale ; l'anglais via \eng{..} / textebox
\usepackage[a4paper,margin=2cm,top=3.4cm,bottom=2.8cm,headheight=1.4cm,footskip=1.3cm]{geometry}
\usepackage{amsmath,amssymb,amsfonts}
\usepackage{mathtools} % \xmapsto, \coloneqq... souvent utilisés par les modèles
\usepackage{cancel} % simplifications barrées (bilans de forces)
% \abs{z} (module, valeur absolue) et \norm{u} (norme), utilisés par les modèles
\providecommand{\abs}{}\let\abs\relax\DeclarePairedDelimiter{\abs}{\lvert}{\rvert}
\providecommand{\norm}{}\let\norm\relax\DeclarePairedDelimiter{\norm}{\lVert}{\rVert}
\usepackage[table]{xcolor} % [table] : fournit aussi \rowcolors (alternance de lignes)
\usepackage{pagecolor}
\usepackage{tikz}
\usetikzlibrary{arrows.meta,positioning,calc,shapes.geometric,shapes.misc,shapes.arrows,decorations.pathmorphing,decorations.pathreplacing,decorations.markings,patterns,shadows,mindmap,trees,fit,backgrounds,matrix,intersections,angles,quotes}
\usepackage{pgfplots}
\usepackage[version=4]{mhchem} % équations nucléaires \ce{^{235}_{92}U}
\usepackage[european,siunitx]{circuitikz} % schémas électriques
\pgfplotsset{compat=1.18}
\usepgfplotslibrary{fillbetween,groupplots} % aires entre courbes (calcul intégral)
\usepackage{enumitem}
\usepackage{fancyhdr}
% [most] charge toutes les bibliothèques tcolorbox (listings, minted, raster,
% documentation...) et gonfle inutilement la mémoire TeX sur un document long
% (~300 boîtes) : on ne charge que ce qui est réellement utilisé.
\usepackage[skins,breakable]{tcolorbox}
\usepackage{graphicx}
\usepackage{colortbl}
\usepackage{booktabs}
\usepackage{tabularx}
\usepackage{diagbox} % case d'en-tête diagonale des tableaux à double entrée
% Colonnes extensibles centrée / gauche / droite (C, L, R), souvent utilisées par les modèles
\newcolumntype{C}{>{\centering\arraybackslash}X}
\newcolumntype{L}{>{\raggedright\arraybackslash}X}
\newcolumntype{R}{>{\raggedleft\arraybackslash}X}
\usepackage{array}
\usepackage{multicol}
\usepackage{multirow}
\usepackage{siunitx}
% parse-numbers=false : accepte aussi \SI{2,5 \times 10^{-3}}{...}
\sisetup{locale=FR,output-decimal-marker={,},group-minimum-digits=4,parse-numbers=false}
\DeclareSIUnit\ohm{\ensuremath{\symup{\Omega}}} % U+2126 absent de la police
\DeclareSIUnit\division{div} % réglages d'oscilloscope (ms/div, V/div)
\DeclareSIUnit\tr{tr} % tours (vitesse de rotation, ex. \SI{1500}{\tr\per\minute})
\DeclareSIUnit\molar{M} % concentration molaire (ex. \SI{3}{\molar}, \SI{1}{\micro\molar})
\DeclareSIUnit\an{an} % année (le modèle confond parfois avec \year, un compteur TeX) — ex. \SI{5}{\kilo\meter\per\an}
\DeclareSIUnit\FCFA{FCFA} % franc CFA (ex. \SI{500}{\FCFA\per\kilo\gram})
\DeclareSIUnit\degreeBrix{\textdegree Brix} % teneur en sucre d'un sirop/jus (réfractométrie)
\DeclareSIUnit\month{mois} % le modèle utilise parfois \month, un compteur TeX, comme unité de durée
\usepackage{qrcode}
% \degree : commande fréquemment utilisée par les modèles pour un angle ou une
% température en dehors de \SI{}{} (non fournie par les paquets chargés).
\providecommand{\degree}{^{\circ}}
% \qed (fin de démonstration), utilisé par les modèles sans amsthm
\providecommand{\qed}{\hfill$\square$}
% \barre : double barre des valeurs interdites dans un tableau de variations (tkz-tab)
\providecommand{\barre}{\ensuremath{\|}}
% environnement proof (amsthm non chargé) utilisé par les modèles
\ifdefined\proof\else\newenvironment{proof}[1][Démonstration]{\par\noindent\textit{#1.}\ }{\qed\par}\fi
\usepackage{lastpage}
\usepackage{hyperref}
\hypersetup{hidelinks}

% ============================================================
% Palette « Pop » OnBuch+ — mêmes hex que la classe P (plus_ui.dart)
% ============================================================
\definecolor{poporange}{HTML}{F59321}
\definecolor{popdark}{HTML}{E07A0C}
\definecolor{popcream}{HTML}{FFF6E8}
\definecolor{popink}{HTML}{1C1714}
\definecolor{poppurple}{HTML}{7A5AE0}
\definecolor{popgreen}{HTML}{2E9E6A}
\definecolor{poppink}{HTML}{E0457B}
\definecolor{popblue}{HTML}{2F7DE1}
\definecolor{popgold}{HTML}{C98A12}
\definecolor{popmuted}{HTML}{8A7F76}
\definecolor{popline}{HTML}{EADFD2}
\definecolor{popgray}{HTML}{8A7F76}
\colorlet{popgrey}{popgray}
% Alias tolérants (le modèle nomme parfois les couleurs différemment)
\colorlet{popwhite}{white}
\colorlet{popblack}{popink}
\colorlet{popred}{poppink}
\colorlet{popyellow}{popgold}
% Teintes claires (fonds de tableaux/encadrés) — le modèle nomme parfois
% les couleurs avec un suffixe L (ex. popblueL) sans les avoir définies.
\colorlet{poporangeL}{poporange!20}
\colorlet{popdarkL}{popdark!20}
\colorlet{poppurpleL}{poppurple!20}
\colorlet{popgreenL}{popgreen!20}
\colorlet{poppinkL}{poppink!20}
\colorlet{popblueL}{popblue!20}
\colorlet{popgoldL}{popgold!20}
\colorlet{popredL}{popred!20}
\colorlet{popgrayL}{popgray!20}
% Teintes claires pour fonds de tableaux / zones
\colorlet{popblueL}{popblue!10!white}
\colorlet{poporangeL}{poporange!14!white}
\colorlet{popgreenL}{popgreen!12!white}
\colorlet{poppurpleL}{poppurple!12!white}
\colorlet{poppinkL}{poppink!12!white}
\colorlet{popgoldL}{popgold!14!white}

\pagecolor{popcream}

% ============================================================
% Typographie
% ============================================================
\defaultfontfeatures{Ligatures=TeX}
\setmainfont{PlusJakartaSans-Medium.ttf}[Path=../../../fonts/,BoldFont=PlusJakartaSans-Bold.ttf]
% Symboles, API et emojis absents de la police principale : repli sur DejaVu Sans (caractères actifs)
\newfontface\symfont{DejaVuSans.ttf}[Path=../../../fonts/]
\makeatletter
\newcommand\symactive[1]{\catcode#1=13 \begingroup\lccode`\~=#1 \lowercase{\endgroup\def~}{{\symfont\char#1}}}
\newcommand\symblank[1]{\catcode#1=13 \begingroup\lccode`\~=#1 \lowercase{\endgroup\def~}{}}
\newcommand\symhyph[1]{\catcode#1=13 \begingroup\lccode`\~=#1 \lowercase{\endgroup\def~}{\mbox{-}}}
\newcount\symc
\newcommand\symrange[2]{\symc=#1 \@whilenum\symc<#2 \do{\expandafter\symactive\expandafter{\the\symc}\advance\symc by 1 }}
\newcommand\symblankrange[2]{\symc=#1 \@whilenum\symc<#2 \do{\expandafter\symblank\expandafter{\the\symc}\advance\symc by 1 }}
\makeatother
\symrange{"0250}{"02FF}\symactive{"2713}\symactive{"2714}\symactive{"2717}\symactive{"2718}\symactive{"2610}\symactive{"2611}\symactive{"2612}
\symhyph{"2011}\symhyph{"2E1A}\symblankrange{"1F300}{"1FAFF}\symblank{"FE0F}
% Maths en sans-serif (Fira Math) avec chiffres et lettres droites de Plus
% Jakarta, pour que les formules se fondent dans le texte.
\usepackage[math-style=ISO,bold-style=ISO]{unicode-math}
\setmathfont{FiraMath-Regular.otf}[Path=../../../fonts/]
\setmathfont{PlusJakartaSans-Medium.ttf}[Path=../../../fonts/,range={up/num,up/{Latin,latin},\mathup}]
\usepackage{icomma} % virgule décimale sans espace en mode math
% Caractères mathématiques Unicode absents des polices de texte (Plus Jakarta,
% Archivo Black) : rendus par leur équivalent mathématique, en texte comme en
% formule — sinon ils s'affichent en carré vide (ex. « Arithmétique dans ℤ »).
\usepackage{newunicodechar}
\newunicodechar{ℝ}{\ensuremath{\mathbb{R}}}
\newunicodechar{ℤ}{\ensuremath{\mathbb{Z}}}
\newunicodechar{ℂ}{\ensuremath{\mathbb{C}}}
\newunicodechar{ℕ}{\ensuremath{\mathbb{N}}}
\newunicodechar{ℚ}{\ensuremath{\mathbb{Q}}}
\newunicodechar{𝔻}{\ensuremath{\mathbb{D}}}
\newunicodechar{α}{\ensuremath{\alpha}}
\newunicodechar{β}{\ensuremath{\beta}}
\newunicodechar{γ}{\ensuremath{\gamma}}
\newunicodechar{δ}{\ensuremath{\delta}}
\newunicodechar{ε}{\ensuremath{\varepsilon}}
\newunicodechar{θ}{\ensuremath{\theta}}
\newunicodechar{λ}{\ensuremath{\lambda}}
\newunicodechar{μ}{\ensuremath{\mu}}
\newunicodechar{σ}{\ensuremath{\sigma}}
\newunicodechar{τ}{\ensuremath{\tau}}
\newunicodechar{φ}{\ensuremath{\varphi}}
\newunicodechar{ψ}{\ensuremath{\psi}}
\newunicodechar{ω}{\ensuremath{\omega}}
\newunicodechar{Σ}{\ensuremath{\Sigma}}
% majuscules produites par \MakeUppercase dans les titres (α-aminés -> Α)
\newunicodechar{Α}{\ensuremath{\alpha}}
\newunicodechar{Β}{\ensuremath{\beta}}
\newunicodechar{Γ}{\ensuremath{\Gamma}}
\newunicodechar{Δ}{\ensuremath{\Delta}}
\newunicodechar{Ε}{\ensuremath{\varepsilon}}
\newunicodechar{Θ}{\ensuremath{\Theta}}
\newunicodechar{Λ}{\ensuremath{\Lambda}}
\newunicodechar{Μ}{\ensuremath{\mu}}
\newunicodechar{Τ}{\ensuremath{\tau}}
\newunicodechar{Φ}{\ensuremath{\Phi}}
\newunicodechar{Ψ}{\ensuremath{\Psi}}
\newunicodechar{Ω}{\ensuremath{\Omega}}
\newunicodechar{↦}{\ensuremath{\mapsto}}
\newunicodechar{∈}{\ensuremath{\in}}
\newunicodechar{∉}{\ensuremath{\notin}}
\newunicodechar{∧}{\ensuremath{\wedge}}
\newunicodechar{⇔}{\ensuremath{\Leftrightarrow}}
\newunicodechar{⟺}{\ensuremath{\Longleftrightarrow}}
\newunicodechar{⇒}{\ensuremath{\Rightarrow}}
\newunicodechar{⟹}{\ensuremath{\Longrightarrow}}
\newunicodechar{∘}{\ensuremath{\circ}}
\newunicodechar{∪}{\ensuremath{\cup}}
\newunicodechar{∩}{\ensuremath{\cap}}
\newunicodechar{≡}{\ensuremath{\equiv}}
\newunicodechar{⊂}{\ensuremath{\subset}}
\newunicodechar{⊥}{\ensuremath{\perp}}
\newunicodechar{∃}{\ensuremath{\exists}}
\newunicodechar{∀}{\ensuremath{\forall}}
\newunicodechar{∛}{\ensuremath{\sqrt[3]{\,}}}
\newunicodechar{∜}{\ensuremath{\sqrt[4]{\,}}}
\newunicodechar{⃗}{\ensuremath{{}^{\to}}}
\newunicodechar{ⁿ}{\ifmmode^{n}\else\textsuperscript{\ensuremath{n}}\fi}
\newunicodechar{ˣ}{\ifmmode^{x}\else\textsuperscript{\ensuremath{x}}\fi}
\newunicodechar{ᵇ}{\ifmmode^{b}\else\textsuperscript{\ensuremath{b}}\fi}
\newunicodechar{ʳ}{\ifmmode^{r}\else\textsuperscript{\ensuremath{r}}\fi}
\newunicodechar{ᵘ}{\ifmmode^{u}\else\textsuperscript{\ensuremath{u}}\fi}
\newunicodechar{ʸ}{\ifmmode^{y}\else\textsuperscript{\ensuremath{y}}\fi}
\newunicodechar{ᵃ}{\ifmmode^{a}\else\textsuperscript{\ensuremath{a}}\fi}
\newunicodechar{ᵉ}{\ifmmode^{e}\else\textsuperscript{\ensuremath{e}}\fi}
\newunicodechar{ᵏ}{\ifmmode^{k}\else\textsuperscript{\ensuremath{k}}\fi}
\newunicodechar{ᵖ}{\ifmmode^{p}\else\textsuperscript{\ensuremath{p}}\fi}
\newunicodechar{ᴺ}{\ifmmode^{N}\else\textsuperscript{\ensuremath{N}}\fi}
\newunicodechar{ᶜ}{\ifmmode^{c}\else\textsuperscript{\ensuremath{c}}\fi}
\newunicodechar{ʰ}{\ifmmode^{h}\else\textsuperscript{\ensuremath{h}}\fi}
\newunicodechar{ᵠ}{\ifmmode^{q}\else\textsuperscript{\ensuremath{q}}\fi}
\newunicodechar{ₙ}{\ifmmode_{n}\else\textsubscript{\ensuremath{n}}\fi}
\newunicodechar{ₐ}{\ifmmode_{a}\else\textsubscript{\ensuremath{a}}\fi}
\newunicodechar{ᵢ}{\ifmmode_{i}\else\textsubscript{\ensuremath{i}}\fi}
\newunicodechar{ᵣ}{\ifmmode_{r}\else\textsubscript{\ensuremath{r}}\fi}
\newunicodechar{ₖ}{\ifmmode_{k}\else\textsubscript{\ensuremath{k}}\fi}
\newunicodechar{ᵦ}{\ifmmode_{\beta}\else\textsubscript{\ensuremath{\beta}}\fi}
\newunicodechar{ₓ}{\ifmmode_{x}\else\textsubscript{\ensuremath{x}}\fi}
\newfontfamily\popdisplay{ArchivoBlack-Regular.ttf}[Path=../../../fonts/]
\newfontfamily\poplogo{SpaceGrotesk-Bold.ttf}[Path=../../../fonts/]
\newfontfamily\popsemi{PlusJakartaSans-SemiBold.ttf}[Path=../../../fonts/]
\newfontfamily\popxbold{PlusJakartaSans-ExtraBold.ttf}[Path=../../../fonts/]
\newfontfamily\popxboldls{PlusJakartaSans-ExtraBold.ttf}[Path=../../../fonts/,LetterSpace=3.0]

\setlist[enumerate]{leftmargin=1.7em,itemsep=5pt,topsep=4pt}
\setlist[itemize]{leftmargin=1.7em,itemsep=4pt,topsep=4pt,label={\color{poporange}\textbullet}}
\setlength{\parindent}{0pt}
\setlength{\parskip}{5pt}
\renewcommand{\arraystretch}{1.25}

% pgfplots : style Pop par défaut
\pgfplotsset{
  every axis/.append style={axis line style={popink,line width=1pt},tick style={popink},
    grid style={popline,line width=0.5pt},label style={font=\small},tick label style={font=\footnotesize}},
  popcurve/.style={poporange,line width=1.8pt,no markers,smooth},
}
% Même style disponible dans un \draw TikZ simple (hors pgfplots)
\tikzset{popcurve/.style={poporange,line width=1.8pt}}
\tikzset{popgrid/.style={popline,very thin}}
\colorlet{popcurve}{poporange} % le nom du style est parfois utilisé comme couleur

% ============================================================
% Pill / squiggle
% ============================================================
\newcommand{\poppill}[3]{%
  \tikz[baseline=-0.55ex]\node[draw=popink,line width=1.2pt,rounded corners=2.6pt,%
    fill=#2,inner xsep=6.5pt,inner ysep=3pt]{%
    \popxboldls\fontsize{7}{7}\selectfont\color{#3}\MakeUppercase{#1}};%
}
\newcommand{\popsquiggle}[1]{%
  \tikz[baseline=-0.4ex]{
    \draw[#1,line width=2.6pt,line cap=round]
      plot [smooth] coordinates {(0,0.09) (0.34,0.01) (0.68,0.21) (1.02,0.01) (1.36,0.10)};
  }%
}
% Mot-clé surligné (marqueur orange sous le texte)
\newcommand{\cle}[1]{{\popxbold\color{popdark}#1}}

% ============================================================
% En-tête / pied
% ============================================================
\pagestyle{fancy}
\fancyhf{}
\renewcommand{\headrulewidth}{0pt}
\renewcommand{\footrulewidth}{0pt}
\lhead{\raisebox{-0.15ex}{\poplogo\fontsize{13}{13}\selectfont\color{popink}OnBuch\color{poporange}+}}
\rhead{\poppill{\DOCMATIERE\ \textbullet\ \DOCNIVEAU\ \textbullet\ Leçon \DOCLECON}{white}{popink}}
\lfoot{\popxbold\fontsize{7.6}{7.6}\selectfont\color{popmuted}\MakeUppercase{Cours signé OnBuch+}\ \textbullet\ {\popsemi\DOCTITRE}}
\rfoot{\tikz[baseline=-0.6ex]\node[rounded corners=8.5pt,fill=popink,minimum height=17pt,minimum width=17pt,inner xsep=5pt,inner sep=0]{\popxbold\fontsize{8}{8}\selectfont\color{popcream}\thepage\,/\,\pageref*{LastPage}};}

% ============================================================
% Titres
% ============================================================
\newcommand{\chaptitre}[2]{%
  \begin{tcolorbox}[enhanced,colback=poporange,colframe=popink,boxrule=2.6pt,arc=11pt,
      drop shadow=popink,left=15pt,right=15pt,top=12pt,bottom=11pt,before skip=3pt,after skip=18pt]
    \poppill{#2}{popink}{popcream}\par\vspace{9pt}
    {\raggedright\hyphenpenalty=10000\exhyphenpenalty=10000\popdisplay\fontsize{22}{24}\selectfont\color{popcream}\MakeUppercase{#1}\par}
  \end{tcolorbox}
}
% Section numérotée : pastille orange avec le numéro + titre Archivo
\newcounter{popsec}
\newcommand{\coursec}[1]{%
  \stepcounter{popsec}\setcounter{popsub}{0}%
  \par\vspace{16pt}\noindent
  \begin{minipage}{\linewidth}
  \tikz[baseline=-0.7ex]\node[draw=popink,line width=1.6pt,fill=poporange,rounded corners=4pt,minimum size=22pt,inner sep=0]{\popdisplay\fontsize{12}{12}\selectfont\color{popcream}\thepopsec};%
  \hspace{8pt}{\popdisplay\fontsize{15}{17}\selectfont\color{popink}\MakeUppercase{#1}}\par
  \vspace{4pt}\noindent{\color{poporange}\rule{36pt}{3.6pt}}
  \end{minipage}\par\nopagebreak\vspace{8pt}
}
\newcounter{popsub}
\newcommand{\courssub}[1]{%
  \stepcounter{popsub}%
  \par\vspace{9pt}\noindent{\popxbold\fontsize{11.5}{13}\selectfont\color{popink}{\color{poporange}\thepopsec.\thepopsub}\hspace{6pt}#1}\par\nopagebreak\vspace{4pt}
}

% ============================================================
% Boîtes pédagogiques — même recette : fond blanc, bordure épaisse colorée,
% ombre dure encre, badge plein. Titre optionnel [#1].
% ============================================================
\tcbset{popbase/.style={breakable,enhanced,colback=white,boxrule=2.3pt,arc=9pt,
  shadow={3pt}{-3pt}{0pt}{popink},left=14pt,right=14pt,top=11pt,bottom=11pt,before skip=11pt,after skip=15pt,
  fontupper=\popsemi\fontsize{10.6}{14.5}\selectfont\color{popink},
  fontlower=\popsemi\fontsize{10.6}{14.5}\selectfont\color{popink}}}
% Coche dessinée (le glyphe ✓ n'existe pas dans les polices du document)
\newcommand{\popcheck}[1]{\tikz[baseline=-0.5ex]\draw[#1,line width=1.6pt,line cap=round,line join=round](-0.13,0.0)--(-0.04,-0.09)--(0.13,0.1);}
\providecommand{\coche}{\popcheck{popgreen}}
\providecommand{\resultat}[1]{\par\smallskip\noindent\fcolorbox{popgreen}{popgreenL}{\parbox{\dimexpr\linewidth-2\fboxsep-2\fboxrule\relax}{\textbf{Résultat.} #1}}\par\smallskip}
\newcommand{\popboxhead}[3]{\poppill{#1}{#2}{white}\ifx\relax#3\relax\else\hspace{6pt}{\popxbold\fontsize{10.6}{12}\selectfont\color{popink}#3}\fi\par\vspace{7pt}}

\newtcolorbox{aretenir}[1][]{popbase,colframe=popgreen,before upper={\popboxhead{À retenir}{popgreen}{#1}}}
% Texte en anglais (document de lecture, dialogue, extrait) : cadre
% d'étude. Un titre optionnel (source, auteur) s'affiche dans l'en-tête.
\newtcolorbox{textebox}[1][]{popbase,colframe=popblue,before upper={\popboxhead{Text}{popblue}{#1}\begin{otherlanguage}{english}},after upper={\end{otherlanguage}}}
% Marques ✓ / ✗ (forme correcte / fautive) souvent utilisées par les modèles
\providecommand{\cross}{\textcolor{poppink}{\ensuremath{\times}}}
\providecommand{\xmark}{\cross}\providecommand{\crossmark}{\cross}\providecommand{\cmark}{\ensuremath{\checkmark}}
% Ligne de réponse / justification à compléter (inventée par les modèles)
\providecommand{\justification}[1]{\par\noindent\textit{Justification~:} \dotfill\par}
% \textipa (package tipa non chargé) : la police ne couvre pas l'API -> simple passage
\providecommand{\textipa}[1]{#1}
% Soulignement (ulem non chargé) : \uline, \ul
\providecommand{\uline}[1]{\underline{#1}}\providecommand{\ul}[1]{\underline{#1}}
% \glossaire{\item ...} inventée par les modèles : liste de glossaire
\providecommand{\glossaire}[1]{\begin{itemize}#1\end{itemize}}
% \alert (beamer) inventée par les modèles : texte mis en évidence
\providecommand{\alert}[1]{\textbf{\textcolor{poppink}{#1}}}
% variantes ASCII d'ordinaux écrites en macro
\providecommand{\iere}{\textsuperscript{re}}\providecommand{\ieme}{\textsuperscript{e}}\providecommand{\ere}{\textsuperscript{re}}\providecommand{\eme}{\textsuperscript{e}}\providecommand{\er}{\textsuperscript{er}}\providecommand{\ie}{\textsuperscript{e}}
% Ordinaux français écrits en macro par les modèles : 1\ière, 3\ième
\newcommand{\ière}{\textsuperscript{re}}\newcommand{\ième}{\textsuperscript{e}}
% \exemple (inventée par les modèles) : étiquette « Exemple : »
\providecommand{\exemple}{\textbf{Exemple~:}\ }
% \square utilisable aussi hors mode math (cases à cocher)
\AtBeginDocument{\let\squareMath\square\renewcommand{\square}{\ensuremath{\squareMath}}}
% Mot ou phrase en anglais à l'intérieur d'un paragraphe français
\newcommand{\eng}[1]{\ifmmode\text{\textit{#1}}\else\begingroup\selectlanguage{english}\itshape #1\endgroup\fi}
\let\esp\eng % alias toléré
% flèches d'intonation inventées par les modèles
\providecommand{\textsearrow}{}\renewcommand{\textsearrow}{\ensuremath{\searrow}}\providecommand{\textnearrow}{}\renewcommand{\textnearrow}{\ensuremath{\nearrow}}
% \fr{..} : mot français (inventé par les modèles)
\providecommand{\fr}[1]{\textit{#1}}
\newtcolorbox{exemplebox}[1][]{popbase,colframe=poporange,before upper={\popboxhead{Exemple}{poporange}{#1}}}
\newtcolorbox{definition}[1][]{popbase,colframe=popblue,before upper={\popboxhead{Définition}{popblue}{#1}}}
\newtcolorbox{propriete}[1][]{popbase,colframe=poppurple,before upper={\popboxhead{Propriété}{poppurple}{#1}}}
\newtcolorbox{methode}[1][]{popbase,colframe=popgold,before upper={\popboxhead{Méthode}{popgold}{#1}}}
\newtcolorbox{attention}[1][]{popbase,colframe=poppink,before upper={\popboxhead{Attention}{poppink}{#1}}}
\newtcolorbox{experience}[1][]{popbase,colframe=popblue,colback=popblueL,before upper={\popboxhead{Expérience}{popblue}{#1}}}
% « Le savais-tu ? » : carte violette pleine (contexte, culture, Cameroun)
\newtcolorbox{savaistu}[1][]{popbase,colback=poppurple,colframe=popink,
  fontupper=\popsemi\fontsize{10.4}{14.2}\selectfont\color{white},
  before upper={\poppill{Le savais-tu\,?}{white}{poppurple}\ifx\relax#1\relax\else\hspace{6pt}{\popxbold\color{white}#1}\fi\par\vspace{7pt}}}
% Exercice résolu : énoncé en haut, \tcblower puis la solution sur fond crème
\newcounter{popexo}\newcounter{popexr}
\newtcolorbox{exoresolu}[1][]{popbase,colframe=poporange,colbacklower=popcream,
  segmentation style={popink,line width=1.2pt,dashed},
  before upper={\stepcounter{popexr}\popboxhead{Exercice résolu \thepopexr}{poporange}{#1}},
  before lower={\poppill{Solution}{popink}{popcream}\par\vspace{6pt}}}
% Exercice (non résolu en place) — la correction est dans la partie Corrigés
\newtcolorbox{exercice}[1][]{popbase,colframe=popink,
  before upper={\stepcounter{popexo}\popboxhead{Exercice \thepopexo}{popink}{#1}}}
% Corrigé
\newtcolorbox{corrige}[1][]{popbase,colframe=popgreen,colback=popgreenL,
  before upper={\popboxhead{Corrigé}{popgreen}{#1}}}
% Objectifs / prérequis / bilan
\newtcolorbox{objectifs}{popbase,colframe=popink,colback=poporangeL,
  before upper={\popboxhead{Objectifs de la leçon}{popink}{}}}
\newtcolorbox{prerequis}{popbase,colframe=popink,colback=popblueL,
  before upper={\popboxhead{Prérequis}{popblue}{}}}
\newtcolorbox{bilan}{popbase,colframe=popink,colback=white,boxrule=2.8pt,
  before upper={\popboxhead{Fiche bilan}{poporange}{L'essentiel à connaître pour l'examen}}}
% Figure encadrée avec légende
\newtcolorbox{popfigure}[1][]{enhanced,breakable=false,colback=white,colframe=popline,boxrule=1.4pt,arc=8pt,
  left=10pt,right=10pt,top=10pt,bottom=8pt,before skip=10pt,after skip=12pt,halign=center,
  lower separated=false,halign lower=center,
  fontlower=\popsemi\fontsize{9}{11.5}\selectfont\color{popmuted},
  #1}
\newcommand{\legende}[1]{\par\vspace{6pt}{\popsemi\fontsize{9}{11.5}\selectfont\color{popmuted}#1}}

% ============================================================
% Signature OnBuch+
% ============================================================
\newcommand{\poplogoinline}[1]{{\poplogo\fontsize{#1}{#1}\selectfont\color{popink}OnBuch\color{poporange}+}}
\newcommand{\signatureonbuch}{%
  \begin{tcolorbox}[enhanced,colback=popink,colframe=popink,boxrule=2.2pt,arc=10pt,
      drop shadow=poporange,left=14pt,right=14pt,top=10pt,bottom=10pt,before skip=0pt,after skip=0pt]
    \tikz[baseline=-0.7ex]\node[circle,fill=popgreen,draw=popcream,line width=1.3pt,minimum size=22pt,inner sep=0]{\popcheck{white}};%
    \hspace{9pt}\begin{minipage}[c]{0.62\linewidth}
      {\popxbold\fontsize{12}{13}\selectfont\color{popcream}Signé\ \textbullet\ {\poplogo OnBuch\color{poporange}+}}\par\vspace{2pt}
      {\raggedright\popsemi\fontsize{8}{10.5}\selectfont\color{popline}\MakeUppercase{Équipe pédagogique OnBuch+}\\\MakeUppercase{Conforme au programme officiel MINESEC}\par}
    \end{minipage}\hfill
    {\poplogo\fontsize{22}{22}\selectfont\color{popcream}OnBuch\color{poporange}+}
  \end{tcolorbox}%
}

% ============================================================
% Page de garde
% #1 titre · #2 matière · #3 niveau · #4 module · #5 n° leçon · #6 durée ·
% #7 description / objectifs (texte ou itemize)
% ============================================================
\newcommand{\pagegarde}[7]{%
  \thispagestyle{empty}
  \newgeometry{a4paper,margin=1.7cm,top=1.6cm,bottom=1.4cm}%
  \begin{tikzpicture}[remember picture,overlay]
    \fill[poporange!18] ([xshift=-3.2cm,yshift=-3.6cm]current page.north east) circle (4.6cm);
    \fill[poppurple!14] ([xshift=2.4cm,yshift=7.6cm]current page.south west) circle (3.8cm);
  \end{tikzpicture}%
  {\poplogo\fontsize{30}{30}\selectfont\color{popink}OnBuch\color{poporange}+}\hfill
  \raisebox{0.6ex}{\poppill{Cours signé \textbullet\ Premium}{popink}{popcream}}\par
  \vspace{1pt}\popsquiggle{poppurple}\par
  \vspace{8pt}
  \begin{tcolorbox}[enhanced,colback=poporange,colframe=popink,boxrule=2.8pt,arc=13pt,
      drop shadow=popink,left=18pt,right=18pt,top=12pt,bottom=13pt,after skip=10pt]
    \poppill{#2 \textbullet\ #3}{popink}{popcream}\hspace{5pt}\poppill{#4}{white}{popink}\par\vspace{8pt}
    {\popdisplay\fontsize{14}{14}\selectfont\color{popink}LEÇON #5}\par\vspace{4pt}
    {\raggedright\hyphenpenalty=10000\exhyphenpenalty=10000\popdisplay\fontsize{25}{27}\selectfont\color{popcream}\MakeUppercase{#1}\par}
    \vspace{8pt}
    {\popxbold\fontsize{9.5}{9.5}\selectfont\color{popink}\MakeUppercase{Durée conseillée : #6}}
  \end{tcolorbox}
  \begin{tcolorbox}[enhanced,colback=white,colframe=popink,boxrule=2.2pt,arc=10pt,
      drop shadow=popink,left=16pt,right=16pt,top=10pt,bottom=10pt,after skip=8pt]
    \poppill{Dans cette leçon}{poppurple}{white}\par\vspace{5pt}
    \popsemi\fontsize{9.3}{12.6}\selectfont\color{popink}#7
  \end{tcolorbox}
  \noindent\begin{minipage}[t]{0.487\linewidth}
    \begin{tcolorbox}[enhanced,colback=popgreen,colframe=popink,boxrule=2.2pt,arc=10pt,
        drop shadow=popink,left=11pt,right=11pt,top=10pt,bottom=10pt,height=3.1cm,valign=center]
      \tcbox[colback=white,colframe=popink,boxrule=1.3pt,arc=5pt,boxsep=0pt,nobeforeafter,
        left=3pt,right=3pt,top=3pt,bottom=3pt,valign=center]{\qrcode[height=1.9cm]{https://play.google.com/store/apps/details?id=com.luvvix.onbuch&hl=fr}}%
      \hspace{8pt}\begin{minipage}[c]{0.5\linewidth}
        {\popdisplay\fontsize{11}{12.5}\selectfont\color{white}\MakeUppercase{Télécharge OnBuch}}\par\vspace{4pt}
        {\raggedright\popsemi\fontsize{8.4}{11}\selectfont\color{white}Gratuit \textbullet\ Google Play\\Tuteur IA, annales, concours, résultats\par}
      \end{minipage}
    \end{tcolorbox}
  \end{minipage}\hfill
  \begin{minipage}[t]{0.487\linewidth}
    \begin{tcolorbox}[enhanced,colback=poppurple,colframe=popink,boxrule=2.2pt,arc=10pt,
        drop shadow=popink,left=11pt,right=11pt,top=10pt,bottom=10pt,height=3.1cm,valign=center]
      \tikz[baseline=-0.7ex]\node[circle,fill=white,draw=popink,line width=1.3pt,minimum size=1.6cm,inner sep=0]{\popdisplay\fontsize{24}{24}\selectfont\color{poppurple}+};%
      \hspace{8pt}\begin{minipage}[c]{0.6\linewidth}
        {\popdisplay\fontsize{11}{12.5}\selectfont\color{white}\MakeUppercase{Passe à OnBuch+}}\par\vspace{4pt}
        {\raggedright\popsemi\fontsize{8.4}{11}\selectfont\color{white}Tous les cours signés, Tuteur IA illimité, fiches PDF — onglet OnBuch+ de l'app\par}
      \end{minipage}
    \end{tcolorbox}
  \end{minipage}
  \vspace{8pt}
  \popcontenu
  \vfill
  \signatureonbuch
  \restoregeometry
  \newpage
}

% Rangée « Ce cours contient » (page de garde)
\newcommand{\poptile}[3]{% #1 couleur #2 gros texte #3 légende
  \begin{tcolorbox}[enhanced,colback=#1,colframe=popink,boxrule=2pt,arc=9pt,shadow={2.5pt}{-2.5pt}{0pt}{popink},
      left=6pt,right=6pt,top=9pt,bottom=9pt,halign=center,valign=center,height=1.75cm,width=\linewidth,nobeforeafter]
    {\popdisplay\fontsize{12.5}{13}\selectfont\color{white}\MakeUppercase{#2}}\par\vspace{3pt}
    {\centering\popsemi\fontsize{7.8}{9.5}\selectfont\color{white}#3\par}
  \end{tcolorbox}}
\newcommand{\popcontenu}{%
  \par\noindent{\popxboldls\fontsize{7.5}{7.5}\selectfont\color{popmuted}CE COURS SIGNÉ CONTIENT}\par\vspace{7pt}
  \noindent\begin{minipage}[t]{0.235\linewidth}\poptile{popblue}{Cours}{complet, illustré, pas à pas}\end{minipage}\hfill
  \begin{minipage}[t]{0.235\linewidth}\poptile{poporange}{Méthodes}{et exercices résolus}\end{minipage}\hfill
  \begin{minipage}[t]{0.235\linewidth}\poptile{poppink}{Exercices}{type examen + corrigés}\end{minipage}\hfill
  \begin{minipage}[t]{0.235\linewidth}\poptile{popgreen}{Fiche bilan}{l'essentiel à retenir}\end{minipage}\par}

% Sommaire
\newtcolorbox{sommaire}{enhanced,colback=white,colframe=popink,boxrule=2.2pt,arc=10pt,
  drop shadow=popink,left=18pt,right=18pt,top=13pt,bottom=13pt,before skip=6pt,after skip=20pt,
  before upper={\poppill{Sommaire}{poppurple}{white}\par\vspace{9pt}\popsemi\fontsize{10.6}{14.5}\selectfont\color{popink}}}

% Signature de fin de cours — poussée tout en bas de la dernière page
\newcommand{\findecours}{%
  \par\vfill\nopagebreak
  \begin{center}\popsquiggle{poporange}\end{center}\vspace{4pt}
  \signatureonbuch
}
\AtBeginDocument{\def\llbracket{\mathopen{[\mkern-3.5mu[}}\def\rrbracket{\mathclose{]\mkern-3.5mu]}}} % intervalles d'entiers (glyphe ⟦ absent de la police)
% Glyphes absents de Fira Math : redéfinis à partir de glyphes disponibles
\AtBeginDocument{%
  \def\setminus{\mathbin{\backslash}}\let\smallsetminus\setminus
  \def\vdots{\mathord{\vbox{\baselineskip3.2pt\lineskiplimit0pt\kern4pt\hbox{.}\hbox{.}\hbox{.}}}}%
  \def\ddots{\mathinner{\mkern1mu\raise6.5pt\hbox{.}\mkern2mu\raise3.6pt\hbox{.}\mkern2mu\raise0.7pt\hbox{.}\mkern1mu}}%
  \def\triangle{\mathord{\scalebox{0.9}{$\symup{\Delta}$}}}%
  \def\nabla{\mathord{\raisebox{\depth}{\scalebox{1}[-1]{$\symup{\Delta}$}}}}%
  \def\checkmark{\popcheck{popdark}}%
  \def\star{\mathbin{\ast}}\def\bigstar{\mathord{\ast}}%
  \def\diamond{\mathbin{\scalebox{0.6}{$\Diamond$}}}%
  \def\bigcup{\mathop{\mathchoice{\raisebox{-0.3em}{\scalebox{1.7}{$\cup$}}}{\scalebox{1.25}{$\cup$}}{\cup}{\cup}}\displaylimits}%
  \def\bigcap{\mathop{\mathchoice{\raisebox{-0.3em}{\scalebox{1.7}{$\cap$}}}{\scalebox{1.25}{$\cap$}}{\cap}{\cap}}\displaylimits}%
  \def\lesseqqgtr{\mathrel{\lessgtr}}\def\bigoplus{\mathop{\oplus}\displaylimits}%
}
\providecommand{\ssi}{si et seulement si} % abréviation fréquente
\AtBeginDocument{\providecommand{\wideparen}{\overparen}} % arc de cercle

\begin{document}

\pagegarde{Lesson 49 — Listening: Texts about subscribing to service packages (telephone/internet services)}{Anglais}{Tle SH}{Chapter 5 : Media and Communication}{EN49}{2 h}{Cette leçon de compréhension orale entraîne les élèves à écouter des textes sur la souscription à des forfaits téléphoniques et internet. Elle développe les stratégies d'écoute (repérage du thème, des mots-clés, des chiffres et des noms propres), enrichit le lexique des télécommunications et prépare aux questions de compréhension de type Bac.
\vspace{4pt}
\begin{multicols}{2}\small
\begin{itemize}
\item À la fin de cette leçon, je sais identifier le thème général d'un enregistrement sur les forfaits téléphone/internet dès la première écoute.
\item Je sais repérer et noter les mots-clés, les chiffres (prix, volumes de données, durées) et les noms propres (opérateurs, offres) dans un texte écouté.
\item Je sais répondre à des questions de compréhension variées : vrai/faux, QCM et réponses courtes.
\item Je connais et j'emploie le lexique essentiel des télécommunications : subscribe, bundle, data, airtime, top up, validity, etc.
\item Je sais distinguer les informations essentielles des détails secondaires dans une annonce commerciale.
\item Je sais anticiper le contenu d'un enregistrement à partir du titre et des questions avant l'écoute.
\end{itemize}\end{multicols}}

\begin{sommaire}
\begin{multicols}{2}
\begin{enumerate}
\item Avant l'écoute : préparation et anticipation
\item Le texte support : script d'écoute
\item Compréhension : questions variées
\item Lexique thématique : télécommunications et forfaits
\item Méthode d'écoute : stratégies et exemple traité
\item Production guidée : réemploi du lexique
\item Activité d'intégration
\item Méthodes et exercices résolus
\item Exercices
\item Corrigés des exercices
\item Fiche bilan
\end{enumerate}
\end{multicols}
\end{sommaire}

\begin{objectifs}
\begin{itemize}
\item À la fin de cette leçon, je sais identifier le thème général d'un enregistrement sur les forfaits téléphone/internet dès la première écoute.
\item Je sais repérer et noter les mots-clés, les chiffres (prix, volumes de données, durées) et les noms propres (opérateurs, offres) dans un texte écouté.
\item Je sais répondre à des questions de compréhension variées : vrai/faux, QCM et réponses courtes.
\item Je connais et j'emploie le lexique essentiel des télécommunications : subscribe, bundle, data, airtime, top up, validity, etc.
\item Je sais distinguer les informations essentielles des détails secondaires dans une annonce commerciale.
\item Je sais anticiper le contenu d'un enregistrement à partir du titre et des questions avant l'écoute.
\item Je sais réemployer le lexique du thème dans de courtes phrases personnelles.
\item Je sais comprendre un dialogue de souscription entre un client et un agent commercial.
\end{itemize}
\end{objectifs}

\begin{prerequis}
\begin{itemize}
\item Les nombres en anglais (cardinaux, prix, pourcentages) vus en Première
\item Le vocabulaire de base des médias et de la communication (Chapter 5, leçons précédentes)
\item Les questions en WH- (what, how much, how long) et leur formation
\item Le présent simple et le futur avec 'will' pour les offres et promotions
\item La prise de notes rapide pendant l'écoute (abréviations, symboles)
\end{itemize}
\end{prerequis}

\newpage

\coursec{Avant l'écoute : préparation et anticipation}

\courssub{Situation d'accroche : le SMS d'Aminatou}

Imaginez la scène. Il est 18 h à Buea, sur les flancs du mont Cameroun. Aminatou, élève de Terminale A, révise ses notes d'histoire quand son téléphone vibre. Un SMS de son opérateur s'affiche :

\begin{textebox}[Le SMS reçu par Aminatou]
“Dear customer, subscribe to our weekly bundle: 2 GB + 100 SMS for only 1,000 FCFA. Dial *123\# to activate. Offer valid for 7 days. Terms and conditions apply.”
\par
\emph{« Cher client, souscrivez à notre forfait hebdomadaire : 2 Go + 100 SMS pour seulement 1 000 FCFA. Composez le *123\# pour activer. Offre valable 7 jours. Conditions générales applicables. »}
\end{textebox}

Faut-il souscrire ? Que comprend exactement ce forfait ? Les 2 GB sont-ils valables jour et nuit ? Les 100 SMS sont-ils vers tous les réseaux ? Que se passe-t-il au bout de 7 jours si le crédit n'est pas épuisé ? Pour répondre, Aminatou doit \cle{comprendre précisément} un message commercial en anglais — exactement la compétence que cette leçon développe.

Au Royaume-Uni comme au Nigeria, au Ghana comme au Cameroun anglophone, choisir un \eng{service package} (un forfait téléphone ou internet) exige de comprendre des annonces publicitaires, des messages automatiques et des conversations avec un agent commercial. À l'examen du Baccalauréat, l'épreuve de compréhension orale repose souvent sur ce type de document authentique : une annonce, un dialogue en agence, un message vocal.

\textbf{Problématique de la leçon :} comment préparer son écoute pour comprendre un texte sur la souscription à un forfait téléphonique ou internet, et repérer rapidement les informations essentielles (prix, volume de données, durée de validité, mode de souscription) ?

\courssub{Brainstorming : les opérateurs que vous connaissez déjà}

Avant toute écoute, un bon auditeur active ses connaissances. Vous en possédez déjà beaucoup sur ce thème, car vous utilisez un téléphone tous les jours. Faisons l'inventaire en anglais.

\begin{exemplebox}[Opérateurs et fournisseurs d'accès : Cameroun et monde anglophone]
\begin{itemize}
\item \textbf{Cameroon:} MTN Cameroon, Orange Cameroon, Camtel, Nexttel — \emph{les opérateurs présents au Cameroun.}
\item \textbf{Nigeria:} MTN Nigeria, Airtel Nigeria, Glo (Globacom), 9mobile — \emph{le marché nigérian, géant des télécoms africaines.}
\item \textbf{United Kingdom:} Vodafone UK, EE, O2, Three — \emph{les grands réseaux britanniques.}
\item \textbf{United States:} Verizon, AT\&T, T-Mobile — \emph{les opérateurs américains.}
\item \textbf{Ghana / Kenya:} Vodafone Ghana, Safaricom (Kenya) — \emph{Safaricom est célèbre pour le paiement mobile M-Pesa.}
\end{itemize}
\end{exemplebox}

\begin{definition}[Le vocabulaire de base du thème]
\begin{itemize}
\item \cle{operator} / \cle{provider} (n.) : opérateur / fournisseur d'accès. \eng{MTN is the biggest mobile operator in Cameroon.} « MTN est le plus grand opérateur mobile au Cameroun. »
\item \cle{service package} (n.) : forfait, offre groupée de services. \eng{This package includes calls, SMS and data.} « Ce forfait comprend les appels, les SMS et les données. »
\item \cle{bundle} (n.) : forfait groupé (terme courant en Afrique anglophone et au Cameroun). \eng{I bought a weekly data bundle.} « J'ai acheté un forfait internet hebdomadaire. »
\item \cle{to subscribe (to)} (v.) : souscrire (à), s'abonner (à). \eng{She subscribed to a monthly internet plan.} « Elle a souscrit à un forfait internet mensuel. »
\item \cle{subscription} (n.) : abonnement, souscription. \eng{My subscription expires tomorrow.} « Mon abonnement expire demain. »
\end{itemize}
\end{definition}

\begin{savaistu}[Plan, contract, bundle : un mot pour chaque pays]
Le même objet porte des noms différents selon les pays anglophones. Au Royaume-Uni, on parle de \eng{mobile plan} ou de \eng{contract} (« contrat », souvent de 12 ou 24 mois, avec un téléphone inclus). Aux États-Unis, on dit \eng{cell phone plan}. Au Cameroun, au Nigeria et au Ghana, le mot courant est \eng{bundle} ou \eng{package}, car la plupart des clients achètent des forfaits prépayés (\eng{prepaid}) plutôt que des abonnements mensuels facturés (\eng{postpaid}). Au Bac, les trois termes peuvent apparaître : sachez qu'ils désignent la même réalité.
\end{savaistu}

\courssub{Prédire le contenu à partir du titre}

Le titre de notre écoute est : \eng{Subscribing to service packages}. C'est une mine d'informations. Le verbe \eng{subscribe} annonce une action du client ; \eng{service packages} annonce le thème. À partir de ce seul titre, on peut prédire les mots que l'on va probablement entendre.

\begin{methode}[Anticiper avant l'écoute : la méthode des 5 mots]
\begin{enumerate}
\item \textbf{Lisez le titre} et soulignez les mots porteurs de sens (ici : \eng{subscribing}, \eng{service packages}).
\item \textbf{Lisez toutes les questions} de compréhension avant la première écoute : elles vous indiquent quelles informations chercher (un prix ? une durée ? un code à composer ?).
\item \textbf{Listez au moins 5 mots} que vous vous attendez à entendre. Exemple ici : \eng{price}, \eng{data}, \eng{validity}, \eng{dial}, \eng{customer care}.
\item \textbf{Traduisez mentalement} ces mots et rappelez-vous leur prononciation : vous les reconnaîtrez plus vite à l'oreille.
\item \textbf{Formulez des hypothèses} : qui parle ? (un agent commercial ? une voix automatique ?) Où ? (en agence ? au téléphone ?)
\end{enumerate}
Cette démarche active vos connaissances et facilite considérablement la compréhension : le cerveau entend mieux ce qu'il attend.
\end{methode}

\begin{attention}[Le piège de la première écoute]
Ne cherchez pas à comprendre \textbf{chaque mot} à la première écoute. C'est l'erreur la plus fréquente des élèves francophones : dès qu'un mot inconnu passe, ils bloquent, paniquent et perdent la suite du texte. À la première écoute, visez uniquement le \cle{sens général} (le \eng{gist}) : \emph{Who is speaking? What is the text about? Where does it take place?} Les détails (chiffres, noms, codes) seront relevés à la deuxième écoute.
\end{attention}

\courssub{Écoute globale et écoute sélective : deux écoutes, deux objectifs}

\begin{definition}[Gist listening et selective listening]
\begin{itemize}
\item \cle{Gist listening} (écoute globale) : comprendre le sens général, le thème, la situation, les personnages — sans s'arrêter sur les détails.
\item \cle{Selective listening} ou \cle{listening for specific information} (écoute sélective) : repérer une information précise annoncée par les questions — un prix, un numéro, une date de validité. En anglais, on parle aussi de \eng{scanning for details}.
\end{itemize}
\end{definition}

\begin{center}
\begin{tabularx}{\linewidth}{|X|X|X|}
\hline
\rowcolor{popblueL} Écoute globale (gist) & Écoute sélective (détails) & Question type \\
\hline
Who is speaking? & How much does the bundle cost? & “The agent offers a package for 1,000 FCFA.” \\
\hline
What is the topic? & How many gigabytes are included? & “The package includes 2 GB of data.” \\
\hline
Where does it happen? & How long is the offer valid? & “The offer is valid for 7 days.” \\
\hline
What is the speaker's purpose? & Which code must the customer dial? & “To subscribe, dial *123\#.” \\
\hline
\end{tabularx}
\end{center}

\begin{popfigure}
\begin{tikzpicture}[node distance=0.9cm and 1.4cm,
box/.style={draw=popblue, fill=popblueL, rounded corners, thick, align=center, minimum height=0.9cm, text width=3.1cm},
arr/.style={-{Stealth[length=3mm]}, thick, popdark}]
\node[box, align=center] (titre){\textbf{Titre} \\ “Subscribing to service packages”};
\node[box, right=of titre, align=center] (pred){\textbf{Prédiction} \\ 5 mots attendus};
\node[box, right=of pred, align=center] (ec1){\textbf{1re écoute} \\ gist : who? what? where?};
\node[box, below=of ec1, align=center] (ec2){\textbf{2e écoute} \\ détails : prix, data, durée, code};
\node[box, left=of ec2, align=center] (verif){\textbf{Vérification} \\ réponses aux questions};
\draw[arr] (titre) -- (pred);
\draw[arr] (pred) -- (ec1);
\draw[arr] (ec1) -- (ec2);
\draw[arr] (ec2) -- (verif);
\end{tikzpicture}
\legende{La démarche d'écoute en cinq étapes : du titre à la vérification.}
\end{popfigure}

\courssub{La carte mentale du thème : organiser le lexique attendu}

Pour prédire efficacement, il est utile de visualiser le champ lexical du thème. La carte mentale ci-dessous rassemble les mots que l'on entend presque toujours dans un texte sur les forfaits.

\begin{popfigure}
\begin{tikzpicture}[mindmap, grow cyclic,
every node/.style={concept, align=center},
root concept/.append style={concept color=popblue, font=\bfseries},
level 1/.append style={level distance=3.2cm, sibling angle=120, concept color=poporange},
level 2/.append style={level distance=2.4cm, sibling angle=45, concept color=popgreenL}]
\node[root concept]{Service packages}
child { node {phone}
  child { node[align=center]{price \\ cost}}
  child { node[align=center]{validity \\ expiry}}
}
child { node {internet}
  child { node[align=center]{data \\ GB / MB}}
  child { node[align=center]{how to subscribe \\ dial *123\#}}
}
child { node {TV}
  child { node[align=center]{channels \\ decoder}}
  child { node {monthly fee}}
};
\end{tikzpicture}
\legende{Carte mentale lexicale : les mots-clés à attendre dans un texte sur les forfaits.}
\end{popfigure}

\courssub{Entraînement de l'oreille : chiffres, prix, données et durées}

Dans ce type de texte, les informations les plus souvent demandées aux questions sont des \cle{chiffres}. Or les nombres sont un piège classique pour les francophones : \eng{thirteen} et \eng{thirty}, \eng{fifteen} et \eng{fifty} se ressemblent ; les prix se disent différemment en anglais.

\begin{propriete}[Dire les prix, les volumes de données et les durées]
\begin{itemize}
\item \textbf{Prix :} on énonce le nombre puis la devise. \eng{1,000 FCFA} se dit \eng{one thousand francs} ou \eng{one thousand CFA francs}. \eng{2,500 FCFA} : \eng{two thousand five hundred francs}.
\item \textbf{Données internet :} \eng{2 GB} se dit \eng{two gigabytes} (prononcez « tou gui-ga-baïts ») ; \eng{500 MB} : \eng{five hundred megabytes}. Dans la conversation, on entend souvent simplement \eng{two gigs} (« tou guiz »).
\item \textbf{Durées de validité :} \eng{30 days} « trente jours » ; \eng{a weekly bundle} « un forfait hebdomadaire » ; \eng{a monthly plan} « un forfait mensuel » ; \eng{valid for 7 days} « valable 7 jours ».
\item \textbf{Codes USSD :} les chiffres se dictent un par un : \eng{*123\#} se dit \eng{star one two three hash}.
\end{itemize}
\end{propriete}

\begin{attention}[Chiffres : les erreurs typiques des francophones]
\begin{itemize}
\item \eng{thirteen} (13) / \eng{thirty} (30) : l'accent tombe sur la \textbf{deuxième} syllabe de \eng{thir\emph{teen}}, sur la \textbf{première} de \eng{\emph{thir}ty}. Entraînez-vous : \eng{fifteen GB} (15) contre \eng{fifty SMS} (50).
\item Ne dites jamais \eng{one thousand and five hundred} pour 1 500 : dites \eng{one thousand five hundred} ou \eng{fifteen hundred}. Le \eng{and} ne s'emploie qu'avant les dizaines : \eng{two thousand and fifty} (2 050).
\item En anglais, le séparateur des milliers est la virgule et le séparateur décimal est le point : \eng{1,000 FCFA} = 1 000 FCFA ; \eng{2.5 GB} = 2,5 Go. Attention aux inversions !
\end{itemize}
\end{attention}

\courssub{Lecture préalable des questions : cibler son écoute}

Voici le texte support de la leçon — une conversation en agence — que le professeur lira à haute voix à la section suivante. Avant l'écoute, observez d'abord les questions : que vous apprennent-elles sur ce que vous devez écouter ?

\begin{textebox}[At the MTN service centre in Buea — script d'écoute]
\textbf{Agent:} Good morning, Madam. Welcome to MTN Buea. How may I help you?

\textbf{Aminatou:} Good morning. I received a text message about a weekly bundle, but I'd like to understand it better before I subscribe.

\textbf{Agent:} Of course. We have three weekly packages at the moment. The first one costs five hundred francs and gives you seven hundred megabytes of data, valid for seven days. The second one is our most popular offer: for one thousand francs, you get two gigabytes of data plus one hundred SMS to all networks, also valid for seven days.

\textbf{Aminatou:} I see. And is there a bigger package? I use the internet a lot for my school research.

\textbf{Agent:} Yes, there is. For two thousand five hundred francs, you can have five gigabytes, unlimited SMS, and fifty minutes of calls to all networks. This one is valid for thirty days — it's actually a monthly package.

\textbf{Aminatou:} Hmm, that's interesting, but it's a bit expensive for me. How do I subscribe to the one-thousand-franc bundle?

\textbf{Agent:} It's very simple. Just dial star one two three hash on your phone, then choose option two. You will receive a confirmation message immediately. Please make sure you have at least one thousand francs of airtime on your account.

\textbf{Aminatou:} All right. And what happens if I don't use all my data before the seven days?

\textbf{Agent:} Unfortunately, the remaining data expires. But if you renew the bundle before it expires, your unused data will be carried over to the next week.

\textbf{Aminatou:} Perfect. Thank you very much for your help.

\textbf{Agent:} You're welcome, Madam. Have a nice day!
\end{textebox}

\textbf{Glossaire :} \emph{bundle} (n.) forfait ; \emph{airtime} (n.) crédit téléphonique ; \emph{to expire} (v.) expirer ; \emph{to renew} (v.) renouveler ; \emph{unused} (adj.) non utilisé ; \emph{to be carried over} (v.) être reporté ; \emph{option} (n.) option, choix.

\textbf{Questions à lire AVANT l'écoute :}
\begin{enumerate}
\item How many weekly packages does the agent mention?
\item How much does the most popular bundle cost?
\item What do you get for 2,500 FCFA?
\item Which code must Aminatou dial to subscribe?
\item What happens to unused data if the bundle is renewed before it expires?
\end{enumerate}

Remarquez comme la lecture préalable des questions transforme votre écoute : la question 2 vous dit d'écouter un \textbf{prix}, la question 4 un \textbf{code}, la question 5 une \textbf{condition}. Vous savez déjà quoi chercher.

\begin{exoresolu}[Prédire les mots-clés à partir des questions]
Énoncé : sans écouter le texte, lisez la question \eng{“How much does the most popular bundle cost?”} et indiquez : (a) quel type d'information chercher ; (b) trois mots anglais que vous vous attendez à entendre ; (c) la réponse, après lecture du script.
\tcblower
Solution : (a) On cherche un \textbf{prix}, donc un nombre suivi d'une devise. (b) Les mots attendus sont par exemple \eng{costs}, \eng{francs}, \eng{popular}, \eng{thousand}. (c) Le script dit : \eng{“The second one is our most popular offer: for one thousand francs...”} La réponse est donc : \eng{It costs one thousand francs (1,000 FCFA).} « Il coûte mille francs. » La prédiction était exacte : le mot \eng{popular} dans la question est repris tel quel dans le texte — c'est un \cle{signal d'écoute} classique.
\end{exoresolu}

\begin{exercice}[À vous de jouer : préparez votre écoute]
\begin{enumerate}
\item Listez cinq opérateurs téléphoniques : deux camerounais, deux nigérians, un britannique.
\item À partir du titre \eng{“Choosing an internet package”}, listez cinq mots anglais que vous vous attendez à entendre.
\item Écrivez en toutes lettres en anglais : 500 FCFA ; 5 GB ; 30 days ; *126\#.
\item Pour chaque question, dites si elle demande une écoute globale ou sélective : (a) \eng{What is the conversation about?} (b) \eng{How many SMS are included?}
\end{enumerate}
\end{exercice}

\begin{corrige}[Exercice : préparez votre écoute]
\begin{enumerate}
\item Exemples : MTN Cameroon, Orange Cameroon ; MTN Nigeria, Glo ; Vodafone UK.
\item Par exemple : \eng{data, price, gigabytes, subscribe, validity, dial} (cinq au choix).
\item \eng{five hundred francs} ; \eng{five gigabytes} ; \eng{thirty days} ; \eng{star one two six hash}.
\item (a) Écoute globale (le thème général) ; (b) écoute sélective (un chiffre précis).
\end{enumerate}
\end{corrige}

\begin{aretenir}[L'essentiel de la préparation à l'écoute]
\begin{itemize}
\item Avant d'écouter : lisez le \textbf{titre} et les \textbf{questions}, puis prédisez au moins \textbf{5 mots-clés}.
\item Première écoute = \textbf{sens général} (who? what? where?) ; deuxième écoute = \textbf{détails} (prix, volumes, durées, codes).
\item Les chiffres sont les cibles favorites des questions : entraînez votre oreille à \eng{thirteen/thirty}, aux prix en \eng{francs} et aux volumes en \eng{gigabytes}.
\item Le lexique du thème : \eng{operator, bundle, package, plan, to subscribe, validity, airtime, to expire}.
\end{itemize}
\end{aretenir}

\coursec{Avant l'écoute : préparation et anticipation}

Avant d'aborder les scripts d'écoute, il est essentiel de mettre en place des stratégies de \cle{pré-écoute} (pre-listening). Cette phase permet d'activer les connaissances antérieures, de lever les obstacles lexicaux et de formuler des hypothèses pour écouter de manière active et ciblée.

\begin{methode}[Trois étapes pour préparer l'écoute]
\begin{enumerate}
\item \textbf{Analyser le contexte et les consignes.} Lisez le titre de la leçon (\eng{Subscribing to service packages}), identifiez le thème (télécommunications, forfaits) et le type de documents attendus (annonce, dialogue service client). Lisez les questions de compréhension \textbf{avant} la première écoute : elles orientent votre attention vers les informations à repérer (prix, volumes, codes, durées).
\item \textbf{Activer le lexique thématique (brainstorming).} En classe ou seul, listez en 2 minutes tous les mots anglais qui vous viennent à l'esprit sur le thème : \eng{data, bundle, plan, subscribe, recharge, credit, balance, network, prepaid, postpaid, unlimited, validity, dial, USSD code}. Vérifiez la prononciation des mots-clés (voir boîtes \eng{Prononciation} ci-dessous).
\item \textbf{Anticiper le contenu (prédiction).} À partir du titre \eng{MTN Mega Bundle} et \eng{Camtel Call Centre}, formulez des hypothèses : quels chiffres allez-vous entendre ? (Prix en FCFA, Go, minutes, jours). Quelles actions ? (\eng{dial, subscribe, activate, pay}). Écrivez 3 questions que vous vous attendez à voir résolues.
\end{enumerate}
\end{methode}

\begin{experience}[Activité de classe : \eng{Mind Map} « Telecom Packages »]
Au tableau, l'enseignant écrit \eng{TELECOM PACKAGES} au centre. Les élèves viennent ajouter des branches en anglais : \eng{Mobile Data} (\eng{GB, MB, 4G, 5G}), \eng{Voice} (\eng{minutes, calls, all networks}), \eng{SMS} (\eng{texts, messages}), \eng{Money} (\eng{FCFA, price, cost, recharge, balance}), \eng{Time} (\eng{validity, daily, weekly, monthly}), \eng{Actions} (\eng{subscribe, dial, check, activate, renew}). Cette carte mentale reste visible pendant toute la leçon.
\end{experience}

\begin{propriete}[Prononciation des mots-clés techniques]
\begin{center}
\begin{tabularx}{\linewidth}{|l|l|X|}\hline
\rowcolor{popblueL} \textbf{Mot} & \textbf{Prononciation (approx. fr.)} & \textbf{Note} \\ \hline
\eng{bundle} & « beun-deul » & \eng{u} = /ʌ/ (comme \eng{cup}) \\ \hline
\eng{data} & « déi-ta » & \eng{a} final = /ə/ (schwa) \\ \hline
\eng{subscribe} & « sèbe-scri-be » & accent sur \eng{scribe} \\ \hline
\eng{validity} & « va-li-di-ti » & 4 syllabes, accent sur \eng{li} \\ \hline
\eng{dial} & « daï-èl » & diphtongue /aɪə/ \\ \hline
\eng{unlimited} & « an-li-mi-tid » & \eng{un-} = /ʌn/ \\ \hline
\eng{activation} & « èk-ti-vei-chn » & \eng{tion} = /ʃn/ \\ \hline
\eng{USSD} & « you-ès-ès-di » & épelé lettre par lettre \\ \hline
\eng{prepaid / postpaid} & « pri-peïd / peust-peïd » & \eng{paid} = /peɪd/ (comme \eng{paid}) \\ \hline
\end{tabularx}
\end{center}
\end{propriete}

\coursec{Le texte support : scripts d'écoute}

Cette section présente les deux documents audio authentiques (ou lus par l'enseignant) qui serviront de support aux trois écoutes successives. Le premier est une annonce publicitaire radio/téléphone (monologue), le second est une interaction réelle au centre d'appel (dialogue).

\courssub{Script 1 : une annonce publicitaire (MTN Mega Bundle)}

\begin{textebox}[Script 1 — MTN Mega Bundle (annonce publicitaire)]
“Welcome to MTN Mega Bundle! Get 5 GB of data, 200 SMS and 60 minutes of calls to all networks for only 2,500 FCFA, valid for 7 days. To subscribe, dial *157\# and follow the instructions. You can check your balance anytime by dialling *157*1\#. Offer available to all prepaid customers.”
\end{textebox}

\textbf{Glossaire du script 1}

\begin{center}
\begin{tabularx}{\linewidth}{|l|X|l|}\hline
\rowcolor{popblueL} \textbf{Mot anglais} & \textbf{Traduction française} & \textbf{Nature} \\ \hline
\eng{bundle} & forfait, bouquet de services & nom \\ \hline
\eng{data} & données (internet mobile) & nom \\ \hline
\eng{to all networks} & vers tous les réseaux & expression \\ \hline
\eng{valid for 7 days} & valable 7 jours & expression \\ \hline
\eng{to subscribe} & souscrire, s'abonner & verbe \\ \hline
\eng{to dial} & composer (un numéro/code) & verbe \\ \hline
\eng{to check your balance} & consulter votre solde & expression \\ \hline
\eng{prepaid customers} & clients prépayés & nom \\ \hline
\eng{instructions} & instructions, consignes & nom (pluriel) \\ \hline
\eng{anytime} & à tout moment & adverbe \\ \hline
\end{tabularx}
\end{center}

\begin{definition}[Prepaid ou postpaid ?]
Un service \cle{prepaid} (prononcé « pri-peïd ») est un service \textbf{payé à l'avance} : le client achète du crédit avant de consommer. C'est le système dominant au Cameroun (cartes prépayées, rechargements \eng{top-up}). À l'opposé, un service \cle{postpaid} (« peust-peïd ») est \textbf{facturé après consommation} : le client reçoit une facture (\eng{bill} ou \eng{invoice}) à la fin du mois, comme c'est souvent le cas en France ou aux États-Unis pour les forfaits avec engagement.
\end{definition}

\begin{exemplebox}[Phrases-clés du script 1, traduites et analysées]
\begin{itemize}
\item \eng{Get 5 GB of data for only 2,500 FCFA.} $\rightarrow$ « Obtenez 5 Go de données pour seulement 2 500 FCFA. » (Impératif publicitaire, \eng{Get} = obtenez/recevez).
\item \eng{To subscribe, dial *157\#.} $\rightarrow$ « Pour souscrire, composez le *157\#. » (Infinitif de but \eng{To subscribe} + impératif \eng{dial}).
\item \eng{The offer is valid for 7 days.} $\rightarrow$ « L'offre est valable 7 jours. » (\eng{valid for} + durée).
\item \eng{You can check your balance anytime.} $\rightarrow$ « Vous pouvez consulter votre solde à tout moment. » (\eng{can} capacité/permission + \eng{anytime} = \eng{any time}).
\end{itemize}
\end{exemplebox}

\begin{attention}[Attention aux chiffres à l'écoute : pièges \eng{-teen} vs \eng{-ty}]
Les francophones confondent souvent à l'oral les nombres en \eng{-teen} (adolescence, 13--19) et en \eng{-ty} (dizaines, 20--90).
\begin{itemize}
\item \eng{Fifteen} (15) /fɪfˈtiːn/ $\neq$ \eng{Fifty} (50) /ˈfɪfti/. Accent sur 2e syllabe pour \eng{-teen}, sur 1re pour \eng{-ty}.
\item \eng{Sixteen} (16) /sɪkˈstiːn/ $\neq$ \eng{Sixty} (60) /ˈsɪksti/.
\item Dans le script 1 : \eng{60 minutes} = « sik-sti » (accent début). \eng{2,500} = \eng{two thousand five hundred} (jamais \eng{twenty-five hundred} dans un texte formel).
\item \eng{7 days} = « seven days » (ne pas confondre avec \eng{seventy} = 70).
\end{itemize}
\end{attention}

\courssub{Script 2 : un dialogue au centre d'appel (Camtel)}

\begin{textebox}[Script 2 — At the Camtel call centre (dialogue)]
Agent: “Good morning, Camtel customer service. How may I help you?”

Customer: “Good morning. I'd like to subscribe to an internet package for my home.”

Agent: “We have three offers: the Basic plan at 10,000 FCFA per month with 20 GB, the Family plan at 20,000 FCFA with 50 GB, and the Premium plan at 35,000 FCFA with unlimited data.”

Customer: “How long does activation take?”

Agent: “Your connection will be activated within 48 hours after payment.”
\end{textebox}

\textbf{Glossaire du script 2}

\begin{center}
\begin{tabularx}{\linewidth}{|l|X|l|}\hline
\rowcolor{popblueL} \textbf{Mot anglais} & \textbf{Traduction française} & \textbf{Nature} \\ \hline
\eng{customer service} & service client / relation client & nom \\ \hline
\eng{How may I help you?} & Comment puis-je vous aider ? & formule de politesse \\ \hline
\eng{I'd like to...} & Je voudrais... & formule de demande (polie) \\ \hline
\eng{package / plan} & forfait, offre, formule & nom \\ \hline
\eng{unlimited data} & données illimitées & expression \\ \hline
\eng{activation} & activation, mise en service & nom \\ \hline
\eng{within 48 hours} & sous 48 heures (délai max) & expression \\ \hline
\eng{payment} & paiement & nom \\ \hline
\eng{per month} & par mois / mensuel & expression \\ \hline
\eng{Basic / Family / Premium} & Basique / Famille / Premium & adjectifs (noms d'offres) \\ \hline
\end{tabularx}
\end{center}

\begin{aretenir}[Les formules de politesse du service client (\eng{Customer Care})]
\begin{itemize}
\item \textbf{Ouverture (Agent)} : \eng{Good morning/afternoon, [Company] customer service. How may I help you?} / \eng{How can I help you?} / \eng{What can I do for you today?} — \textbf{Jamais} \eng{What do you want?} (très impoli).
\item \textbf{Demande (Client)} : \eng{I'd like to...} (= \eng{I would like to...}) — contraction \textbf{obligatoire} à l'oral (\eng{I'd} /aɪd/). \eng{I want to...} est trop direct, familier, à éviter dans un contexte formel.
\item \textbf{Demander une durée} : \eng{How long does [action] take?} (Ex. \eng{How long does activation take?}) $\rightarrow$ « Combien de temps prend l'activation ? / Combien de temps faut-il pour activer ? »
\item \textbf{Clôture} : \eng{Thank you for calling Camtel. Have a nice day.}
\end{itemize}
\end{aretenir}

\begin{propriete}[Comparaison structurelle des deux scripts]
\begin{center}
\begin{tabularx}{\linewidth}{|X|X|X|}\hline
\rowcolor{popblueL} \textbf{Critère} & \textbf{Script 1 : Annonce MTN} & \textbf{Script 2 : Dialogue Camtel} \\ \hline
\textbf{Type de texte} & Monologue publicitaire (one-way) & Dialogue interactionnel (two-way) \\ \hline
\textbf{Locuteurs} & Une voix (speaker/annonceur) & Deux voix : Agent + Client \\ \hline
\textbf{Canal / Lieu} & Radio, TV, SMS, IVR (serveur vocal) & Centre d'appel (téléphone) \\ \hline
\textbf{Fonction} & Informer, persuader, inciter à l'action (\eng{Call to action}) & Renseigner, conseiller, vendre, rassurer \\ \hline
\textbf{Registre de langue} & Publicitaire, direct, impératifs (\eng{Get, Dial, Check}) & Formel, poli, modalaux (\eng{may, would, will be}) \\ \hline
\textbf{Chiffres clés} & 5 GB, 200 SMS, 60 min, 2,500 FCFA, 7 jours, *157\# & 10k/20k/35k FCFA, 20/50 GB/unlimited, 48h \\ \hline
\textbf{Grammaire notable} & Impératif, Infinitif de but, \eng{valid for} & \eng{I'd like to}, \eng{How long... take?}, Passif futur (\eng{will be activated}), \eng{within} \\ \hline
\end{tabularx}
\end{center}
\end{propriete}

\coursec{Compréhension : les trois écoutes successives}

La méthode des trois écoutes est la stratégie standard au Baccalauréat et dans l'approche par les compétences. Elle structure l'effort d'attention : global $\rightarrow$ sélectif $\rightarrow$ vérification.

\courssub{Première écoute : identification globale (Global Listening)}

\begin{methode}[Première écoute : les 3 questions repères]
\begin{enumerate}
\item \textbf{What is the text about?} (Thème) $\rightarrow$ Mots-recurrents : \eng{bundle, data, subscribe, plan, package, internet, offer}.
\item \textbf{Who is speaking?} (Locuteurs) $\rightarrow$ Une voix (annonce) OU deux voix (dialogue : agent/client). Rôles ?
\item \textbf{Where / Through which channel?} (Lieu/Canal) $\rightarrow$ Indices : \eng{customer service, call centre, dial, welcome to MTN, radio}.
\end{enumerate}
\end{methode}

\begin{exemplebox}[Réponses attendues après 1ère écoute]
\begin{itemize}
\item \textbf{Script 1} : \eng{It's an advertisement for a mobile phone bundle called Mega Bundle.} (C'est une pub pour un forfait mobile appelé Mega Bundle.)
\item \textbf{Script 2} : \eng{It's a conversation between a customer and a Camtel agent about home internet packages.} (C'est une conversation entre un client et un agent Camtel sur les forfaits internet domicile.)
\end{itemize}
\end{exemplebox}

\courssub{Deuxième écoute : relevé sélectif des informations chiffrées (Selective Listening)}

C'est l'écoute la plus importante pour la note. On ne cherche \textbf{que} les données factuelles (prix, volumes, codes, durées, noms propres). Préparez un tableau de prise de notes \textbf{avant} d'écouter.

\begin{popfigure}
\begin{center}
\begin{tabular}{|l|c|c|c|c|c|}\hline
\rowcolor{poporangeL} \textbf{Offer Name} & \textbf{Price (FCFA)} & \textbf{Data} & \textbf{Calls / SMS} & \textbf{Validity / Duration} & \textbf{Action / Code} \\ \hline
\rule[-0.35cm]{0pt}{0.9cm} MTN Mega Bundle & & & & & \\ \hline
\rule[-0.35cm]{0pt}{0.9cm} Camtel Basic & & & -- & & \\ \hline
\rule[-0.35cm]{0pt}{0.9cm} Camtel Family & & & -- & & \\ \hline
\rule[-0.35cm]{0pt}{0.9cm} Camtel Premium & & & -- & & \\ \hline
\end{tabular}
\end{center}
\legende{Tableau de prise de notes vierge à compléter pendant la 2ème écoute. Utilisez des abréviations : GB, min, SMS, d (days), mo (month), h (hours).}
\end{popfigure}

\begin{exoresolu}[Prise de notes complète après 2ème écoute]
Énoncé : Complétez le tableau ci-dessus pour les deux scripts.
\tcblower
Solution détaillée :
\begin{itemize}
\item \textbf{MTN Mega Bundle} : Price = 2,500 ; Data = 5 GB ; Calls/SMS = 60 min + 200 SMS ; Validity = 7 days ; Code = *157\# (Balance: *157*1\#).
\item \textbf{Camtel Basic} : Price = 10,000/mo ; Data = 20 GB ; Validity = 1 month (implied \eng{per month}) ; Action = Subscribe at call centre.
\item \textbf{Camtel Family} : Price = 20,000/mo ; Data = 50 GB ; Validity = 1 month.
\item \textbf{Camtel Premium} : Price = 35,000/mo ; Data = Unlimited ; Activation = within 48h after payment.
\end{itemize}
\textbf{Astuce méthode} : Écrivez les chiffres en \textbf{chiffres arabes} (2,500 pas \eng{two thousand five hundred}). Notez les unités (GB, min, FCFA, days) pour éviter les confusions (ex: 7 $\neq$ 7 FCFA).
\end{exoresolu}

\begin{attention}[Pièges fréquents de la 2ème écoute]
\begin{itemize}
\item \textbf{Unité de temps} : \eng{Per month} (facturation mensuelle, script 2) $\neq$ \eng{Valid for 7 days} (durée du forfait, script 1). Ne mélangez pas « 10 000 FCFA pour 7 jours » !
\item \textbf{Unlimited} = \textbf{Illimité} (préfixe \eng{un-} = négation). Ne pas traduire par « limité ».
\item \textbf{Within 48 hours} = « \textbf{sous} 48 heures / \textbf{au plus tard dans} 48 heures » (délai maximum). Pas « dans exactement 48 heures » (\eng{in 48 hours}).
\item \textbf{Passif} : \eng{will be activated} (sera activée) $\neq$ \eng{will activate} (s'activera toute seule / activera qqch d'autre). Le sujet \eng{connection} subit l'action.
\end{itemize}
\end{attention}

\courssub{Troisième écoute : vérification et détails fins (Detailed Listening / Checking)}

\begin{methode}[Troisième écoute : vérification ciblée]
\begin{enumerate}
\item Relisez votre tableau : entourez les cases \textbf{vides} ou \textbf{incertaines}.
\item Concentrez-vous \textbf{uniquement} sur ces trous pendant l'écoute.
\item Vérifiez la \textbf{cohérence} : Prix en FCFA ? Volume en GB ? Durée en jours/mois ? Code USSD commence par * et finit par \# ?
\item Captez les \textbf{détails secondaires} (conditions, étapes suivantes, formules de politesse) souvent demandés en questions « Vrai/Faux » ou « Justifiez ».
\end{enumerate}
\end{methode}

\begin{exemplebox}[Détails souvent manqués, relevés à la 3ème écoute]
\begin{itemize}
\item \eng{Offer available to all prepaid customers.} $\rightarrow$ Condition d'éligibilité (pas pour postpaid).
\item \eng{...and follow the instructions.} $\rightarrow$ Étape 2 après composition du code (menu interactif).
\item \eng{Your connection will be activated within 48 hours after payment.} $\rightarrow$ Condition déclenchante : \eng{after payment}. Futur passif.
\item \eng{How may I help you?} $\rightarrow$ Marque de politesse formelle (\eng{may} > \eng{can}).
\end{itemize}
\end{exemplebox}

\coursec{Questions de compréhension variées (Entraînement Bac)}

Voici des exercices types du Baccalauréat (Vrai/Faux, QCM, Réponses courtes, Complétion) basés sur les deux scripts. Répondez en anglais.

\begin{exercice}[Compréhension de l'oral -- Scripts 1 \& 2]
\textbf{A. True or False? Justify with a quote from the text.}
\begin{enumerate}
\item The MTN Mega Bundle is valid for one month.
\item The Camtel Basic plan costs 10,000 FCFA per week.
\item To check your MTN balance, you dial *157\#.
\item The Camtel Premium plan offers limited data.
\item Activation of the Camtel connection is immediate after payment.
\end{enumerate}

\textbf{B. Multiple Choice Questions (MCQ).}
\begin{enumerate}
\item Who is the target audience of the MTN Mega Bundle?
      \begin{itemize}
      \item[a)] Postpaid customers only.
      \item[b)] All prepaid customers.
      \item[c)] Camtel agents.
      \item[d)] Family plan subscribers.
      \end{itemize}
\item What does the customer want in Script 2?
      \begin{itemize}
      \item[a)] A mobile data bundle.
      \item[b)] An internet package for home.
      \item[c)] To check his balance.
      \item[d)] To complain about activation.
      \end{itemize}
\item How many offers does the Camtel agent present?
      \begin{itemize}
      \item[a)] Two.
      \item[b)] Three.
      \item[c)] Four.
      \item[d)] Five.
      \end{itemize}
\end{enumerate}

\textbf{C. Short Answers (Answer in complete sentences).}
\begin{enumerate}
\item How much data does the MTN Mega Bundle provide?
\item What is the price of the Camtel Family plan?
\item What code must you dial to subscribe to the MTN Mega Bundle?
\item How long does the Camtel agent say activation takes?
\item Which plan offers unlimited data at Camtel?
\end{enumerate}

\textbf{D. Gap-fill (Complete the summary).}
The MTN Mega Bundle gives you 5 \_\_\_\_\_ of data, 200 \_\_\_\_\_ and 60 \_\_\_\_\_ of calls for 2,500 FCFA. It is \_\_\_\_\_ for 7 days. To \_\_\_\_\_, dial *157\#. At Camtel, the \_\_\_\_\_ plan costs 35,000 FCFA and gives \_\_\_\_\_ data. Activation takes \_\_\_\_\_ 48 hours \_\_\_\_\_ payment.
\end{exercice}

\begin{corrige}[Exercice : Compréhension de l'oral]
\textbf{A. True / False}
\begin{enumerate}
\item \textbf{False.} Quote: \eng{“valid for 7 days”} (not one month).
\item \textbf{False.} Quote: \eng{“the Basic plan at 10,000 FCFA per month”} (not per week).
\item \textbf{False.} Quote: \eng{“You can check your balance anytime by dialling *157*1\#”} (The code *157\# is for subscribing).
\item \textbf{False.} Quote: \eng{“the Premium plan at 35,000 FCFA with unlimited data”} (Unlimited = illimité).
\item \textbf{False.} Quote: \eng{“Your connection will be activated within 48 hours after payment”} (Not immediate).
\end{enumerate}

\textbf{B. MCQ}
\begin{enumerate}
\item \textbf{b)} (\eng{Offer available to all prepaid customers}).
\item \textbf{b)} (\eng{I'd like to subscribe to an internet package for my home}).
\item \textbf{b)} (Basic, Family, Premium).
\end{enumerate}

\textbf{C. Short Answers}
\begin{enumerate}
\item \eng{It provides 5 GB of data.}
\item \eng{The Camtel Family plan costs 20,000 FCFA per month.}
\item \eng{You must dial *157\# to subscribe.}
\item \eng{Activation takes within 48 hours after payment.}
\item \eng{The Premium plan offers unlimited data.}
\end{enumerate}

\textbf{D. Gap-fill}
The MTN Mega Bundle gives you 5 \textbf{GB} of data, 200 \textbf{SMS} and 60 \textbf{minutes} of calls for 2,500 FCFA. It is \textbf{valid} for 7 days. To \textbf{subscribe}, dial *157\#. At Camtel, the \textbf{Premium} plan costs 35,000 FCFA and gives \textbf{unlimited} data. Activation takes \textbf{within} 48 hours \textbf{after} payment.
\end{corrige}

\coursec{Lexique thématique : télécommunications et forfaits (Consolidation)}

Ce lexique synthétique regroupe le vocabulaire essentiel des deux scripts et l'étend aux notions connexes du programme (facturation, réclamations, résiliation). À mémoriser par catégories.

\begin{definition}[Lexique thématique organisé : \eng{Telecom Packages \& Services}]
\begin{center}
\begin{tabularx}{\linewidth}{|l|X|l|}\hline
\rowcolor{popblueL} \textbf{Catégorie} & \textbf{Mots / Expressions (English)} & \textbf{Traduction / Contexte} \\ \hline
\textbf{Types de forfaits} & \eng{bundle, package, plan, offer, deal, subscription} & forfait, offre, abonnement \\ \hline
& \eng{prepaid (pay-as-you-go), postpaid (contract)} & prépayé, postpayé (avec engagement) \\ \hline
& \eng{basic / standard / family / premium / unlimited} & basique, standard, famille, premium, illimité \\ \hline
\textbf{Volumes \& Unités} & \eng{data, mobile data, internet data} & données (internet) \\ \hline
& \eng{GB (Gigabyte), MB (Megabyte)} & Go, Mo \\ \hline
& \eng{minutes (mins), calls, voice minutes} & minutes d'appel \\ \hline
& \eng{SMS, texts, text messages} & SMS, textos \\ \hline
& \eng{unlimited / capped / throttled} & illimité / plafonné / bridé (débit réduit) \\ \hline
\textbf{Prix \& Facturation} & \eng{price, cost, rate, fee, charge} & prix, coût, tarif, frais \\ \hline
& \eng{per month / monthly, per day / daily, per week / weekly} & par mois / mensuel, par jour, par semaine \\ \hline
& \eng{FCFA, CFA franc} & Franc CFA \\ \hline
& \eng{recharge, top-up, credit, balance} & rechargement, crédit, solde \\ \hline
& \eng{bill, invoice, billing cycle, due date} & facture, cycle de facturation, date d'échéance \\ \hline
\textbf{Actions \& Verbes} & \eng{subscribe (to), sign up (for), register} & s'abonner, souscrire, s'inscrire \\ \hline
& \eng{dial, call, ring} & composer, appeler \\ \hline
& \eng{activate, activation, deactivate} & activer, activation, désactiver \\ \hline
& \eng{renew, auto-renewal} & renouveler, renouvellement automatique \\ \hline
& \eng{check (balance), monitor, track} & consulter (solde), surveiller \\ \hline
& \eng{upgrade / downgrade (plan)} & passer à un forfait supérieur / inférieur \\ \hline
& \eng{cancel, unsubscribe, opt out} & résilier, se désabonner \\ \hline
\textbf{Service Client} & \eng{customer service / care / support, call centre, agent} & service client, centre d'appel, agent \\ \hline
& \eng{query, complaint, issue, troubleshooting} & demande, réclamation, problème, dépannage \\ \hline
& \eng{refund, compensation} & remboursement, compensation \\ \hline
\textbf{Technique / USSD} & \eng{USSD code, shortcode, star (*), hash (\#) / pound (US)} & code USSD, code court, étoile, dièse \\ \hline
& \eng{network, coverage, signal, roaming} & réseau, couverture, signal, itinérance \\ \hline
& \eng{validity, expiry date, expiration} & validité, date d'expiration \\ \hline
\end{tabularx}
\end{center}
\end{definition}

\begin{attention}[Faux-amis et pièges de traduction (Spécial Telecom)]
\begin{itemize}
\item \eng{Subscription} $\neq$ « souscription » (acte unique) mais souvent « \textbf{abonnement} » (état durable). \eng{To subscribe} = s'abonner.
\item \eng{Plan} $\neq$ « plan » (carte, projet) mais « \textbf{forfait / formule} » commercial.
\item \eng{Deal} $\neq$ « marché / donne » mais « \textbf{bonne affaire / offre promotionnelle} ».
\item \eng{Credit} $\neq$ « crédit » (prêt bancaire) mais « \textbf{crédit d'appel / unité} » (airtime).
\item \eng{Balance} $\neq$ « balance » (instrument de pesée) mais « \textbf{solde} » (compte).
\item \eng{Package} $\neq$ « paquet » (colis) mais « \textbf{forfait / pack} » de services.
\item \eng{Validity} $\neq$ « validité » (qualité de valide) mais « \textbf{durée de validité} » (période).
\item \eng{Within} $\neq$ « dans » (localisation) mais « \textbf{sous / en moins de / au plus tard dans} » (délai).
\item \eng{To dial} $\neq$ « dialer » (argot) mais « \textbf{composer un numéro} ».
\end{itemize}
\end{attention}

\coursec{Méthode d'écoute : stratégies et exemple traité}

Cette section récapitule la démarche experte pour réussir l'épreuve de compréhension de l'oral (Listening) au Bac.

\begin{methode}[La check-list de l'auditeur expert (Avant / Pendant / Après)]
\begin{enumerate}
\item \textbf{AVANT (30 sec -- 1 min)} : Lire le titre, les consignes, les questions. Souligner les \textbf{mots-clés} à écouter (chiffres, noms propres, verbes d'action). Prédire le vocabulaire. Dessiner le tableau de notes.
\item \textbf{PENDANT (Écoute 1)} : Ne pas écrire. Juste écouter pour valider les hypothèses (Thème, Locuteurs, Lieu). Noter mentalement la structure (Intro $\rightarrow$ Détails $\rightarrow$ Conclusion / Appel à l'action).
\item \textbf{PENDANT (Écoute 2)} : Remplir le tableau de notes. \textbf{Écrire en chiffres + abréviations}. Si vous ratez un chiffre, laissez la case vide, continuez, ne bloquez pas.
\item \textbf{PENDANT (Écoute 3)} : Cibler \textbf{uniquement} les cases vides. Vérifier la cohérence (unités, ordre de grandeur). Noter les connecteurs logiques (\eng{but, however, also, after that}).
\item \textbf{APRÈS} : Transférer les réponses sur la copie finale. Phrases complètes pour les questions ouvertes. Citations exactes pour les justifications Vrai/Faux.
\end{enumerate}
\end{methode}

\begin{experience}[Entraînement chronométré : \eng{Speed Note-Taking}]
L'enseignant lit le Script 1 à vitesse normale une seule fois. Les élèves doivent remplir le tableau de notes en \textbf{temps réel} sans deuxième écoute. Comparaison en binôme puis correction collective. Objectif : automatiser la prise de notes chiffrées.
\end{experience}

\coursec{Production guidée : réemploi du lexique (Speaking \& Writing)}

Le but est de réactiver le vocabulaire et les structures (impératif, \eng{I'd like to}, passif, \eng{within}) dans des tâches communicatives réalistes.

\courssub{Expression orale (Speaking) : Jeu de rôle \eng{At the Call Centre}}

\begin{experience}[Jeu de rôle guidé : \eng{Subscribe to a Package}]
\textbf{Consigne} : Par binômes (A = Client, B = Agent Camtel/MTN/Orange/Nexttel). Durée : 3 min par rôle. Utilisez le lexique de la leçon.
\begin{enumerate}
\item \textbf{Préparation (2 min)} : L'élève A choisit un besoin (ex: « Je veux 20 Go pour 2 semaines, budget 5 000 FCFA »). L'élève B prépare ses 3 offres (inventées mais réalistes : prix, data, validité, code USSD).
\item \textbf{Interaction} :
      \begin{itemize}
      \item Agent : \eng{Good morning, [Operator] customer service. How may I help you?}
      \item Client : \eng{Good morning. I'd like to subscribe to a data package.}
      \item Agent : \eng{We have... [Offre 1], ... [Offre 2], ... [Offre 3].}
      \item Client : \eng{How much is the [Name] plan? / How long is it valid? / How do I subscribe?}
      \item Agent : \eng{It costs... / It is valid for... / To subscribe, dial...}
      \item Client : \eng{Thank you. I'll take the [Name] plan. Goodbye.}
      \item Agent : \eng{You're welcome. Thank you for calling [Operator]. Have a nice day.}
      \end{itemize}
\item \textbf{Échange des rôles} et répétition avec un autre besoin.
\item \textbf{Retour (Feedback)} : L'enseignant note au tableau les phrases bien formulées et corrige 2--3 erreurs fréquentes (ex: \eng{*I want} $\rightarrow$ \eng{I'd like}, \eng{*valid during} $\rightarrow$ \eng{valid for}, \eng{*activate} $\rightarrow$ \eng{activated}).
\end{enumerate}
\end{experience}

\courssub{Expression écrite (Writing) : Rédaction d'une annonce publicitaire}

\begin{methode}[Rédiger une annonce radio / SMS pour un forfait (Format court)]
\begin{enumerate}
\item \textbf{Accroche (Hook)} : Nom de l'opérateur + Nom de l'offre + \eng{Welcome to...} / \eng{Introducing...}
\item \textbf{Offre (Core)} : \eng{Get [Volume] of [Data/Calls/SMS] for [Price] FCFA.} (Impératif \eng{Get} / \eng{Enjoy}).
\item \textbf{Validité} : \eng{Valid for [Duration].}
\item \textbf{Appel à l'action (Call to Action)} : \eng{To subscribe, dial [USSD Code].} / \eng{Dial [Code] now!}
\item \textbf{Condition / Info légale} : \eng{Offer available for [Target] customers.} / \eng{Terms and conditions apply.}
\end{enumerate}
\end{methode}

\begin{exercice}[Production écrite : Votre propre forfait]
Imaginez que vous êtes chef marketing pour un opérateur fictif \eng{“CamConnect”}. Rédigez une annonce publicitaire (style radio, 30--40 mots) pour un nouveau forfait \eng{“Student Boost”} ciblant les lycéens et étudiants.
\textbf{Contraintes} :
\begin{itemize}
\item Prix : 1,000 FCFA.
\item Contenu : 3 GB Data + 100 SMS + 30 min calls (all networks).
\item Validité : 30 days.
\item Code USSD : *123\#.
\item Cible : \eng{prepaid students}.
\item Utilisez : \eng{Get, Valid for, Dial, Offer available to}.
\end{itemize}
\end{exercice}

\begin{corrige}[Exercice : Rédaction annonce \eng{Student Boost}]
\textbf{Proposition de réponse (modèle) :}
\begin{textebox}[Student Boost — CamConnect Advert]
“Welcome to CamConnect Student Boost! Get 3 GB of data, 100 SMS and 30 minutes of calls to all networks for only 1,000 FCFA. Valid for 30 days. To subscribe, dial *123\#. Offer available to all prepaid students. Terms and conditions apply.”
\end{textebox}
\textbf{Analyse} : Respect de la structure (Hook, Core, Validity, CTA, Condition). Vocabulaire précis. Impératifs corrects. Chiffres exacts. Registre publicitaire.
\end{corrige}

\courssub{Expression écrite : Email formel de réclamation (Complaint Email)}

\begin{methode}[Structure d'un email de réclamation formel (\eng{Formal Complaint Email})]
\begin{enumerate}
\item \textbf{Subject Line} : Clair, avec référence client/numéro de ligne. Ex: \eng{Subject: Complaint -- Activation Delay -- Line +237 6XX XX XX XX -- Ref: 12345}
\item \textbf{Salutation} : \eng{Dear Sir/Madam,} (ou \eng{Dear Customer Service Team,})
\item \textbf{Opening} : \eng{I am writing to complain about...} / \eng{I wish to express my dissatisfaction with...}
\item \textbf{Facts (Chronologique)} : Date, action effectuée (paiement, souscription), problème constaté (non-activation, débit lent, facturation erronée). Utilisez \eng{On [date], I subscribed to... / I paid... / However, ...}
\item \textbf{Impact} : \eng{This prevents me from... / This causes me great inconvenience.}
\item \textbf{Request (Action attendue)} : \eng{I would like you to...} (activate immediately / refund / investigate). \eng{I look forward to your prompt response.}
\item \textbf{Closing} : \eng{Yours faithfully,} (si \eng{Dear Sir/Madam}) + \eng{Full Name} + \eng{Phone Number} + \eng{Customer ID}.
\end{enumerate}
\end{methode}

\begin{exercice}[Production écrite : Email de réclamation]
Vous avez souscrit au forfait \eng{Camtel Family Plan} (20,000 FCFA) il y a 72 heures. Votre connexion n'est toujours pas activée. Vous avez payé via Mobile Money (Réf: MM123456). Rédigez un email formel à \eng{support@camtel.cm} pour demander l'activation immédiate ou le remboursement. (80--100 mots).
\end{exercice}

\begin{corrige}[Exercice : Email de réclamation]
\begin{textebox}[Email de réclamation -- Modèle]
Subject: Complaint -- Activation Delay -- Family Plan -- Line +237 6XX XX XX XX -- Ref: MM123456

Dear Sir/Madam,

I am writing to complain about the non-activation of my Camtel Family Plan subscription.

On [Date], I subscribed to the Family Plan (20,000 FCFA/month, 50 GB) and made the payment via Mobile Money (Reference: MM123456). According to your agent, activation should occur within 48 hours after payment. However, 72 hours have passed and my connection is still not active.

This delay prevents me from working remotely and accessing online educational resources. I would like you to activate my connection immediately or process a full refund.

I look forward to your prompt response.

Yours faithfully,
[Your Full Name]
Customer ID: [Number]
Phone: +237 6XX XX XX XX
\end{textebox}
\end{corrige}

\begin{savaistu}[Culture numérique : Les codes USSD au Cameroun et en Afrique]
Les codes \cle{USSD} (\eng{Unstructured Supplementary Service Data}) sont la colonne vertébrale des services mobiles en Afrique, là où la data est chère ou le réseau 3G/4G instable. Ils fonctionnent sur \textbf{tout téléphone} (même un \eng{feature phone} / \eng{kabato} à 5 000 FCFA), sans internet.
\begin{itemize}
\item \textbf{Syntaxe} : \eng{*[Code]\#} (ex: \eng{*157\#}). À l'oral : \eng{star one five seven hash} (UK) ou \eng{pound} (US).
\item \textbf{Usages} : Souscription forfaits (\eng{*157\#}), Solde (\eng{*157*1\#}), Transfert crédit (\eng{*150\#}), Mobile Money (\eng{*126\#} MTN MoMo, \eng{*150\#} Orange Money), \eng{Call me back} (\eng{*121*Num\#}).
\item \textbf{Session} : Une session USSD a un timeout court (souvent 20 sec). Il faut répondre vite (chiffre 1, 2, 3...).
\item \textbf{Interopérabilité} : Au Cameroun, la portabilité du numéro permet de garder son numéro en changeant d'opérateur, mais les codes USSD changent.
\end{itemize}
\textbf{Anecdote} : Le \eng{*157\#} de MTN est si connu qu'il est devenu un verbe : \eng{“Did you 157 the bundle?”} (As-tu souscrit le forfait via le 157 ?).
\end{savaistu}

\begin{savaistu}[Le Cameroun et la connectivité : Enjeux actuels]
Le Cameroun compte plus de 20 millions d'abonnés mobiles (taux de pénétration > 80\%). Les deux opérateurs historiques \textbf{MTN Cameroon} et \textbf{Orange Cameroun} se partagent le marché mobile, tandis que \textbf{Camtel} (public) domine la fibre optique et l'internet fixe (ADSL/Fibre). L'arrivée de la 4G (2016) puis de la 5G (tests en cours à Yaoundé/Douala) transforme les usages : streaming, e-learning, e-commerce (\eng{Jumia, Glovo}), Mobile Money (\eng{MoMo, Orange Money}). Le coût de la data reste un sujet de débat public (\eng{“Data is expensive”}). Le gouvernement a lancé le projet \eng{“Cameroun Digital 2020”} pour baisser les coûts et étendre la couverture rurale. Comprendre les offres (\eng{bundles}) en anglais est donc une compétence citoyenne et professionnelle majeure.
\end{savaistu}

\coursec{Compréhension : questions variées}

Après avoir découvert le script d'écoute de la section précédente — l'annonce publicitaire d'une compagnie de téléphonie mobile présentant ses forfaits \eng{Mega Bundle} et \eng{Family Plan} — nous passons maintenant à l'étape décisive : \textbf{vérifier que nous avons réellement compris}. C'est exactement ce que l'on vous demandera au baccalauréat : écouter (ou lire) un document, puis répondre à des questions de compréhension de types variés. Dans cette section, nous allons nous entraîner sur les quatre formats les plus fréquents : le \textbf{vrai/faux justifié}, le \textbf{QCM}, les \textbf{questions ouvertes à réponse courte} (WH-questions) et la \textbf{complétion d'un tableau récapitulatif}.

\courssub{Le vrai/faux justifié : prouver ce que l'on affirme}

\begin{definition}[True or False?]
L'exercice \eng{True or False} consiste à lire une affirmation et à décider si elle est \cle{vraie} (\eng{true}) ou \cle{fausse} (\eng{false}) par rapport au document. Au baccalauréat, la consigne exige presque toujours de \textbf{justifier} la réponse en citant une phrase ou un passage du texte : \eng{Justify your answer with a sentence from the text}.
\end{definition}

\begin{textebox}[True or False? — Justify your answers]
1. The Mega Bundle costs 5,000 FCFA. \textbf{(F)} — “The Mega Bundle is yours for only 2,500 FCFA.”

2. The offer includes unlimited data. \textbf{(F)} — “You get 5 GB of internet data.”

3. Customers can check their balance. \textbf{(T)} — “To check your balance at any time, simply dial *157\#.”

4. The offer is valid for one month. \textbf{(F)} — “This bundle is valid for seven days.”
\end{textebox}

\begin{methode}[Réussir un vrai/faux justifié]
\begin{enumerate}
\item \textbf{Lisez d'abord toutes les affirmations} avant l'écoute : elles vous indiquent quels détails surveiller (prix, durée, contenu de l'offre).
\item Pendant l'écoute, \textbf{notez les chiffres et les mots-clés} qui correspondent à chaque affirmation.
\item Répondez \textbf{T} ou \textbf{F}, puis \textbf{soulignez dans le script la phrase exacte} qui prouve votre réponse.
\item Rédigez la justification entre guillemets, en recopiant fidèlement le texte : ne paraphrasez pas, citez.
\item Au Bac, \textbf{une réponse sans justification perd des points}, même si elle est correcte. La justification vaut souvent la moitié de la note.
\end{enumerate}
\end{methode}

\begin{attention}[Les affirmations « presque vraies »]
Les concepteurs d'examens adorent les affirmations qui reprennent les mots du texte en changeant \textbf{un seul détail} : \eng{The Mega Bundle costs 5,000 FCFA} ressemble beaucoup à la réalité (2,500 FCFA), mais elle est fausse. Autre piège classique : \eng{unlimited data} (« données illimitées ») alors que le texte dit précisément \eng{5 GB}. Ne vous fiez jamais à votre impression générale : confrontez chaque mot de l'affirmation au texte.
\end{attention}

\courssub{Le QCM : choisir entre plusieurs options}

\begin{definition}[Multiple Choice Questions]
Dans un \cle{QCM} (\eng{multiple choice questions}), chaque question propose trois ou quatre options (a, b, c, d). Une seule est correcte ; les autres sont des \cle{distracteurs}, c'est-à-dire de fausses réponses plausibles, souvent construites à partir d'autres chiffres ou noms présents dans le texte.
\end{definition}

\begin{textebox}[Multiple Choice Questions]
1. The Family Plan costs:
a) 10,000 FCFA \quad b) 20,000 FCFA \quad c) 35,000 FCFA

2. Activation of the Family Plan takes:
a) 24 hours \quad b) 48 hours \quad c) one week

3. To subscribe to the Mega Bundle, customers dial:
a) *157\# \quad b) *150\# \quad c) *175\#

\textbf{Answers:} 1. c) — “For just 35,000 FCFA, the Family Plan connects your whole household.” / 2. b) — “Activation takes only 48 hours.” / 3. a) — “Dial *157\# to subscribe.”
\end{textebox}

\begin{attention}[Méfiez-vous des distracteurs numériques]
Dans les QCM portant sur les prix et les durées, le texte mentionne presque toujours \textbf{plusieurs chiffres} : ici 10,000 FCFA (le forfait \eng{Student Pack}), 20,000 FCFA (le forfait \eng{Business Pack}) et 35,000 FCFA (le \eng{Family Plan}). Si vous n'avez retenu qu'un seul chiffre, vous avez une chance sur trois de tomber dans le piège. \textbf{Stratégie} : pendant l'écoute, notez \emph{tous} les chiffres entendus et associez immédiatement chacun à la bonne offre, sous forme de petit tableau au brouillon.
\end{attention}

\courssub{Les questions ouvertes : WH-questions et réponses courtes}

Les questions ouvertes commencent par un mot interrogatif (\cle{WH-words}) et attendent une \textbf{réponse courte mais complète}, construite à partir du texte. Rappelons les principaux mots interrogatifs et ce qu'ils demandent :

\begin{center}
\begin{tabularx}{\linewidth}{|l|l|X|X|}
\hline
\rowcolor{popblueL} \textbf{WH-word} & \textbf{Français} & \textbf{Demande…} & \textbf{Exemple de réponse} \\
\hline
Who & Qui & une personne, un groupe & \eng{Students and young professionals.} \\
\hline
What & Quoi / Que & une chose, une information & \eng{5 GB of data, 200 SMS and 60 minutes of calls.} \\
\hline
How much & Combien (prix, quantité indénombrable) & un prix, une quantité & \eng{It costs 2,500 FCFA.} \\
\hline
How many & Combien (dénombrable) & un nombre & \eng{200 SMS.} \\
\hline
How long & Combien de temps & une durée & \eng{For seven days.} \\
\hline
How & Comment & une manière, un procédé & \eng{By dialling *157\#.} \\
\hline
Where & Où & un lieu & \eng{At any agency in Douala or Yaoundé.} \\
\hline
When & Quand & un moment & \eng{Within 48 hours.} \\
\hline
Why & Pourquoi & une raison & \eng{Because it is cheaper than buying data daily.} \\
\hline
\end{tabularx}
\end{center}

\begin{propriete}[Rédiger une réponse courte correcte]
\begin{itemize}
\item \textbf{Reprenez la structure de la question} dans la réponse : \eng{How much does the Mega Bundle cost?} $\rightarrow$ \eng{The Mega Bundle costs 2,500 FCFA.} (et non simplement « 2,500 FCFA » si une phrase complète est exigée).
\item Respectez l'\textbf{ordre anglais} : sujet + verbe + complément. Ne dites jamais \eng{*Costs 2,500 FCFA} sans sujet.
\item \eng{How much} s'emploie pour les prix et les quantités indénombrables ; \eng{how many} pour les noms dénombrables : \eng{How many SMS are included?} — \eng{Two hundred.}
\end{itemize}
\end{propriete}

\begin{exemplebox}[WH-questions sur le script]
\eng{1. Who is the Mega Bundle designed for?} — \eng{It is designed for students and young people who use social media every day.} (« Pour qui le Mega Bundle est-il conçu ? »)

\eng{2. What does the Mega Bundle include?} — \eng{It includes 5 GB of internet data, 200 SMS and 60 minutes of calls.} (« Que comprend le Mega Bundle ? »)

\eng{3. How long is the bundle valid?} — \eng{It is valid for seven days.} (« Combien de temps le forfait est-il valable ? »)

\eng{4. How can customers check their balance?} — \eng{They can check their balance by dialling *157\#.} (« Comment les clients peuvent-ils consulter leur solde ? »)
\end{exemplebox}

\courssub{Complétion d'un tableau récapitulatif (grid completion)}

La \cle{complétion de tableau} est un exercice très fréquent : on vous donne une grille partiellement remplie et vous devez la compléter avec les informations du document. C'est un excellent entraînement car il vous oblige à \textbf{organiser} les informations entendues.

\begin{exoresolu}[Complete the grid with information from the script]
Énoncé — Complete the table about the Mega Bundle: price, data, SMS and calls, validity, subscription code.\\[2mm]
\textit{Grille à compléter :} Mega Bundle | ? FCFA | ? GB | ? SMS + ? min | ? days | *?\#
\tcblower
Solution détaillée — En réécoutant (ou relisant) le script, on relève : « \eng{The Mega Bundle is yours for only 2,500 FCFA} » (prix) ; « \eng{You get 5 GB of internet data} » (données) ; « \eng{200 SMS and 60 minutes of calls to all networks} » (SMS et appels) ; « \eng{valid for seven days} » (validité) ; « \eng{dial *157\# to subscribe} » (code).\\[2mm]
\textbf{Tableau complété :} Mega Bundle | \textbf{2,500 FCFA} | \textbf{5 GB} | \textbf{200 SMS + 60 min} | \textbf{7 days} | \textbf{*157\#}\\[2mm]
\textit{Astuce :} chaque case correspond à un mot-clé du script. Si une case reste vide, c'est que vous avez manqué une information — réécoutez le passage concerné.
\end{exoresolu}

\begin{popfigure}
\begin{center}
\begin{tabular}{|l|c|c|c|c|c|}
\hline
\rowcolor{popgreenL} \textbf{Package} & \textbf{Price} & \textbf{Data} & \textbf{SMS + Calls} & \textbf{Validity} & \textbf{Code} \\
\hline
Mega Bundle & 2,500 FCFA & 5 GB & 200 SMS + 60 min & 7 days & *157\# \\
\hline
\end{tabular}
\end{center}
\legende{Tableau récapitulatif complété en correction collective : les informations essentielles du Mega Bundle.}
\end{popfigure}

\begin{popfigure}
\begin{center}
\begin{tikzpicture}[mindmap, grow cyclic, every node/.style={concept}, concept color=popblueL,
root concept/.append style={concept color=popblue, text=white, font=\bfseries},
level 1/.append style={level distance=3.2cm, sibling angle=72},
level 2/.append style={level distance=2.2cm, sibling angle=45}]
\node[root concept]{Mega Bundle}
  child[concept color=poporangeL] { node {Price} child { node {2,500 FCFA} } }
  child[concept color=popgreenL] { node {Data} child { node {5 GB} } }
  child[concept color=poppurpleL] { node {Calls} child { node {60 min} } }
  child[concept color=poppinkL] { node {SMS} child { node {200 SMS} } }
  child[concept color=popgoldL] { node {Validity} child { node {7 days} } };
\end{tikzpicture}
\end{center}
\legende{Carte mentale de l'offre : une stratégie visuelle pour retenir les cinq informations-clés pendant l'écoute.}
\end{popfigure}

\courssub{Correction collective : justifier en citant le script}

La dernière étape se déroule en classe entière : les élèves donnent leurs réponses à voix haute et \textbf{citent le script} pour les justifier. C'est un moment d'expression orale autant que de compréhension.

\begin{methode}[Participer à la correction collective]
\begin{enumerate}
\item Annoncez d'abord votre réponse clairement : \eng{Number one is false.}
\item Citez ensuite la preuve : \eng{The script says: “The Mega Bundle is yours for only 2,500 FCFA.”}
\item Si un camarade n'est pas d'accord, défendez votre réponse : \eng{I think it is true because the speaker says…}
\item En cas d'erreur, repérez \textbf{pourquoi} vous vous êtes trompé : distracteur numérique, mot mal entendu, affirmation « presque vraie ».
\end{enumerate}
\end{methode}

\begin{exercice}[Your turn — Mixed comprehension]
Répondez aux questions suivantes sur le script de la section 2.
\begin{enumerate}
\item True or False? Justify with a sentence from the text.\\
a) The Family Plan connects the whole household.\\
b) The Mega Bundle is valid for thirty days.\\
c) Customers can subscribe by dialling a code.
\item Choose the correct answer. The Mega Bundle includes: a) 2 GB b) 5 GB c) 10 GB of data.
\item Answer with a full sentence.\\
a) How much does the Mega Bundle cost?\\
b) How long does activation of the Family Plan take?\\
c) How can customers check their balance?
\end{enumerate}
\end{exercice}

\begin{corrige}[Exercice « Your turn »]
\begin{enumerate}
\item a) \textbf{T} — “The Family Plan connects your whole household.” b) \textbf{F} — “This bundle is valid for seven days.” c) \textbf{T} — “Dial *157\# to subscribe.”
\item b) 5 GB — “You get 5 GB of internet data.”
\item a) \eng{The Mega Bundle costs 2,500 FCFA.} b) \eng{Activation of the Family Plan takes 48 hours.} c) \eng{Customers can check their balance by dialling *157\#.}
\end{enumerate}
\end{corrige}

\begin{aretenir}[L'essentiel de la section]
\begin{itemize}
\item \textbf{Vrai/faux} : toujours justifier en citant une phrase exacte du texte, entre guillemets.
\item \textbf{QCM} : attention aux distracteurs numériques ; notez tous les chiffres et associez-les à la bonne offre.
\item \textbf{WH-questions} : répondez par une phrase complète qui reprend la structure de la question ; \eng{how much} (prix, indénombrable) ≠ \eng{how many} (dénombrable).
\item \textbf{Grid completion} : chaque case du tableau correspond à un mot-clé du script.
\end{itemize}
\end{aretenir}

\coursec{Lexique thématique : télécommunications et forfaits}

Après avoir répondu aux questions de compréhension sur le script d'écoute, arrêtons-nous maintenant sur les mots qui ont permis de comprendre le texte. Une bonne écoute repose d'abord sur un bon vocabulaire : si vous connaissez les termes des télécommunications, vous reconnaîtrez instantanément les informations essentielles (prix, durée, contenu du forfait) dans n'importe quel enregistrement. Organisons ce lexique en cinq champs, comme les branches d'une carte mentale.

\begin{popfigure}
\begin{center}
\begin{tikzpicture}[mindmap, grow cyclic,
  every node/.style={concept, text=white, font=\small\bfseries},
  root concept/.append style={concept color=popblue, font=\bfseries},
  level 1/.append style={level distance=3.4cm, sibling angle=72},
  level 1 concept/.append style={concept color=poporange}]
\node[root concept]{Telecom vocabulary}
  child { node {Subscribing and activating} }
  child { node {Package contents} }
  child { node {Prices and payment} }
  child { node {Duration and validity} }
  child { node {Problems and customer care} };
\end{tikzpicture}
\end{center}
\legende{Carte mentale : les cinq champs lexicaux des télécommunications}
\end{popfigure}

\courssub{Champ 1 : souscription et activation}

Observons d'abord ces phrases tirées de situations réelles au Cameroun :

\eng{To subscribe to the weekly bundle, dial *123\# and follow the instructions.} « Pour souscrire au forfait hebdomadaire, composez le *123\# et suivez les instructions. »

\eng{You must register your SIM card before you can activate any service.} « Vous devez enregistrer votre carte SIM avant de pouvoir activer un service. »

\begin{definition}[Subscribing and activating]
\begin{itemize}
\item \cle{to subscribe (to)} : souscrire, s'abonner (à). Attention à la préposition : on dit \eng{to subscribe \textbf{to} a package}, jamais \eng{to subscribe a package}.
\item \cle{to activate} : activer (un service, une carte SIM). Nom : \eng{activation}.
\item \cle{to dial (a code)} : composer (un code, un numéro). \eng{Dial *123\# to check your balance.}
\item \cle{to register} : enregistrer, s'enregistrer. \eng{Register your SIM card with your ID card.}
\item \cle{a subscription} : un abonnement, une souscription.
\end{itemize}
\end{definition}

\begin{attention}[Piège : la préposition]
En français, on « souscrit un forfait » ou on « s'abonne à un forfait ». En anglais, le verbe \eng{to subscribe} exige toujours la préposition \eng{to} : \eng{I subscribed to a monthly plan.} « Je me suis abonné à un forfait mensuel. » Ne dites jamais \eng{I subscribed a plan}.
\end{attention}

\courssub{Champ 2 : contenu des forfaits}

\begin{definition}[Bundle]
Un \cle{bundle} est un forfait groupé combinant plusieurs services (données + SMS + appels) à un prix unique. Exemple : \eng{a weekly bundle} = un forfait hebdomadaire ; \eng{a data bundle} = un forfait internet.
\end{definition}

\begin{definition}[Package contents]
\begin{itemize}
\item \cle{data} : les données internet (mégaoctets, gigaoctets). \eng{This bundle gives you 2 GB of data.}
\item \cle{airtime / credit} : le crédit d'appel. \eng{I have no airtime left.}
\item \cle{SMS / text messages} : les SMS. \eng{The plan includes 100 free SMS.}
\item \cle{minutes} : les minutes d'appel. \eng{You get 50 minutes of calls to all networks.}
\item \cle{a plan} : un forfait, une formule (souvent mensuelle, par abonnement). \eng{a prepaid plan} (prépayé), \eng{a postpaid plan} (postpayé).
\end{itemize}
\end{definition}

\courssub{Champ 3 : prix et paiement}

\begin{definition}[Prices and payment]
\begin{itemize}
\item \cle{a fee} : des frais, un tarif. \eng{There is no subscription fee.} « Il n'y a pas de frais de souscription. »
\item \cle{a rate} : un taux, un tarif (par unité). \eng{The call rate is 50 francs per minute.}
\item \cle{to top up / to recharge} : recharger (son téléphone, son compte). \eng{I need to top up my phone.}
\item \cle{a balance} : un solde. \eng{Check your balance by dialling *155\#.}
\item \cle{a payment} : un paiement. \eng{You can make the payment by mobile money.}
\end{itemize}
\end{definition}

\begin{exemplebox}[Exemples traduits]
\begin{itemize}
\item \eng{I need to top up my phone before I call my uncle in Bamenda.} « Je dois recharger mon téléphone avant d'appeler mon oncle à Bamenda. »
\item \eng{What is the rate for international calls?} « Quel est le tarif des appels internationaux ? »
\item \eng{My balance is too low to buy a data bundle.} « Mon solde est trop bas pour acheter un forfait internet. »
\item \eng{She paid the monthly fee by mobile money.} « Elle a payé les frais mensuels par mobile money. »
\end{itemize}
\end{exemplebox}

\courssub{Champ 4 : durée et validité}

\begin{definition}[Duration and validity]
\begin{itemize}
\item \cle{validity} : la validité (durée pendant laquelle le forfait est utilisable). \eng{The validity of this bundle is seven days.}
\item \cle{to expire} : expirer, arriver à expiration. \eng{My bundle has expired.}
\item \cle{to renew} : renouveler. \eng{You can renew your plan automatically.}
\item \cle{unlimited} : illimité. \eng{unlimited calls} « appels illimités » ; \eng{unlimited data} « données illimitées ».
\end{itemize}
\end{definition}

\begin{exemplebox}[Exemples traduits]
\begin{itemize}
\item \eng{My bundle has expired, so I cannot browse the internet.} « Mon forfait a expiré, donc je ne peux pas naviguer sur internet. »
\item \eng{This offer is valid for thirty days.} « Cette offre est valable trente jours. »
\item \eng{If you do not renew your subscription, the service will stop.} « Si vous ne renouvelez pas votre abonnement, le service s'arrêtera. »
\end{itemize}
\end{exemplebox}

\courssub{Champ 5 : problèmes et service client}

\begin{definition}[Problems and customer care]
\begin{itemize}
\item \cle{network coverage} : la couverture réseau. \eng{There is no network coverage in the village.}
\item \cle{a connection} : une connexion. \eng{The connection is very slow today.}
\item \cle{customer care / customer service} : le service client. \eng{Call customer care on 800.}
\item \cle{a complaint} : une plainte, une réclamation. Verbe : \eng{to complain (about)}. \eng{She made a complaint about the poor network.}
\item \cle{a network failure} : une panne de réseau.
\end{itemize}
\end{definition}

\begin{exemplebox}[Exemples traduits]
\begin{itemize}
\item \eng{There is no network coverage in the village, so we cannot make calls.} « Il n'y a pas de couverture réseau au village, donc nous ne pouvons pas passer d'appels. »
\item \eng{I complained to customer care because my data disappeared.} « Je me suis plaint au service client parce que mes données ont disparu. »
\item \eng{The internet connection is down in the whole neighbourhood.} « La connexion internet est en panne dans tout le quartier. »
\end{itemize}
\end{exemplebox}

\courssub{Collocations fréquentes et tableau récapitulatif}

En anglais, certains verbes vont obligatoirement avec certains noms : ce sont les \cle{collocations}. Les apprendre par blocs, et non mot par mot, est la clé d'une expression naturelle.

\begin{center}
\begin{tabularx}{\linewidth}{|X|X|X|}
\hline
\rowcolor{popblueL} Collocation anglaise & Traduction française & Exemple \\
\hline
to subscribe to a package & souscrire à un forfait & \eng{She subscribed to a monthly package.} \\
\hline
to top up your phone & recharger son téléphone & \eng{He topped up his phone with 1,000 francs.} \\
\hline
to check your balance & consulter son solde & \eng{Dial *155\# to check your balance.} \\
\hline
to run out of credit & être à court de crédit & \eng{I have run out of credit.} \\
\hline
to dial a code & composer un code & \eng{Dial the code to activate the bundle.} \\
\hline
to make a complaint & déposer une plainte & \eng{They made a complaint about the fees.} \\
\hline
to renew a subscription & renouveler un abonnement & \eng{I renewed my subscription yesterday.} \\
\hline
\end{tabularx}
\end{center}

\begin{attention}[Faux amis et pièges lexicaux]
\begin{itemize}
\item \eng{data} se prononce « déï-teu » et désigne les \textbf{données internet}, pas une date. « Une date » se dit \eng{a date}.
\item \eng{airtime} signifie ici « crédit d'appel », et non « temps d'antenne » (sens réservé aux médias).
\item \eng{to expire} = arriver à expiration ; « la date d'expiration » se dit \eng{the expiry date} (et non « expiration date » en anglais britannique).
\item \eng{to recharge} : on recharge un téléphone ou un compte, mais « recharger une batterie » se dit \eng{to charge} (ou \eng{to recharge a battery}).
\item Ne confondez pas \eng{actually} (« en fait ») et « actuellement » (\eng{currently, at present}).
\end{itemize}
\end{attention}

\begin{savaistu}[Le même objet, des mots différents]
Au Nigeria et au Ghana, la carte de recharge s'appelle \eng{a recharge card} ; au Royaume-Uni, on parle plutôt de \eng{a top-up voucher}. Le verbe \eng{to flash someone} (bipper quelqu'un, lui faire un appel bref pour qu'il rappelle) est typique de l'anglais africain : un Britannique dirait plutôt \eng{to give someone a missed call}. Ces différences montrent que l'anglais est une langue plurielle, adaptée à chaque région du monde.
\end{savaistu}

\begin{textebox}[Choosing the right package]
Amina is a student in Yaoundé. Every Monday morning, she tops up her phone with 2,000 francs and subscribes to a weekly bundle. The bundle gives her 1.5 GB of data, 30 minutes of calls and 50 SMS. To activate it, she simply dials a short code and confirms the payment from her balance. Last week, however, Amina had a problem. Her bundle expired on Sunday night, but she did not notice it. On Monday, she browsed the internet for two hours using her airtime instead of her data bundle. When she checked her balance, all her credit was gone. She was very disappointed and decided to complain. She called customer care and explained the situation politely. The agent told her that she could activate an automatic renewal option, so that her bundle would renew itself every week. He also advised her to dial the balance code regularly to avoid running out of credit without warning. Amina thanked him and activated the option immediately. Now she never worries about her subscription any more. She even taught her younger brother how to check the validity of his own plan and how to compare the rates of different packages before choosing one. In her neighbourhood, many people still buy expensive daily plans because they do not know that weekly bundles are much cheaper. Amina enjoys explaining to them how mobile services work.
\end{textebox}

\textbf{Glossaire :} \emph{to top up} (v., recharger) ; \emph{bundle} (n., forfait groupé) ; \emph{to expire} (v., expirer) ; \emph{airtime} (n., crédit d'appel) ; \emph{balance} (n., solde) ; \emph{customer care} (n., service client) ; \emph{renewal} (n., renouvellement) ; \emph{rate} (n., tarif).

\textbf{Quick check :} 1) What does Amina's weekly bundle include? 2) Why did she lose all her credit? 3) What solution did the customer care agent suggest?

\begin{exoresolu}[Complétez avec le mot qui convient]
Complétez les phrases avec : \eng{balance, expired, top up, coverage, dial, unlimited}.
\begin{enumerate}
\item \eng{Please \ldots\ *123\# to activate the bundle.}
\item \eng{There is no network \ldots\ in my village.}
\item \eng{I must \ldots\ my phone before calling Douala.}
\item \eng{My data bundle has \ldots\ ; I need a new one.}
\item \eng{This plan offers \ldots\ SMS to all networks.}
\item \eng{Check your \ldots\ before you buy a package.}
\end{enumerate}
\tcblower
\textbf{Solution détaillée :}
\begin{enumerate}
\item \eng{dial} : on « compose » un code (\eng{to dial a code}).
\item \eng{coverage} : \eng{network coverage} = couverture réseau (collocation obligatoire).
\item \eng{top up} : \eng{to top up one's phone} = recharger son téléphone.
\item \eng{expired} : participe passé de \eng{to expire}, avec l'auxiliaire \eng{has} (present perfect).
\item \eng{unlimited} : adjectif placé devant le nom \eng{SMS} ; « SMS illimités ».
\item \eng{balance} : \eng{to check your balance} = consulter son solde.
\end{enumerate}
\end{exoresolu}

\begin{exercice}[À vous de jouer]
Traduisez en anglais : 1) « Je dois recharger mon téléphone. » 2) « Mon forfait a expiré hier. » 3) « Il n'y a pas de couverture réseau au village. » 4) « Composez ce code pour vérifier votre solde. » 5) « Ce forfait hebdomadaire comprend des appels illimités. »
\end{exercice}

\begin{corrige}[Exercice : traduction]
1) \eng{I need to top up my phone.} 2) \eng{My bundle expired yesterday.} 3) \eng{There is no network coverage in the village.} 4) \eng{Dial this code to check your balance.} 5) \eng{This weekly bundle includes unlimited calls.} Attention : n'oubliez pas la préposition dans \eng{to subscribe \textbf{to}}, et n'écrivez jamais \eng{expired date} pour « date d'expiration » (\eng{expiry date}).
\end{corrige}

\begin{aretenir}[L'essentiel du lexique]
Retenez les cinq champs et leurs collocations : \eng{subscribe to a package}, \eng{activate a service}, \eng{dial a code}, \eng{top up your phone}, \eng{check your balance}, \eng{run out of credit}, \eng{a bundle expires}, \eng{renew a subscription}, \eng{network coverage}, \eng{call customer care}, \eng{make a complaint}. Ces blocs de mots vous permettront de comprendre n'importe quel enregistrement sur les forfaits et de réemployer le lexique dans vos propres phrases.
\end{aretenir}

\coursec{Méthode d'écoute : stratégies et exemple traité}

Vous venez d'apprendre le vocabulaire des télécommunications et des forfaits : \eng{data bundle}, \eng{prepaid customer}, \eng{validity}, \eng{dial}, \eng{top up}. Mais connaître les mots ne suffit pas : à l'écoute, tout va très vite, on ne peut pas revenir en arrière, et le stress fait perdre des informations pourtant simples. Cette section vous donne une \cle{méthode d'écoute} en quatre stratégies, puis un exemple traité pas à pas sur le Script 1 (l'annonce \eng{Mega Bundle}).

\courssub{Pourquoi une méthode d'écoute ?}

Écouter en anglais n'est pas comme lire. À la lecture, vous pouvez relire une phrase, chercher un mot, prendre votre temps. À l'écoute, le son disparaît aussitôt prononcé. La bonne nouvelle : pour réussir une épreuve de compréhension orale, il n'est \emph{pas nécessaire de tout comprendre}. Il faut comprendre \emph{ce qui est utile} : le thème, les chiffres, les noms propres, les réponses aux questions. Toute la méthode qui suit repose sur cette idée : \textbf{écouter avec un objectif, pas écouter pour tout saisir}.

\begin{propriete}[Compréhension orale au Bac : ce qui est attendu]
Au baccalauréat, la compréhension (orale ou écrite) est évaluée par des \cle{questions objectives} : vrai/faux, QCM, réponses courtes. Une réponse courte et exacte — un chiffre, un nom, une date — suffit presque toujours. Inutile de rédiger de longues phrases : \eng{2,500 FCFA} est une réponse parfaite à \eng{How much does the bundle cost?}. Écrire \eng{The bundle costs two thousand five hundred francs CFA, which is a very good price for customers} n'apporte rien et fait perdre du temps. Allez à l'essentiel, en anglais correct.
\end{propriete}

\courssub{Stratégie 1 : identifier le type de texte dès les premières secondes}

Les cinq à dix premières secondes d'un enregistrement sont décisives. Elles vous disent \emph{à quel genre de texte} vous avez affaire, et chaque genre a ses propres codes.

\begin{center}
\begin{tabularx}{\linewidth}{|X|X|X|}
\hline
\rowcolor{popblueL} Type de texte & Signaux reconnaissables & Ce qu'on attend \\
\hline
Annonce publicitaire (\eng{advert}) & voix enthousiaste, musique, \eng{Get...!}, \eng{only...!} & nom de l'offre, prix, contenu, validité, code \\
\hline
Dialogue en agence (\eng{dialogue}) & deux voix, questions/réponses, \eng{How can I help you?} & le problème du client, la solution proposée \\
\hline
Message automatique (\eng{recorded message}) & voix neutre, \eng{Press 1 to...}, \eng{Please hold the line} & les options, les numéros à taper, les horaires \\
\hline
\end{tabularx}
\end{center}

Dès que vous entendez \eng{Hello! Get 5 GB of data...}, vous savez que c'est une \cle{annonce} : votre cerveau doit immédiatement se préparer à capter un \emph{nom d'offre}, un \emph{prix}, un \emph{contenu}, une \emph{validité} et un \emph{code}. C'est exactement la structure d'une annonce de forfait, au Cameroun comme ailleurs.

\courssub{Stratégie 2 : repérer les mots porteurs de sens}

Une phrase anglaise contient deux sortes de mots : les \cle{mots porteurs de sens} (noms, verbes, adjectifs, chiffres) et les \cle{mots-outils} (articles, prépositions, auxiliaires). À l'oral, les mots porteurs de sens sont \emph{accentués} ; les mots-outils sont avalés, à peine prononcés. Bonne nouvelle : ce sont précisément les mots accentués qui comptent pour la compréhension.

\begin{exemplebox}[Les mots-clés d'une phrase d'annonce]
Écoutez (ou lisez) cette phrase extraite du Script 1 :

\eng{Get 5 GB of data, 200 SMS and 60 minutes for only 2,500 FCFA, valid for 7 days.}

Les mots-clés — les seuls indispensables — sont :

\begin{center}
\begin{tabular}{|l|l|l|l|l|}
\hline
\rowcolor{popgreenL} 5 GB & 200 SMS & 60 minutes & 2,500 FCFA & 7 days \\
\hline
\end{tabular}
\end{center}

Tout le reste (\eng{get}, \eng{of}, \eng{for}, \eng{only}, \eng{valid}) est secondaire : il relie les informations mais n'en contient aucune. Si vous avez noté ces cinq éléments, vous pouvez répondre à toutes les questions de contenu de l'annonce.
\end{exemplebox}

\courssub{Stratégie 3 : noter les chiffres en chiffres arabes}

Pendant l'écoute, chaque seconde compte. Si vous écrivez \eng{two thousand five hundred} en toutes lettres, vous ratez la phrase suivante. Notez \textbf{2500}. Les chiffres arabes sont universels, rapides, et acceptés dans les réponses.

\begin{attention}[Fifteen ou fifty ? Le piège des nombres à l'oral]
À l'oral, certains nombres anglais se ressemblent dangereusement : \eng{fifteen} (15) et \eng{fifty} (50), \eng{thirteen} (13) et \eng{thirty} (30), \eng{sixteen} (16) et \eng{sixty} (60). Le moyen sûr de les distinguer est \cle{l'accent tonique} :
\begin{itemize}
\item pour les nombres en \eng{-teen}, l'accent tombe sur la \textbf{fin} : fif-\textbf{TEEN}, thir-\textbf{TEEN} ;
\item pour les nombres en \eng{-ty}, l'accent tombe sur le \textbf{début} : \textbf{FIF}-ty, \textbf{THIR}-ty.
\end{itemize}
Autre piège : les prix. \eng{2,500 FCFA} se dit \eng{two thousand five hundred} ou \eng{twenty-five hundred}. Entraînez-vous à noter directement 2500, quelle que soit la formulation.
\end{attention}

\courssub{Stratégie 4 : la deuxième écoute sert à confirmer}

Beaucoup d'élèves « réécoutent tout » à la deuxième écoute, comme la première fois. C'est une erreur. La deuxième écoute a une fonction précise : \cle{vérifier} les réponses déjà trouvées et combler \emph{uniquement} les trous.

\begin{enumerate}
\item Première écoute : identifier le type de texte, noter les mots-clés et les chiffres, répondre à ce qui est clair.
\item Entre les deux écoutes : relire vos notes, repérer ce qui manque (par exemple : « je n'ai pas le code de souscription »).
\item Deuxième écoute : tendre l'oreille \emph{seulement} vers l'information manquante, et confirmer mentalement chaque réponse déjà notée.
\end{enumerate}

\begin{popfigure}
\begin{center}
\begin{tikzpicture}[node distance=0.55cm and 0.9cm,
etape/.style={rectangle, rounded corners=4pt, draw=popblue, fill=popblueL, thick, align=center, minimum width=3.1cm, minimum height=1.1cm, font=\small},
fleche/.style={-{Stealth[length=3mm]}, thick, popdark}]
\node[etape, align=center] (e1){\textbf{1. Type de texte}\\ annonce, dialogue,\\ message automatique};
\node[etape, right=of e1, fill=poporangeL, draw=poporange, align=center] (e2){\textbf{2. Mots-clés}\\ noms, verbes, chiffres\\ (ignorer les mots-outils)};
\node[etape, right=of e2, fill=popgreenL, draw=popgreen, align=center] (e3){\textbf{3. Chiffres notés}\\ en chiffres arabes\\ 2500, pas en lettres};
\node[etape, right=of e3, fill=poppurpleL, draw=poppurple, align=center] (e4){\textbf{4. Vérification}\\ 2\textsuperscript{e} écoute :\\ confirmer, combler};
\draw[fleche] (e1) -- (e2);
\draw[fleche] (e2) -- (e3);
\draw[fleche] (e3) -- (e4);
\end{tikzpicture}
\end{center}
\legende{Organigramme de la méthode d'écoute en quatre étapes}
\end{popfigure}

\courssub{Exemple traité pas à pas : le Script 1 (annonce Mega Bundle)}

Appliquons maintenant les quatre stratégies au Script 1 de cette leçon. Lisez-le d'abord comme si vous l'entendiez.

\begin{textebox}[Script 1 — “Mega Bundle” advertisement]
Hello! Are you a prepaid customer looking for more data for less money? This message is for you! Get the new \textbf{Mega Bundle} from StarCom Cameroon: 5 GB of data, 200 SMS and 60 minutes of calls to all networks, for only 2,500 FCFA. The bundle is valid for 7 days, so you can browse, chat and call all week long. To subscribe, simply dial *123\# on your phone and follow the instructions. You can pay with your airtime or with mobile money. And that's not all: if you subscribe today, you will receive a bonus of 500 MB, free! The Mega Bundle is perfect for students who need internet for research, for workers who follow the news online, and for families who want to stay in touch. Remember: 5 GB, 200 SMS, 60 minutes, for only 2,500 FCFA, valid for 7 days. Dial *123\# now and enjoy the Mega Bundle from StarCom Cameroon. StarCom — always connected, everywhere in Cameroon.
\end{textebox}

\begin{definition}[Glossaire du Script 1]
\begin{itemize}
\item \eng{prepaid customer} (nom) : client prépayé (qui recharge son crédit)
\item \eng{bundle} (nom) : forfait, bouquet de services
\item \eng{airtime} (nom) : crédit téléphonique
\item \eng{to subscribe} (verbe) : souscrire, s'abonner
\item \eng{bonus} (nom) : avantage supplémentaire offert
\item \eng{to stay in touch} (expression) : rester en contact
\end{itemize}
\end{definition}

\begin{methode}[Traiter une annonce de forfait en 4 questions]
Face à une annonce d'opérateur téléphonique, posez-vous dans l'ordre :
\begin{enumerate}
\item \textbf{Qui parle ?} — Un opérateur (ici : StarCom Cameroon).
\item \textbf{À qui ?} — Aux clients prépayés (\eng{prepaid customers}).
\item \textbf{Quoi ?} — Une offre, définie par 5 informations : le \emph{nom} de l'offre, le \emph{prix}, le \emph{contenu}, la \emph{validité}, le \emph{code} de souscription.
\item \textbf{Notez ces 5 informations dans un tableau pendant l'écoute} — c'est votre fiche de réponses.
\end{enumerate}
\end{methode}

Voici le tableau rempli pendant l'écoute, en appliquant les stratégies 2 et 3 :

\begin{center}
\begin{tabularx}{\linewidth}{|l|X|l|}
\hline
\rowcolor{popblueL} Information & Question type & Réponse notée \\
\hline
Nom de l'offre & What is the name of the bundle? & Mega Bundle \\
\hline
Prix & How much does it cost? & 2,500 FCFA \\
\hline
Contenu & What does it include? & 5 GB + 200 SMS + 60 min \\
\hline
Validité & How long is it valid? & 7 days \\
\hline
Code & How can customers subscribe? & *123\# \\
\hline
Bonus & Is there a special offer? & 500 MB free (today) \\
\hline
\end{tabularx}
\end{center}

Remarquez que toutes les réponses sont \emph{courtes} : un chiffre, un nom, un code. C'est exactement ce qu'attend le correcteur (voir la propriété plus haut).

\begin{exoresolu}[Questions de compréhension sur le Script 1]
Énoncé — Répondez en anglais, de façon courte et précise.
\begin{enumerate}
\item Who is this advertisement for?
\item How much data does the Mega Bundle offer?
\item How long is the bundle valid?
\item What should customers dial to subscribe?
\item What bonus do customers receive if they subscribe today?
\end{enumerate}
\tcblower
Solution détaillée.
\begin{enumerate}
\item \eng{Prepaid customers.} — Stratégie 1 : la deuxième phrase de l'annonce (\eng{Are you a prepaid customer...?}) identifie le public cible.
\item \eng{5 GB (of data).} — Stratégies 2 et 3 : le chiffre est noté en chiffres arabes dès la première écoute.
\item \eng{7 days. / For 7 days.} — Mot-clé \eng{valid} + chiffre ; attention à ne pas confondre avec le bonus.
\item \eng{*123\#.} — Stratégie 4 : si vous l'avez raté à la première écoute, c'est typiquement l'information à guetter à la deuxième écoute (elle est répétée à la fin : \eng{Dial *123\# now}).
\item \eng{A bonus of 500 MB (free).} — Signal : \eng{And that's not all...}, formule classique qui introduit une information supplémentaire souvent interrogée.
\end{enumerate}
\end{exoresolu}

\begin{attention}[Erreurs fréquentes des francophones à l'écoute]
\begin{itemize}
\item \textbf{Tout vouloir comprendre.} Paniquer sur un mot inconnu (\eng{airtime}) fait perdre la phrase suivante. Laissez passer et restez concentré sur les mots-clés.
\item \textbf{Traduire mentalement en français.} La traduction prend du temps ; pensez directement en anglais avec les mots du script.
\item \textbf{Écrire les chiffres en lettres.} \eng{two thousand five hundred} = 8 mots à écrire ; 2500 = 4 caractères. Le choix est vite fait.
\item \textbf{Confondre \eng{minutes} et \eng{SMS}.} Dans un forfait, vérifiez bien quel chiffre correspond à quel service : 200 SMS, 60 minutes — pas l'inverse.
\item \textbf{Oublier la validité.} Les questions sur \eng{How long...?} sont presque systématiques dans les annonces de forfaits.
\end{itemize}
\end{attention}

\begin{savaistu}[Les forfaits au Cameroun : un anglais du quotidien]
Au Cameroun, les annonces de forfaits (\eng{bundles}) des opérateurs de téléphonie mobile sont diffusées en anglais et en français, à la radio, à la télévision et par SMS. Les expressions de cette leçon — \eng{dial *123\#}, \eng{valid for 7 days}, \eng{bonus}, \eng{mobile money} — sont donc de l'anglais \emph{réel}, que vous entendez déjà chaque semaine à Yaoundé, Douala ou Bamenda. Savoir les comprendre à l'écoute, c'est à la fois préparer le Bac et mieux gérer votre propre téléphone.
\end{savaistu}

\courssub{Auto-évaluation : ai-je réussi mon écoute ?}

Après chaque exercice d'écoute, évaluez-vous honnêtement avec la grille suivante. Cochez ce que vous avez réellement obtenu \emph{sans aide}.

\begin{center}
\begin{tabularx}{\linewidth}{|X|c|}
\hline
\rowcolor{popblueL} Critère de réussite & Oui / Non \\
\hline
J'ai trouvé le thème général (une offre de forfait) dès la première écoute & \\
\hline
J'ai identifié le type de texte (annonce publicitaire) & \\
\hline
J'ai noté au moins 3 chiffres en chiffres arabes (5, 200, 60, 2500, 7, 500) & \\
\hline
J'ai relevé au moins 2 noms propres (Mega Bundle, StarCom Cameroon) & \\
\hline
J'ai répondu correctement à toutes les questions de compréhension & \\
\hline
J'ai utilisé la deuxième écoute pour vérifier, pas pour tout refaire & \\
\hline
\end{tabularx}
\end{center}

Si vous avez coché au moins cinq cases sur six, votre méthode d'écoute est efficace. Sinon, identifiez la stratégie qui vous a manqué et refaites l'exercice en vous concentrant uniquement sur elle. La section suivante vous permettra de réemployer tout ce lexique dans une production guidée.

\coursec{Production guidée : réemploi du lexique}

Après avoir analysé les stratégies d'écoute et travaillé la compréhension de textes authentiques sur les forfaits téléphoniques et internet, nous passons maintenant à l'appropriation active du lexique. L'objectif est de transformer la compréhension passive en compétence productive : vous devez être capable d'utiliser le vocabulaire ciblé (\cle{subscribe}, \cle{bundle}, \cle{top up}, \cle{balance}, \cle{validity}) dans des phrases personnelles, de rédiger une annonce publicitaire concise, et de mener une interaction orale réaliste entre un client et un agent commercial. Cette section suit une progression naturelle : de l'écrit guidé vers l'oral interactif, pour finir par une tâche de type « vie réelle » (rédaction d'un SMS publicitaire) souvent rencontrée à l'examen du Baccalauréat.

\courssub{Réemploi écrit : phrases personnelles}

Avant de produire des énoncés complexes, il est indispensable de fixer le sens et la syntaxe des cinq mots-clés de la leçon. Observons d'abord leur comportement grammatical dans des phrases modèles issues du script d'écoute ou créées pour l'occasion.

\begin{textebox}{Modèles d'observation pour le réemploi lexical}
\textbf{Subscribe} (verbe, intransitif + to / transitif direct) : \eng{Customers can subscribe to the new bundle online.} (Les clients peuvent s'abonner au nouveau forfait en ligne.) \\
\textbf{Bundle} (nom comptable) : \eng{The weekly bundle includes 2 GB of data and 100 minutes.} (Le forfait hebdomadaire inclut 2 Go de données et 100 minutes.) \\
\textbf{Top up} (verbe à particule, séparable / nom : \cle{a top-up}) : \eng{I need to top up my credit before the call.} (J'ai besoin de recharger mon crédit avant l'appel.) / \eng{Buy a top-up at the kiosk.} (Achetez une recharge au kiosque.) \\
\textbf{Balance} (nom, solde) : \eng{Check your balance by dialling *123\#.} (Vérifiez votre solde en composant le *123\#.) \\
\textbf{Validity} (nom, durée de validité) : \eng{The validity of this pack is thirty days.} (La validité de ce pack est de trente jours.)
\end{textebox}

\begin{propriete}[Points grammaticaux clés à respecter]
\begin{itemize}
    \item \textbf{Subscribe to} : ne pas oublier la préposition \eng{to} quand l'objet suit (s'abonner \emph{à} quelque chose). À la voix passive : \eng{The service is subscribed to by many users.}
    \item \textbf{Top up} : verbe à particule. Si l'objet est un pronom, il \textbf{doit} se placer entre le verbe et la particule : \eng{top \textbf{it} up} (et non \eng{top up it}).
    \item \textbf{Bundle} / \textbf{Pack} / \textbf{Plan} : synonymes courants pour « forfait ». \eng{Bundle} insiste sur le regroupement de services (data + voix + SMS).
    \item \textbf{Validity} (nom) vs \textbf{Valid} (adjectif) : \eng{The offer is valid for 7 days} (adj.) / \eng{The validity is 7 days} (nom).
    \item \textbf{Balance} : en télécoms, équivaut à \eng{credit} (crédit) ou \eng{account balance} (solde du compte).
\end{itemize}
\end{propriete}

\begin{exemplebox}[Cinq phrases modèles (à adapter pour votre production)]
\begin{enumerate}
    \item \eng{I decided to \textbf{subscribe} to the monthly internet \textbf{bundle} because it is cheaper.} \par \textit{(J'ai décidé de m'abonner au forfait internet mensuel car il est moins cher.)}
    \item \eng{My data \textbf{balance} is low, so I must \textbf{top up} my account today.} \par \textit{(Mon solde de données est bas, je dois donc recharger mon compte aujourd'hui.)}
    \item \eng{The \textbf{validity} of the daily pack expires at midnight.} \par \textit{(La validité du forfait journalier expire à minuit.)}
    \item \eng{If you \textbf{subscribe} before Friday, you get a bonus \textbf{top-up} of 500 MB.} \par \textit{(Si vous vous abonnez avant vendredi, vous obtenez une recharge bonus de 500 Mo.)}
    \item \eng{This \textbf{bundle} offers unlimited calls during its \textbf{validity} period.} \par \textit{(Ce forfait offre des appels illimités pendant sa période de validité.)}
\end{enumerate}
\end{exemplebox}

\begin{exercice}[Votre production : 5 phrases personnelles]
Rédigez \textbf{cinq phrases originales} en anglais, en utilisant \textbf{une fois chacun} des cinq mots obligatoires : \cle{subscribe}, \cle{bundle}, \cle{top up}, \cle{balance}, \cle{validity}. Vos phrases doivent refléter votre réalité (étudiant à Yaoundé, Douala, Buea, etc.) ou une situation imaginée réaliste. Respectez la grammaire (prépositions, ordre des mots, temps verbaux).
\end{exercice}

\begin{corrige}[Exercice : 5 phrases personnelles -- Proposition de correction]
Voici des exemples de phrases correctes et naturelles. Vos phrases peuvent différer ; l'essentiel est l'exactitude grammaticale et lexicale.
\begin{enumerate}
    \item \eng{Many students in Yaoundé \textbf{subscribe} to the "Student Pack" for affordable internet.} (Beaucoup d'étudiants à Yaoundé s'abonnent au "Student Pack" pour un internet abordable.)
    \item \eng{The weekly \textbf{bundle} costs 1000 FCFA and includes 3 GB + 200 SMS.} (Le forfait hebdomadaire coûte 1000 FCFA et inclut 3 Go + 200 SMS.)
    \item \eng{I usually \textbf{top up} my credit via Mobile Money (MoMo) on my phone.} (Je recharge généralement mon crédit via Mobile Money sur mon téléphone.)
    \item \eng{Before calling, check your \textbf{balance} by dialling *144\#.} (Avant d'appeler, vérifiez votre solde en composant le *144\#.)
    \item \eng{The \textbf{validity} of the night bundle is from 10 PM to 5 AM.} (La validité du forfait nuit est de 22h à 5h.)
\end{enumerate}
\end{corrige}

\courssub{Production d'une mini-annonce publicitaire}

Dans la vie professionnelle comme dans les examens (épreuve d'expression écrite), vous pouvez être amené à rédiger un texte court, incitatif et informatif : une \cle{mini-advert} (mini-annonce). Le format impose la concision : nom accrocheur, prix clair, contenu précis, code USSD d'activation, appel à l'action (\cle{call to action}).

\begin{methode}[Rédiger une mini-annonce efficace (5 étapes)]
\begin{enumerate}
    \item \textbf{Accroche (Hook) :} Un nom de forfait court, en majuscules, avec un mot-clé attractif (\eng{Student}, \eng{Night}, \eng{Family}, \eng{Unlimited}).
    \item \textbf{Offre chiffrée (The Deal) :} Donnez \textbf{Data} (Go/Mo), \textbf{Voix} (min), \textbf{SMS} et \textbf{Validity} (jours). Utilisez le symbole \eng{+} ou \eng{\&}.
    \item \textbf{Prix (Price) :} En FCFA (contexte camerounais) ou \$. Format : \eng{for 500 FCFA} / \eng{at 500 FCFA}.
    \item \textbf{Code USSD (Activation Code) :} \eng{Dial *XXX\#} ou \eng{Send "KEYWORD" to XXXX}.
    \item \textbf{Appel à l'action (Call to Action) :} \eng{Subscribe now!}, \eng{Stay connected!}, \eng{Don't miss out!}
\end{enumerate}
\end{methode}

\begin{textebox}{Modèle de mini-annonce (Modèle de référence)}
\textbf{ORANGE STUDENT PACK!} \\
1 GB Data + 50 SMS + 20 Min for \textbf{500 FCFA}. \\
Validity: \textbf{3 Days}. \\
Dial \textbf{*111\#} now and \textbf{stay connected!}
\end{textebox}

\begin{exoresolu}[Analyse du modèle]
\textbf{Énoncé :} Identifiez les 5 composantes de la méthode dans l'annonce ci-dessus.
\tcblower
\textbf{Solution détaillée :}
\begin{enumerate}
    \item \textbf{Hook :} "ORANGE STUDENT PACK!" (Majuscules, cible "Student").
    \item \textbf{Deal :} "1 GB Data + 50 SMS + 20 Min" (Data, SMS, Voix quantifiés).
    \item \textbf{Price :} "for 500 FCFA" (Préposition \eng{for} pour le prix).
    \item \textbf{Code :} "Dial *111\#" (Impératif + code USSD).
    \item \textbf{CTA :} "now and stay connected!" (Urgence + bénéfice).
\end{enumerate}
\textbf{Note lexicale :} \eng{Data} est indénombrable (pas de 's'), mais on écrit \eng{1 GB of data} ou \eng{1 GB data} (nom attribut). \eng{Min} = abréviation acceptée pour \eng{minutes} dans les annonces. \eng{Validity} est le nom ; l'adjectif serait \eng{Valid for 3 days}.
\end{exoresolu}

\begin{exercice}[Votre mini-annonce]
Imaginez un forfait imaginaire pour un opérateur fictif (\eng{GreenNet}, \cle{BlueCom}, \cle{YellowTel}...). Rédigez votre annonce (4 à 5 lignes max) en respectant la structure en 5 étapes. Inventez un nom, un prix en FCFA, un contenu (Data, Voix, SMS), une validité (jours/heures) et un code USSD (*...\#).
\end{exercice}

\begin{corrige}[Exercice : Mini-annonce -- Proposition]
\textbf{GREENNET NIGHT OWL PACK!} \\
5 GB Night Data + Unlimited On-Net Calls for \textbf{300 FCFA}. \\
Validity: \textbf{10 PM -- 6 AM} (Daily). \\
Dial \textbf{*222\#} to activate. \textbf{Surf all night!}
\end{corrige}

\begin{popfigure}
\begin{tikzpicture}[
    node distance=0.8cm and 1.2cm,
    box/.style={draw=popblue, thick, rounded corners, fill=popblueL, minimum width=3.5cm, minimum height=1.2cm, align=center, font=\small},
    arrow/.style={-Stealth, thick, popdark},
    labelbox/.style={font=\tiny\itshape, text=popmuted, anchor=north}
]
% Nœuds principaux
\node[box, align=center] (hook){\textbf{HOOK}\\\eng{NAME OF PACK}};
\node[box, right=of hook, align=center] (deal){\textbf{DEAL}\\\eng{Data + Voice + SMS}};
\node[box, right=of deal, align=center] (price){\textbf{PRICE}\\\eng{XXX FCFA}};
\node[box, right=of price, align=center] (code){\textbf{CODE}\\\eng{Dial *...\#}};
\node[box, right=of code, align=center] (cta){\textbf{CTA}\\\eng{Action Verb!}};

% Flèches
\draw[arrow] (hook) -- (deal);
\draw[arrow] (deal) -- (price);
\draw[arrow] (price) -- (code);
\draw[arrow] (code) -- (cta);

% Étiquettes sous les boîtes
\node[labelbox] at (hook.south) {Ex: MTN STUDENT PACK};
\node[labelbox] at (deal.south) {Ex: 2GB + 100min + 50SMS};
\node[labelbox] at (price.south) {Ex: 500 FCFA};
\node[labelbox] at (code.south) {Ex: *150\#};
\node[labelbox] at (cta.south) {Ex: Subscribe Now!};
\end{tikzpicture}
\legende{Structure visuelle d'une mini-annonce publicitaire télécom (5 blocs obligatoires).}
\end{popfigure}

\courssub{Dialogue guidé en binômes : Client / Agent commercial}

L'interaction orale (\cle{speaking}) est une compétence évaluée au Baccalauréat. Le scénario type est une \cle{service encounter} (interaction de service) : un client appelle le service client (\cle{customer care}) pour s'informer ou souscrire. La politesse (\cle{politeness strategies}) et la précision des informations chiffrées (\cle{figures: price, data, validity}) sont cruciales.

\begin{methode}[Jeu de rôle : Client vs Agent Commercial -- Règles du jeu]
\begin{enumerate}
    \item \textbf{Préparation (2 min) :} En binôme, désignez le \textbf{Client (A)} et l'\textbf{Agent (B)}. L'Agent a une "fiche offre" (inventée ou issue de l'exercice précédent). Le Client a une "fiche besoin" (budget, data souhaitée, durée).
    \item \textbf{Déroulement (3-4 min) :} Suivez la trame ci-dessous. \textbf{Le Client doit poser au moins 3 questions} (ex: \eng{Price?}, \eng{Data volume?}, \eng{Validity?}, \eng{How to activate?}). \textbf{L'Agent doit donner au moins 3 informations chiffrées} (prix, volume data, durée validité, code USSD).
    \item \textbf{Politesse obligatoire :} Le Client utilise \eng{I'd like to...} / \eng{Could you tell me...?} / \eng{Thank you.}. L'Agent utilise \eng{How may I help you?} / \eng{Certainly.} / \eng{Thank you for calling [Operator Name].}
    \item \textbf{Inversion des rôles :} Rejouez la scène en échangeant les rôles et en changeant de forfait.
\end{enumerate}
\end{methode}

\begin{popfigure}
\begin{tikzpicture}[
    node distance=0.6cm and 0.5cm,
    bubble/.style={draw=poporange, thick, rounded corners, fill=poporangeL, minimum width=6cm, minimum height=1.5cm, align=left, font=\small, text width=5.5cm},
    role/.style={font=\bfseries\small, text=popdark, anchor=east},
    arrow/.style={-Stealth, poporange, thick}
]
% Agent 1
\node[role] (a1r) at (0,0) {\eng{Agent:}};
\node[bubble, right=of a1r] (a1) {Good morning, \textbf{GreenNet Customer Care}, \textit{Amadou speaking}. \textbf{How may I help you?}};

% Client 1
\node[role] (c1r) at (0,-2.5) {\eng{Customer:}};
\node[bubble, right=of c1r] (c1) {Good morning. \textbf{I'd like to subscribe to} a data bundle. \textbf{Could you tell me} what offers you have?};

% Agent 2
\node[role] (a2r) at (0,-5) {\eng{Agent:}};
\node[bubble, right=of a2r] (a2) {Certainly. We have the \textbf{Weekly Pack}: \textbf{2 GB} + 100 SMS for \textbf{1000 FCFA}. \textbf{Valid for 7 days}. Or the \textbf{Monthly Pack}: 10 GB for 3000 FCFA.};

% Client 2
\node[role] (c2r) at (0,-7.5) {\eng{Customer:}};
\node[bubble, right=of c2r] (c2) {\textbf{How much is} the Weekly Pack? \textbf{And how do I activate it?}};

% Agent 3
\node[role] (a3r) at (0,-10) {\eng{Agent:}};
\node[bubble, right=of a3r] (a3) {It costs \textbf{1000 FCFA}. Dial \textbf{*150\#} and select option 1. \textbf{Your balance} will be updated instantly.};

% Client 3
\node[role] (c3r) at (0,-12.5) {\eng{Customer:}};
\node[bubble, right=of c3r] (c3) {Perfect. \textbf{Thank you for your help.} I will top up now. Goodbye.};

% Agent 4
\node[role] (a4r) at (0,-15) {\eng{Agent:}};
\node[bubble, right=of a4r] (a4) {You're welcome. \textbf{Thank you for calling GreenNet.} Have a nice day!};

% Flèches de flux
\draw[arrow] (a1) -- (c1);
\draw[arrow] (c1) -- (a2);
\draw[arrow] (a2) -- (c2);
\draw[arrow] (c2) -- (a3);
\draw[arrow] (a3) -- (c3);
\draw[arrow] (c3) -- (a4);
\end{tikzpicture}
\legende{Trame de dialogue type (Client / Agent) avec exemples de phrases cibles. Les parties en \textbf{gras} sont les éléments lexicaux/grammaticaux obligatoires à produire.}
\end{popfigure}

\begin{exemplebox}[Expressions utiles pour le dialogue (Boîte à outils orale)]
\begin{center}
\begin{tabularx}{\linewidth}{|X|X|}\hline
\rowcolor{popblueL} \textbf{Client (Vous)} & \textbf{Agent (Partenaire)} \\ \hline
\eng{I'd like to subscribe to...} & \eng{How may I help you?} \\
\eng{Could you tell me the price of...?} & \eng{We have several bundles...} \\
\eng{What is the validity period?} & \eng{It is valid for [X] days.} \\
\eng{How much data does it include?} & \eng{It includes [X] GB of data.} \\
\eng{How do I activate / top up?} & \eng{Dial *...\# / Send SMS to...} \\
\eng{Can I check my balance?} & \eng{Dial *...\# to check your balance.} \\
\eng{Thank you for your help.} & \eng{Thank you for calling [Name].} \\ \hline
\end{tabularx}
\end{center}
\end{exemplebox}

\begin{attention}[Pièges à l'oral : Politesse et "Faux Amis" structurels]
\begin{itemize}
    \item \textbf{NE DITES PAS :} \eng{"I want a bundle."} (Trop direct, impoli).
    \item \textbf{DITES :} \eng{"\textbf{I'd like} a bundle."} ou \eng{"\textbf{Could I have} a bundle?"} (\cle{Conditionnel de politesse} / \cle{Modaux}).
    \item \textbf{NE DITES PAS :} \eng{"What is the price?"} (Un peu brutal).
    \item \textbf{DITES :} \eng{"\textbf{How much is} the bundle?"} ou \eng{"\textbf{Could you tell me} the price?"} (\cle{Discours indirect} pour la politesse).
    \item \textbf{NE DITES PAS :} \eng{"Recharge my phone."} (Ordre).
    \item \textbf{DITES :} \eng{"\textbf{I'd like to top up} my account, please."}
    \item \textbf{Faux ami :} \eng{Validity} (nom) $\neq$ \eng{Validité} (orthographe fr). Adjectif = \eng{Valid}. \eng{Expiry date} = Date d'expiration.
    \item \textbf{Prononciation :} \eng{Bundle} $\rightarrow$ « \textbf{ban-dl} » (u = \eng{/$\Lambda$/}, pas « boun-dle »). \eng{Validity} $\rightarrow$ « \textbf{va-LI-di-ti} » (accent sur 3e syllabe). \eng{Subscribe} $\rightarrow$ « \textbf{sab-SKRAIB} » (u = \eng{/$\Lambda$/}).
\end{itemize}
\end{attention}

\begin{experience}[Activité de classe : "Le Marché des Forfaits" (15 min)]
\begin{enumerate}
    \item \textbf{Phase 1 (Prépa) :} Chaque binôme crée \textbf{2 fiches "Offre Agent"} différentes (Nom, Prix, Data, SMS, Validité, Code USSD) sur un papier.
    \item \textbf{Phase 2 (Mise en relation) :} Les binômes se mélangent. Un élève devient "Client itinérant", l'autre reste "Agent fixe" à sa table.
    \item \textbf{Phase 3 (Négociation) :} Le Client visite 3 "Agents" différents en 5 min. Il doit choisir la meilleure offre pour son besoin (ex: "Je veux 5 Go pour 7 jours, budget 1500 FCFA"). Il note les infos clés.
    \item \textbf{Phase 4 (Restitution) :} Le Client revient vers son partenaire initial et justifie son choix en anglais : \eng{"I chose the BlueCom pack because it offers 5 GB for 7 days at 1200 FCFA."}
\end{enumerate}
\end{experience}

\courssub{Restitution orale : Présentation devant la classe}

Deux binômes volontaires (ou tirés au sort) jouent leur dialogue devant la classe. C'est l'évaluation formative de la compétence \cle{Speaking : Interaction}.

\begin{propriete}[Grille d'évaluation orale (Critères Bac / CECRL B1-B2)]
\begin{center}
\begin{tabularx}{\linewidth}{|l|X|X|}\hline
\rowcolor{popblueL} \textbf{Critère} & \textbf{Attendu (Réussi)} & \textbf{Points de vigilance} \\ \hline
\textbf{Interaction} & Prise de parole fluide, tours de parole respectés, réaction aux infos du partenaire. & Silences longs, lecture du texte, pas d'écoute active. \\ \hline
\textbf{Lexique ciblé} & Utilisation correcte et naturelle de : \eng{subscribe, bundle, top up, balance, validity, data, activate, dial}. & Oubli des mots obligatoires, confusion \eng{valid/validity}, \eng{top up/top-up}. \\ \hline
\textbf{Grammaire / Précision} & Questions correctes (\eng{How much...?}, \eng{What is...?}), Politesse (\eng{I'd like, Could you...}), Chiffres bien dits. & \eng{How much cost?} (erreur), \eng{I want}, mauvaises prépositions (\eng{subscribe \textbf{for}}). \\ \hline
\textbf{Prononciation / Intonation} & Mots-clés articulés (\eng{bundle, validity, subscribe}), intonation montante questions, descendante affirmations. & Accent français trop marqué, chiffres inaudibles (ex: \eng{thirteen} vs \eng{thirty}). \\ \hline
\textbf{Tâche accomplie} & Le client obtient les 3 infos chiffrées ; l'agent donne le code USSD. & Infos manquantes, dialogue incomplet. \\ \hline
\end{tabularx}
\end{center}
\end{propriete}

\courssub{Complément Bac : Rédiger un SMS publicitaire (160 caractères)}

Format court, très fréquent dans la vie réelle (marketing direct) et possible à l'écrit du Baccalauréat (production de message court). Contraintes : \textbf{160 caractères max} (espaces compris), langage \cle{SMS} / \cle{Textspeak} toléré mais anglais correct exigé, informations essentielles seulement.

\begin{methode}[Rédiger un SMS Pub -- 4 Règles d'Or]
\begin{enumerate}
    \item \textbf{Identité :} Nom opérateur/forfait au début (\eng{MTN:}, \eng{Orange:}).
    \item \textbf{Offre choc :} Chiffres clés seulement (\eng{5GB}, \eng{1000F}, \eng{7d}).
    \item \textbf{Action immédiate :} Code court (\eng{Dial *123\#}, \eng{Reply YES}).
    \item \textbf{Opt-out légal :} \eng{Stop? Dial *123*0\#} (souvent obligatoire).
\end{enumerate}
\end{methode}

\begin{exemplebox}[Exemples de SMS Pub (comptage caractères)]
\begin{enumerate}
    \item \eng{Orange: Student Pack! 2GB+100SMS 7days=1000F. Dial *111\#. Stop? *111*0\#} \par \textit{(96 caractères -- Parfait. Clair, complet, code USSD, opt-out.)}
    \item \eng{MTN Deal: Night 5GB 10pm-5am only 300F! Valid 1 day. Dial *150\# to sub.} \par \textit{(98 caractères -- Efficace. "Sub" abbr. de subscribe accepté en SMS.)}
    \item \eng{BlueCom: Unlimited Social Media (FB, WhatsApp) 30days 2000F. *200\#. T\&Cs apply.} \par \textit{(102 caractères -- Bon. "T\&Cs" = Terms and Conditions.)}
\end{enumerate}
\end{exemplebox}

\begin{attention}[Erreurs fréquentes en SMS Bac]
\begin{itemize}
    \item \textbf{Trop long :} Dépasser 160 caractères $\rightarrow$ message coupé en 2 (coût double) ou refusé. \textbf{Comptez vos caractères !}
    \item \textbf{Trop vague :} \eng{"Cheap data, call now."} (Pas de prix, pas de volume, pas de code $\rightarrow$ 0 point au contenu).
    \item \textbf{Abrévisions dangereuses :} \eng{"2G"} (ambigu : 2 Go ou réseau 2G ?) $\rightarrow$ Préférez \eng{2GB}. \eng{"1k"} pour 1000 FCFA ? Risqué. Écrivez \eng{1000F} ou \eng{1000FCFA}.
    \item \textbf{Oubli du code USSD :} L'appel à l'action est la finalité commerciale. Sans \eng{Dial *...\#}, l'annonce est inutile.
\end{itemize}
\end{attention}

\begin{exercice}[Bac Blanc : SMS Publicitaire]
Rédigez un SMS publicitaire (max 160 caractères) pour le forfait imaginaire que vous avez créé à l'exercice 2 (Mini-annonce). Respectez les 4 Règles d'Or. Comptez vos caractères (lettres + espaces + symboles) et notez le total à la fin.
\end{exercice}

\begin{corrige}[Exercice Bac : SMS -- Proposition pour "GreenNet Night Owl Pack"]
\eng{GreenNet: NightOwl 5GB+UnlimCalls 10pm-6am 300F. Dial *222\#. Stop *222*0\#} \par \textit{(104 caractères. Conforme : Identité, Offre chiffrée, Code, Opt-out.)}
\end{corrige}

\courssub{Synthèse lexicale et révision finale}

Pour ancrer durablement le vocabulaire de cette leçon, complétez le tableau récapitulatif ci-dessous. Il oppose la forme anglaise, sa nature grammaticale, sa traduction française et un exemple contextuel court. C'est votre "fiche de révision" pour le Bac.

\begin{center}
\begin{tabularx}{\linewidth}{|l|l|l|X|}\hline
\rowcolor{popblueL} \textbf{English Word} & \textbf{Part of Speech} & \textbf{Français} & \textbf{Exemple contextuel (English / Français)} \\ \hline
\cle{Subscribe} & Verb (v.) & S'abonner, souscrire & \eng{Subscribe to the pack.} / S'abonner au forfait. \\ \hline
\cle{Subscription} & Noun (n.) & L'abonnement, la souscription & \eng{Your subscription is active.} / Votre abonnement est actif. \\ \hline
\cle{Bundle} / \cle{Pack} / \cle{Plan} & Noun (n.) & Le forfait, le pack, le plan & \eng{A weekly data bundle.} / Un forfait data hebdo. \\ \hline
\cle{Top up} (v.) / \cle{Top-up} (n.) & Verb / Noun & Recharger / Une recharge & \eng{Top up your credit.} / \eng{Buy a top-up.} \\ \hline
\cle{Balance} & Noun (n.) & Le solde, le crédit restant & \eng{Check your balance.} / Vérifiez votre solde. \\ \hline
\cle{Validity} & Noun (n.) & La validité, la durée & \eng{Validity: 30 days.} / Validité : 30 jours. \\ \hline
\cle{Valid} & Adjective (adj.) & Valide & \eng{The code is valid.} / Le code est valide. \\ \hline
\cle{Activate} & Verb (v.) & Activer & \eng{Activate the bundle.} / Activer le forfait. \\ \hline
\cle{Dial} & Verb (v.) & Composer (un numéro/USSD) & \eng{Dial *123\#.} / Composez le *123\#. \\ \hline
\cle{Data} / \cle{Internet} & Noun (n., U) & Les données, internet mobile & \eng{High-speed data.} / Données haut débit. \\ \hline
\cle{On-net} / \cle{Off-net} & Adj. / Adv. & Vers même réseau / Autre réseau & \eng{Unlimited on-net calls.} / Appels illimités même réseau. \\ \hline
\cle{Customer Care} / \cle{Support} & Noun (n.) & Service client & \eng{Call customer care.} / Appelez le service client. \\ \hline
\end{tabularx}
\end{center}

\begin{savaistu}[Le saviez-vous ? Le langage SMS au Cameroun et dans le monde]
Au Cameroun, le \cle{Mobile Money} (MoMo, Orange Money) a révolutionné le \eng{top-up}. On dit souvent \eng{"I'll MoMo you"} (je t'envoie de l'argent par MoMo) ou \eng{"Buy credit via MoMo"}. Dans les pays anglophones (Nigeria, Ghana, UK, USA), les codes courts (\cle{Short Codes}) comme \eng{*123\#} sont appelés \cle{USSD codes} (Unstructured Supplementary Service Data). Le terme \eng{bundle} est universel en Afrique anglophone (Nigeria, Ghana, Kenya, Afrique du Sud), tandis qu'en France on dit "forfait" et au Québec "plan". Le mot \eng{validity} est très utilisé en Afrique de l'Ouest dans les télécoms (souvent raccourci en \eng{validity period}), alors qu'en Europe on verra plus souvent \eng{valid for X days} ou \eng{expiry date}.
\end{savaistu}

\begin{aretenir}[Check-list finale de la Leçon 49 -- Êtes-vous prêt pour le Bac ?]
\begin{itemize}
    \item[$\square$] Je sais écrire 5 phrases correctes avec \eng{subscribe, bundle, top up, balance, validity}.
    \item[$\square$] Je sais structurer une mini-annonce (Hook, Deal, Price, Code, CTA).
    \item[$\square$] Je sais jouer le dialogue Client/Agent : politesse (\eng{I'd like, Could you}), 3 questions client, 3 infos chiffrées agent.
    \item[$\square$] Je prononce correctement : \eng{bundle} (« ban-dl »), \eng{validity} (« va-LI-di-ti »), \eng{subscribe} (« sab-SKRAIB »).
    \item[$\square$] Je sais rédiger un SMS pub < 160 caractères avec les infos essentielles et le code USSD.
    \item[$\square$] Je distingue \eng{validity} (nom) et \eng{valid} (adj.), et je place le pronom après \eng{top} dans \eng{top it up}.
\end{itemize}
\end{aretenir}

\newpage

\coursec{Activité d'intégration}

\coursec{Lesson 49 — Listening: Texts about subscribing to service packages (telephone/internet services)}

\courssub{Objectifs de l'activité}
Cette activité d'intégration vise à mobiliser l'ensemble des compétences travaillées dans la leçon : compréhension de l'oral (\emph{listening}), prise de notes structurée, réemploi du lexique spécialisé des télécommunications (forfaits, data, validité, codes USSD, recharges) et expression orale en interaction (\emph{speaking}). Concrètement, l'élève doit être capable de :
\begin{itemize}
    \item Écouter activement une annonce promotionnelle (radio, haut-parleur en agence, message vocal) et en extraire les informations chiffrées et techniques essentielles.
    \item Organiser ces informations dans un tableau de synthèse (\textsc{Offer / Price / Data / Validity / USSD Code}).
    \item Engager une conversation téléphonique ou en face-à-face avec un agent de service client (\emph{customer care agent}) pour vérifier, corriger et confirmer ces informations avant de souscrire.
    \item Utiliser les formes de politesse, les structures de vérification (\eng{Could you confirm...?}, \eng{Did you say...?}) et le vocabulaire technique de manière fluide et naturelle.
\end{itemize}

\courssub{Situation}
\textbf{Contexte :} Vous êtes à Yaoundé, au quartier \emph{Mvog-Ada}, devant une agence \textbf{MTN Cameroon} (ou \emph{Orange Cameroon} / \emph{Camtel}). Un haut-parleur diffuse en boucle une annonce promotionnelle pour un nouveau forfait « \emph{Back-to-School} » destiné aux élèves et étudiants. Vous n'avez pas de smartphone pour scanner le QR code affiché, vous devez donc vous fier à votre écoute. Vous décidez d'entrer dans l'agence pour souscrire auprès d'un agent au guichet \emph{Customer Care}.

\textbf{Document 1 : Annonce radio / haut-parleur (Annonce n°1 — Pour l'élève A)}
\begin{textebox}[Script d'écoute — Annonce MTN « Back-to-School Special »]
\eng{``Hello everyone! MTN Cameroon brings you the ultimate Back-to-School bundle! Introducing the \textbf{MTN Scholar Pack}. For only \textbf{2,500 FCFA}, you get a massive \textbf{10 GB} of high-speed internet data valid for \textbf{30 days}. That's enough for your research, online classes and WhatsApp videos. But wait, there's more! You also enjoy \textbf{unlimited MTN-to-MTN calls} for the whole month. To subscribe, simply dial \textbf{*157*99\#} and select option 1. Don't miss out! Offer ends October 31st. Terms and conditions apply. MTN, Everywhere you go.''}
\end{textebox}

\textbf{Document 2 : Fiche technique de l'agent (Pour l'élève B — Vérification)}
\begin{textebox}[Agent's Reference Sheet — MTN Scholar Pack (Real Data)]
\begin{center}
\begin{tabularx}{\linewidth}{|X|X|}\hline
\rowcolor{popblueL} \textbf{Parameter} & \textbf{Details} \\ \hline
Commercial Name & MTN Scholar Pack \\ \hline
Price & 2,500 FCFA (Taxes included) \\ \hline
Internet Data & 10 GB (4G/3G) \\ \hline
Validity & 30 Calendar Days \\ \hline
Voice Bonus & Unlimited On-Net (MTN-to-MTN) calls \\ \hline
Off-Net Calls & Standard rates apply (50 FCFA/min) \\ \hline
Subscription Code & *157*99\# then Option 1 \\ \hline
Eligibility & Prepaid customers only (Student ID required for first sub) \\ \hline
Promo End Date & October 31st \\ \hline
\end{tabularx}
\end{center}
\end{textebox}

\textbf{Document 3 : Annonce radio / haut-parleur (Annonce n°2 — Pour l'élève B, rôles inversés)}
\begin{textebox}[Script d'écoute — Annonce Orange « Campus Connect »]
\eng{``Students of Cameroon, Orange is connecting your future! Get the \textbf{Orange Campus Connect} offer. Pay \textbf{3,000 FCFA} and receive \textbf{15 GB} of 4G data, valid for \textbf{45 days}. Perfect for the full semester! Plus, get \textbf{1 hour of free calls to all networks} (Orange, MTN, Camtel, Fixes). To activate, dial \textbf{*141*50\#}. No ID needed, just dial and learn. Valid until December 31st. Orange, Listen to the future.''}
\end{textebox}

\textbf{Document 4 : Fiche technique de l'agent (Pour l'élève A — Vérification)}
\begin{textebox}[Agent's Reference Sheet — Orange Campus Connect (Real Data)]
\begin{center}
\begin{tabularx}{\linewidth}{|X|X|}\hline
\rowcolor{poporangeL} \textbf{Parameter} & \textbf{Details} \\ \hline
Commercial Name & Orange Campus Connect \\ \hline
Price & 3,000 FCFA \\ \hline
Internet Data & 15 GB (4G+) \\ \hline
Validity & 45 Calendar Days \\ \hline
Voice Bonus & 60 Minutes All Networks (On-net + Off-net) \\ \hline
Subscription Code & *141*50\# \\ \hline
Eligibility & All Prepaid Customers \\ \hline
Promo End Date & December 31st \\ \hline
\end{tabularx}
\end{center}
\end{textebox}

\courssub{Déroulement de l'activité}

\begin{methode}[Stratégies d'écoute active — Repérage de l'information clé]
\textbf{Avant l'écoute (Anticipation) :}
\begin{enumerate}
    \item Lire le tableau de prise de notes (catégories : \eng{Offer Name, Price, Data, Validity, Voice Bonus, USSD Code, Deadline}).
    \item Prédire le type d'information attendue : un nom propre (nom de l'offre), des chiffres (prix, volume, jours), un code (symboles * \#), une date.
    \item Activer le lexique thématique (\cle{bundle}, \cle{data allowance}, \cle{validity}, \cle{on-net/off-net}, \cle{subscribe}, \cle{dial}).
\end{enumerate}
\textbf{Pendant l'écoute (Écoute sélective) :}
\begin{enumerate}
    \item \textbf{Première écoute (globale) :} Identifier le thème (promo étudiante), l'opérateur (MTN / Orange) et l'intention (souscrire).
    \item \textbf{Deuxième écoute (détail) :} Remplir le tableau case par case. Ne pas écrire de phrases complètes : mots-clés et chiffres seulement.
    \item \textbf{Troisième écoute (vérification) :} Contrôler les chiffres (10 vs 15 GB, 30 vs 45 days), le code USSD (*157*99\# vs *141*50\#) et la date limite.
\end{enumerate}
\textbf{Après l'écoute (Transfert) :} Réutiliser les notes pour interagir avec un interlocuteur (jeu de rôle).
\end{methode}

\courssub{Phase 1 : Écoute active et prise de notes (Individuel — 10 min)}
\begin{enumerate}
    \item \textbf{Élève A :} Écoute l'\textbf{Annonce n°1 (MTN Scholar Pack)}. Remplis le tableau ci-dessous \textbf{en anglais} pendant l'écoute.
    \item \textbf{Élève B :} Écoute l'\textbf{Annonce n°2 (Orange Campus Connect)}. Remplis le même tableau.
    \item \textbf{Tableau de prise de notes (à recopier sur votre cahier) :}
\end{enumerate}
\begin{center}
\begin{tabularx}{\linewidth}{|X|X|}\hline
\rowcolor{popblueL} \textbf{Category} & \textbf{Your Notes} \\ \hline
Offer Name & \\ \hline
Price (FCFA) & \\ \hline
Data Volume & \\ \hline
Validity Period & \\ \hline
Voice Bonus (Calls) & \\ \hline
USSD Code / Shortcode & \\ \hline
Special Condition / Deadline & \\ \hline
\end{tabularx}
\end{center}

\begin{experience}[Jeu de rôle — Au guichet Customer Care (En binôme — 15 min par rôle)]
\textbf{Scénario 1 :} L'élève A est le \textbf{Client}. L'élève B est l'\textbf{Agent MTN} (utilise Doc 2).
\begin{itemize}
    \item Le Client salue, dit vouloir souscrire au \eng{Scholar Pack} et lit ses notes pour confirmer.
    \item L'Agent écoute, compare avec sa \textbf{Fiche technique}, valide ou \textbf{corrige} poliment les erreurs/omissions (ex: \eng{Actually, sir, the off-net calls are not free...}, \eng{Please note that a Student ID is required...}).
    \item L'Agent guide le client pour la souscription (composition du code, confirmation par SMS).
    \item Le Client remercie et termine l'appel.
\end{itemize}
\textbf{Scénario 2 (Inversion des rôles) :} L'élève B est le \textbf{Client} (pour l'offre Orange, Doc 3). L'élève A est l'\textbf{Agent Orange} (utilise Doc 4).
\textbf{Contrainte de langue :} Le dialogue se déroule \textbf{uniquement en anglais}. Utilisez au moins 5 mots du lexique de la leçon (\cle{bundle}, \cle{subscribe}, \cle{validity}, \cle{dial}, \cle{on-net/off-net}, \cle{recharge}, \cle{promo}, \cle{terms and conditions}).
\end{experience}

\courssub{Phase 3 : Évaluation par les pairs et synthèse (Collectif — 10 min)}
\begin{enumerate}
    \item Deux binômes volontaires jouent leur dialogue devant la classe.
    \item La classe note sur une fiche d'évaluation : \eng{Clarity of information}, \eng{Politeness}, \eng{Correct vocabulary}, \eng{Accuracy of corrections}.
    \item Vote pour le \eng{Best Customer Care Interaction}.
\end{enumerate}

\courssub{Ressources à mobiliser}

\begin{definition}[Lexique clé — Telecommunications Packages]
\begin{itemize}
    \item \textbf{Noms :} \eng{a bundle / a pack / a plan} (un forfait) ; \eng{data allowance / data volume} (volume de data) ; \eng{validity / validity period} (durée de validité) ; \eng{a subscription / a sub} (un abonnement) ; \eng{a recharge / a top-up} (une recharge) ; \eng{on-net calls} (communications vers le même réseau) ; \eng{off-net calls} (communications vers les autres réseaux/fixes) ; \eng{a shortcode / a USSD code} (code USSD, ex: *157\#) ; \eng{a promo / a special offer} (une promo) ; \eng{terms and conditions (T\&Cs)} (conditions générales).
    \item \textbf{Verbes :} \eng{to subscribe (to)} (s'abonner à) ; \eng{to dial} (composer un numéro/code) ; \eng{to activate / to trigger} (activer) ; \eng{to confirm} (confirmer) ; \eng{to expire} (expirer) ; \eng{to roll over} (reporter le solde non utilisé).
    \item \textbf{Adjectifs :} \eng{unlimited / uncapped} (illimité) ; \eng{capped} (plafonné) ; \eng{eligible} (éligible) ; \eng{prepaid / postpaid} (prépayé / postpayé).
\end{itemize}
\end{definition}

\begin{propriete}[Structures utiles pour la vérification et la correction (Checking \& Correcting)]
\begin{itemize}
    \item \textbf{Demander confirmation :} \eng{Could you confirm the price, please?} / \eng{Did you say 10 GB or 15 GB?} / \eng{Is that 30 days or 30 calendar days?}
    \item \textbf{Corriger poliment (Agent $\to$ Client) :} \eng{Actually, the offer includes...} / \eng{I'm afraid that's not quite right.} / \eng{Please note that off-net calls are charged at standard rates.} / \eng{The correct code is *157*99\#, not *157*98\#.}
    \item \textbf{Corriger ses propres notes (Client) :} \eng{Oh, I see. Let me correct that.} / \eng{Right, 45 days, not 30. Thanks for the clarification.}
    \item \textbf{Guider l'action (USSD) :} \eng{Please dial *157*99\# on your phone.} / \eng{Select option 1.} / \eng{You will receive a confirmation SMS shortly.}
\end{itemize}
\end{propriete}

\begin{attention}[Erreurs fréquentes des francophones — Pièges à éviter]
\begin{itemize}
    \item \textbf{Faux ami :} \eng{``Abonnement''} $\neq$ \eng{Abomination}. Dire \eng{\textbf{Subscription}} ou \eng{\textbf{Plan}}. \eng{Subscription} = l'acte de s'abonner ; \eng{Plan/Package} = le forfait lui-même.
    \item \textbf{Préposition :} \eng{Subscribe \textbf{to} a plan} (pas \eng{subscribe a plan}).
    \item \textbf{Data :} \eng{Data} est \textbf{indénombrable} en anglais technique. Dire \eng{10 GB \textbf{of} data} (pas \eng{10 datas}).
    \item \textbf{Validity :} Nom = \eng{Validity} (la durée). Adjectif = \eng{Valid} (valide). Dire \eng{The validity is 30 days} (correct) ou \eng{The offer is valid \textbf{for} 30 days}.
    \item \textbf{Prononciation :} \eng{USSD} se prononce lettre par lettre « \emph{you-ess-es-di} ». \eng{Giga} dans \eng{Gigabyte} se prononce « \emph{djaï-ga} » (son \eng{dj}), pas « \emph{gi-ga} » à la française.
    \item \textbf{Chiffres :} \eng{2,500} s'écrit \eng{2,500} (virgule anglaise = \emph{comma}) mais se lit \eng{two thousand five hundred} (pas \eng{two thousand five hundred \textbf{FCFA}} si le contexte est clair, sinon \eng{two thousand five hundred \textbf{CFA francs}}).
\end{itemize}
\end{attention}

\courssub{Production guidée et corrigé commenté}

\begin{exoresolu}[Modèle de dialogue — Scénario 1 : MTN Scholar Pack (Élève A = Client / Élève B = Agent)]
\textbf{Énoncé :} Reproduire l'interaction complète : prise de notes $\to$ vérification $\to$ correction $\to$ souscription.

\tcblower

\textbf{Partie 1 : Tableau de prise de notes complété (Élève A — Résultat attendu après écoute)}
\begin{center}
\begin{tabularx}{\linewidth}{|X|X|}\hline
\rowcolor{popblueL} \textbf{Category} & \textbf{Notes (Student A)} \\ \hline
Offer Name & MTN Scholar Pack \\ \hline
Price (FCFA) & 2,500 FCFA \\ \hline
Data Volume & 10 GB \\ \hline
Validity Period & 30 days \\ \hline
Voice Bonus (Calls) & Unlimited MTN-to-MTN calls \\ \hline
USSD Code / Shortcode & *157*99\# (Option 1) \\ \hline
Special Condition / Deadline & Ends Oct 31st \\ \hline
\end{tabularx}
\end{center}

\textbf{Partie 2 : Dialogue modèle (Annoté)}

\begin{tabularx}{\linewidth}{|l|X|}\hline
\rowcolor{popblueL} \textbf{Rôle} & \textbf{Dialogue (English)} \\ \hline
\textbf{Agent} & \eng{Good morning! Welcome to MTN Customer Care. How can I help you today?} \\ \hline
\textbf{Client} & \eng{Good morning. I heard the announcement for the \textbf{Scholar Pack} outside. I would like to \textbf{subscribe to} this \textbf{bundle}.} \\ \hline
\textbf{Agent} & \eng{Certainly. It is an excellent \textbf{promo} for students. Let me \textbf{confirm} the details with you first.} \\ \hline
\textbf{Client} & \eng{Yes, please. From what I noted: the \textbf{price} is 2,500 FCFA. I get \textbf{10 GB of data} with a \textbf{validity} of 30 days. And there are \textbf{unlimited on-net calls}.} \\ \hline
\textbf{Agent} & \eng{That is \textbf{correct} for the data and the on-net calls. \textbf{However}, please \textbf{note} that \textbf{off-net calls} (to Orange or Camtel) are \textbf{not included}; they are charged at the standard rate of 50 FCFA per minute.} \\ \hline
\textbf{Client} & \eng{Ah, I see. \textbf{Did the announcement mention that?} I might have missed it. Thanks for the \textbf{clarification}.} \\ \hline
\textbf{Agent} & \eng{The radio spot is short. Also, \textbf{one important condition}: this \textbf{plan} is for \textbf{prepaid} customers only, and for the first subscription, you must present your \textbf{Student ID}.} \\ \hline
\textbf{Client} & \eng{I have my student card here. What is the \textbf{USSD code} again?} \\ \hline
\textbf{Agent} & \eng{Dial \textbf{*157*99\#}, then select \textbf{Option 1}. You will receive a confirmation \textbf{SMS} instantly.} \\ \hline
\textbf{Client} & \eng{Okay, dialing *157*99\#... Option 1... Done. I just got the SMS: ``\textbf{Subscription successful}. 10 GB activated. Valid for 30 days.''} \\ \hline
\textbf{Agent} & \eng{Perfect. Your \textbf{balance} has been debited. Enjoy your \textbf{internet access} and good luck with your studies!} \\ \hline
\textbf{Client} & \eng{Thank you very much for your help. Have a nice day.} \\ \hline
\textbf{Agent} & \eng{You're welcome. Thank you for choosing MTN.} \\ \hline
\end{tabularx}

\textbf{Commentaire des choix de langue (Pour l'enseignant / Auto-évaluation) :}
\begin{enumerate}
    \item \textbf{Politesse professionnelle :} \eng{Good morning / Welcome to... / How can I help you? / Certainly / Please note that...} (Registre formel attendu au guichet).
    \item \textbf{Vocabulaire technique précis :} \eng{Bundle / Promo / On-net / Off-net / Validity / USSD code / Prepaid / Subscription successful / Balance debited}. L'élève A réemploie exactement les termes de la fiche technique.
    \item \textbf{Structure de correction :} \eng{That is correct... However, please note that...} (Validation partielle + correction ciblée). C'est la structure standard du \emph{Customer Care}.
    \item \textbf{Vérification active (Client) :} \eng{Did the announcement mention that? / What is the USSD code again?} Montre que le client n'est pas passif, il sécurise sa souscription.
    \item \textbf{Grammaire :} \eng{Subscribe \textbf{to}} (préposition) ; \eng{10 GB \textbf{of} data} (indénombrable) ; \eng{Valid \textbf{for} 30 days} (préposition de durée) ; \eng{Charged \textbf{at} 50 FCFA} (préposition de prix/taux).
    \item \textbf{Prononciation (indications pour l'oral) :} \eng{USSD} [« you-ess-es-di »] ; \eng{Gigabyte} [« dji-ga-baït »] ; \eng{Validity} [« va-li-di-ti »] (accent sur 3e syllabe).
\end{enumerate}

\textbf{Partie 3 : Modèle rapide pour le Scénario 2 (Orange Campus Connect — Inversion)}
\begin{itemize}
    \item \textbf{Client (Élève B) :} \eng{``I'd like to activate the \textbf{Campus Connect} offer. I noted 15 GB for 45 days, 3,000 FCFA, with 1 hour \textbf{all networks}.''}
    \item \textbf{Agent (Élève A) :} \eng{``Correct on data and validity. \textbf{Just a precision}: the voice bonus is \textbf{60 minutes total} (on-net + off-net combined), not 1 hour per network. Code is \textbf{*141*50\#}. No ID needed.''}
    \item \textbf{Client :} \eng{``Understood. 60 minutes shared. Dialing *141*50\# now. Confirmed.''}
\end{itemize}

\end{exoresolu}

\courssub{Bilan de l'activité}

\begin{aretenir}[Ce que vous devez retenir de cette intégration]
\begin{itemize}
    \item \textbf{L'écoute sélective} permet de ne noter que l'essentiel (Chiffres, Noms propres, Mots-clés : \eng{Price, Data, Validity, Code}). Ne cherchez pas à tout écrire.
    \item \textbf{La vérification orale} (\emph{checking}) est une compétence vitale : \eng{``Could you confirm...?''}, \eng{``Did you say...?''}. Elle évite les erreurs de facturation ou de souscription.
    \item \textbf{Le lexique technique} (\eng{bundle, on-net, off-net, USSD, validity, prepaid, roll over}) est le passeport pour gérer ses télécoms en anglais (voyages, études à l'étranger, concours).
    \item \textbf{La structure « Validation + Correction polie »} (\eng{That's right, however... / Actually...}) est le standard professionnel anglophone (\emph{Customer Service English}).
    \item \textbf{Contexte camerounais :} Les codes USSD (*157\#, *141\#) et les devises (FCFA) sont des réalités locales. Savoir les dicter et les comprendre en anglais est un atout majeur pour l'insertion professionnelle.
\end{itemize}
\end{aretenir}

\begin{savaistu}[Le saviez-vous ? — Le marché télécoms au Cameroun]
Le Cameroun compte trois opérateurs majeurs : \textbf{MTN Cameroon} (leader historique, filiale du groupe sud-africain MTN Group), \textbf{Orange Cameroun} (filiale d'Orange S.A., France) et \textbf{Camtel} (opérateur historique public, aussi fournisseur d'accès internet via le câble sous-marin WACS). Depuis 2015, la \emph{Number Portability} (portabilité des numéros) permet de changer d'opérateur en gardant son 6xx. Les forfaits « \emph{Student} » ou « \emph{Campus} » sont des offres \emph{segmentées} (\emph{market segmentation}) ciblant la jeunesse (plus de 60\% de la population a moins de 25 ans). La data mobile est le principal moteur de croissance (ARPU \emph{Average Revenue Per User} data > voix). L'anglais est la langue de travail technique dans ces entreprises (réseau, ingénierie, \emph{core network}), d'où l'importance de maîtriser ce vocabulaire dès le lycée.
\end{savaistu}

\newpage

\coursec{Méthodes et exercices résolus}

\coursec{Lesson 49 — Listening: Texts about subscribing to service packages (telephone/internet services)}

\courssub{Avant l'écoute : Préparation et anticipation}

\begin{methode}[Activer ses connaissances et anticiper le contenu]
Avant toute écoute, suivez ces trois étapes pour mettre toutes les chances de votre côté.
\begin{enumerate}
    \item \textbf{Analyser le support visuel / le titre :} S'il y a une image (logo opérateur, smartphone, carte de recharge), identifiez le thème (\cle{telecommunications}, \cle{mobile data}, \cle{subscription}). Le titre \eng{« Subscribing to a 4G Package »} annonce le lexique : \cle{subscribe}, \cle{package}, \cle{bundle}, \cle{data plan}, \cle{recharge}.
    \item \textbf{Mobiliser le lexique thématique :} Faites un \emph{brainstorming} rapide (1 minute) sur les mots anglais associés : \eng{price / cost / rate}, \eng{validity / duration / expiry}, \eng{GB / MB / data allowance}, \eng{activate / dial / USSD code}, \eng{confirmation / SMS / balance}, \eng{prepaid / postpaid}.
    \item \textbf{Prédire les informations chiffrées :} Préparez une grille de prise de notes avec les colonnes : \textsc{Package Name} | \textsc{Price (FCFA)} | \textsc{Data Volume} | \textsc{Validity} | \textsc{USSD Code}. Les chiffres (prix, volumes, codes \eng{*...\#}) sont les cibles prioritaires.
\end{enumerate}
\end{methode}

\begin{experience}[Activité de classe : « Prédictions à partir de mots-clés »]
Le professeur écrit au tableau 10 mots-clés extraits du script (ex. \eng{Agent}, \eng{Customer}, \eng{Daily}, \eng{Weekly}, \eng{Monthly}, \eng{500}, \eng{2000}, \eng{10000}, \eng{dial}, \eng{SMS}). Par binômes, les élèves imaginent le scénario probable (qui parle ? quoi ? combien ?) et rédigent 3 questions qu'ils s'attendent à entendre. Mise en commun orale avant la première écoute.
\end{experience}

\courssub{Le texte support : Script d'écoute}

\begin{textebox}[Script d'écoute : « Subscribing to a 4G Package »]
\eng{Agent:} Good morning, CamTel Plus customer care, my name is Pierre. How may I assist you today?\\
\eng{Customer:} Good morning. My name is Paul Abega. I would like to subscribe to a 4G internet package for my smartphone.\\
\eng{Agent:} Certainly, Mr. Abega. We have three main bundles. The \textbf{Daily Pass} costs 500 FCFA for 1 GB valid for 24 hours. The \textbf{Weekly Bundle} is 2,000 FCFA for 5 GB valid for 7 days. And our best value, the \textbf{Monthly Max}, gives you 25 GB for 10,000 FCFA, valid for 30 days.\\
\eng{Customer:} I think I will go for the Monthly Max. It suits my budget.\\
\eng{Agent:} Excellent choice. Could you please confirm your phone number?\\
\eng{Customer:} Yes, it is 6 77 88 99 00.\\
\eng{Agent:} Thank you. The subscription has been processed. You will receive a confirmation SMS shortly. To activate the data, please dial *155*1\#. Is there anything else I can help you with?\\
\eng{Customer:} No, that's all. Thank you very much. Goodbye.\\
\eng{Agent:} You're welcome, Mr. Abega. Have a nice day.
\end{textebox}

\courssub{Compréhension : Questions variées}

\begin{methode}[Stratégie n°1 : Première écoute globale — Valider les hypothèses et capter l'essentiel]
\begin{enumerate}
    \item \textbf{Ne cherchez pas à tout comprendre.} Concentrez-vous sur : \textbf{Qui ?} (Interlocuteurs), \textbf{Où ?} (Contexte : service client), \textbf{Quoi ?} (Action : souscription), \textbf{Résultat ?} (Choix final + Code activation).
    \item \textbf{Repérez les « mots-ancres » :} Noms propres (\eng{Paul Abega}, \eng{Pierre}), Mots de liaison logiques (\eng{Certainly}, \eng{Excellent choice}, \eng{Thank you}), Chiffres (\eng{500}, \eng{2,000}, \eng{10,000}, \eng{*155*1\#}).
    \item \textbf{Répondez aux questions larges (Wh-questions) :} \eng{What does the customer want?}, \eng{Which package does he choose?}, \eng{What is the activation code?}
\end{enumerate}
\end{methode}

\begin{exoresolu}[Exercice 1 : Prédiction et première écoute globale]
\textbf{Consigne :} Vous allez entendre le dialogue ci-dessus. Après une \textbf{première écoute globale}, répondez aux questions 1 à 3 en anglais.

\textbf{Questions :}
\begin{enumerate}
    \item What is the customer's full name?
    \item Which package does the customer choose?
    \item What is the USSD code to activate the data?
\end{enumerate}

\tcblower
\textbf{Solution détaillée :}
\begin{enumerate}
    \item \textbf{Repérage de l'information (Proper Noun) :} Le client dit : \eng{« My name is Paul Abega. »} $\to$ \textbf{Paul Abega}.
    \item \textbf{Repérage du choix (Mot-clé « go for ») :} Le client dit : \eng{« I will go for the Monthly Max. »} $\to$ \textbf{The Monthly Max (package)}.
    \item \textbf{Repérage du code technique (Chiffres/Symboles) :} L'agent dit : \eng{« dial *155*1\# »} $\to$ \textbf{*155*1\#}.
\end{enumerate}
\textbf{Conclusion :} La première écoute ne demande pas de tout comprendre, mais de valider les hypothèses (contexte : souscription ; interlocuteurs : client/agent) et de capter les \textbf{informations ciblées} (noms, choix, codes).
\end{exoresolu}

\begin{methode}[Stratégie n°2 : Deuxième écoute détaillée — Prise de notes structurées (Tableau comparatif)]
\begin{enumerate}
    \item \textbf{Préparez les colonnes :} \textsc{Package Name} | \textsc{Price (FCFA)} | \textsc{Data Volume} | \textsc{Validity}.
    \item \textbf{Notez en abrégé :} N'écrivez pas de phrases. Ex: \eng{Daily / 500 / 1GB / 24h}.
    \item \textbf{Attention aux distracteurs :} L'énoncé peut citer un prix pour le mauvais forfait, ou inverser la validité.
    \item \textbf{Vérifiez la cohérence :} Le « Monthly Max » est-il vraiment le moins cher au Go ? (Calcul mental rapide : 10 000 / 25 = 400 FCFA/Go vs 500 FCFA/Go pour le Daily).
\end{enumerate}
\end{methode}

\begin{exoresolu}[Exercice 2 : Compréhension détaillée – Vrai / Faux / Non mentionné + Justification]
\textbf{Consigne :} Réécoutez le dialogue. Pour chaque affirmation, cochez \textbf{True (T)}, \textbf{False (F)} ou \textbf{Not Mentioned (NM)}. \textbf{Justifiez chaque réponse par une citation exacte du texte (en anglais).}

\begin{enumerate}
    \item The Daily Pass offers 2 GB of data.
    \item The Weekly Bundle is valid for one week.
    \item The Monthly Max is the cheapest option per day.
    \item Mr. Abega pays with Mobile Money.
    \item The agent's name is Pierre.
\end{enumerate}

\tcblower
\textbf{Solution détaillée et justifications :}

\begin{center}
\begin{tabularx}{\linewidth}{|X|X|X|}
\hline
\rowcolor{popblueL} \textbf{Affirmation} & \textbf{Réponse} & \textbf{Justification (Citation + Traduction)} \\
\hline
1. The Daily Pass offers 2 GB of data. & \textbf{F (False)} & \eng{« The Daily Pass costs 500 FCFA for \textbf{1 GB} »} (Le Daily Pass coûte 500 FCFA pour 1 Go). \\
\hline
2. The Weekly Bundle is valid for one week. & \textbf{T (True)} & \eng{« \textbf{valid for 7 days} »} (valide pendant 7 jours = une semaine). \\
\hline
3. The Monthly Max is the cheapest option per day. & \textbf{NM (Not Mentioned)} & Le texte donne le prix total et la durée, mais \textbf{ne compare pas} le coût journalier des trois offres explicitement. \\
\hline
4. Mr. Abega pays with Mobile Money. & \textbf{NM (Not Mentioned)} & Le texte dit \eng{« The subscription has been processed »} mais \textbf{ne précise pas le mode de paiement} (Mobile Money, carte, facture). \\
\hline
5. The agent's name is Pierre. & \textbf{T (True)} & \eng{« \textbf{my name is Pierre} »} (première réplique de l'agent). \\
\hline
\end{tabularx}
\end{center}

\textbf{Point de méthode :} Pour \textbf{NM}, ne déduisez jamais ! Si l'information n'est pas dans l'audio, c'est NM, même si cela semble logique (ex: paiement Mobile Money très courant au Cameroun).
\end{exoresolu}

\courssub{Lexique thématique : Télécommunications et forfaits}

\begin{methode}[Stratégie n°3 : Maîtriser le lexique technique (Collocations et Verbes opérateurs)]
Le vocabulaire du thème forme des \textbf{blocs lexicaux} (collocations) qu'il faut mémoriser ensemble pour éviter les calques du français.
\begin{enumerate}
    \item \textbf{Verbes d'action :} \cle{to subscribe to} (s'abonner à) $\neq$ \emph{to suscribe} (faute) ; \cle{to recharge / to top up} (recharger) ; \cle{to activate} (activer) ; \cle{to dial} (composer un numéro/code) ; \cle{to unsubscribe / to opt out} (se désabonner).
    \item \textbf{Noms composés clés :} \cle{data allowance / data cap} (plafond de données) ; \cle{validity period} (période de validité) ; \cle{rollover} (report du solde non utilisé) ; \cle{out-of-bundle rate} (tarif hors forfait) ; \cle{USSD code} (code USSD).
    \item \textbf{Adjectifs de description :} \cle{unlimited} (illimité) ; \cle{capped} (plafonné) ; \cle{seamless} (fluide, sans coupure) ; \cle{prepaid / postpaid} (prépayé / postpayé).
    \item \textbf{Structures grammaticales associées :}
    \begin{itemize}
        \item \eng{Thanks \textbf{for calling}.} (Préposition + V-ing)
        \item \eng{You will receive an SMS \textbf{confirming} your subscription.} (Participe présent / Gérondif après nom)
        \item \eng{Please dial *123\# \textbf{to activate} the data.} (Infinitif de finalité)
    \end{itemize}
\end{enumerate}
\end{methode}

\begin{exoresolu}[Exercice 3 : Lexique – Complétion de phrases et Reformulation]
\textbf{Consigne :} Complétez les phrases avec le mot ou l'expression approprié(e) parmi la liste ci-dessous. Mettez le verbe à la forme correcte si nécessaire.
\emph{Liste : \textbf{subscribe to, top up, validity period, rollover, dial, out-of-bundle, prepaid, unlimited.}}

\begin{enumerate}
    \item Most students in Yaoundé use \_\_\_\_\_\_\_\_\_\_ lines because they control their spending better.
    \item If you exceed your data allowance, you will be charged at the \_\_\_\_\_\_\_\_\_\_ rate, which is very expensive.
    \item Don't forget to \_\_\_\_\_\_\_\_\_\_ *155\# to check your balance.
    \item The \_\_\_\_\_\_\_\_\_\_ of the Weekly Bundle is exactly 7 days.
    \item Some operators offer a \_\_\_\_\_\_\_\_\_\_ feature: unused data is added to next month's bundle.
    \item I need to \_\_\_\_\_\_\_\_\_\_ my account before I can buy the Monthly Max package.
    \item Would you like to \_\_\_\_\_\_\_\_\_\_ the « Social Media Pack » for 500 FCFA?
    \item The « Night Owl » package offers \_\_\_\_\_\_\_\_\_\_ downloads between 10 PM and 6 AM.
\end{enumerate}

\tcblower
\textbf{Solution détaillée :}
\begin{enumerate}
    \item \textbf{prepaid} (adj. attribut du nom \eng{lines}). \textit{Astuce :} \eng{prepaid lines} = lignes prépayées.
    \item \textbf{out-of-bundle} (adj. composé invariable devant \eng{rate}). \textit{Piège :} ne pas écrire « out of bundle rate » sans traits d'union.
    \item \textbf{dial} (impératif présent). \textit{Collocation :} \eng{dial a code / a number}.
    \item \textbf{validity period} (nom composé sujet). \textit{Synonyme :} \eng{duration}.
    \item \textbf{rollover} (nom sujet). \textit{Contexte :} « feature » (fonctionnalité) appelle un nom.
    \item \textbf{top up} (verbe phrasal, infinitif après \eng{to}). \textit{Synonyme :} \eng{recharge}.
    \item \textbf{subscribe to} (verbe + préposition \textbf{to} obligatoire). \textit{Erreur fréquente :} \eng{*subscribe the pack}.
    \item \textbf{unlimited} (adjectif épithète de \eng{downloads}). \textit{Collocation :} \eng{unlimited data / calls / downloads}.
\end{enumerate}

\textbf{Reformulation (Entraînement écrit) :}
\eng{Original: « The subscription has been processed. »} $\to$ \textbf{Passif présent perfect.}
\eng{Réécriture active :} « We \textbf{have processed} your subscription. »
\eng{Original: « Please dial *155*1\# to activate the data. »} $\to$ \textbf{Infinitif de finalité.}
\eng{Réécriture :} « Dial *155*1\# \textbf{so that you can activate} the data. » / « \textbf{In order to activate} the data, dial *155*1\#. »
\end{exoresolu}

\begin{propriete}[Tableau récapitulatif : Vocabulaire clé « Telecom Packages »]
\begin{center}
\begin{tabularx}{\linewidth}{|l|X|X|}
\hline
\rowcolor{popblueL} \textbf{English Term} & \textbf{French Translation} & \textbf{Context / Collocation} \\
\hline
\eng{Bundle / Package} & Forfait / Pack & \eng{a Weekly Bundle}, \eng{subscribe to a package} \\
\eng{Data allowance / Cap} & Volume de données / Plafond & \eng{exceed your data allowance}, \eng{capped at 5 GB} \\
\eng{Validity period / Expiry date} & Durée de validité / Date d'expiration & \eng{valid for 30 days}, \eng{expiry date: 15/12/2025} \\
\eng{Rollover / Carry over} & Report (solde) & \eng{unused data rolls over} \\
\eng{Out-of-bundle rate} & Tarif hors forfait & \eng{charged at the out-of-bundle rate} \\
\eng{USSD code / Short code} & Code USSD / Code court & \eng{dial *155*1\#} \\
\eng{Top-up / Recharge (n.)} & Rechargement (carte, crédit) & \eng{buy a top-up}, \eng{a recharge card} \\
\eng{Prepaid / Postpaid} & Prépayé / Postpayé (facturé) & \eng{a prepaid line}, \eng{postpaid contract} \\
\eng{Seamless connection} & Connexion fluide / sans coupure & \eng{enjoy a seamless experience} \\
\hline
\end{tabularx}
\end{center}
\end{propriete}

\courssub{Production guidée : Réemploi du lexique}

\begin{methode}[Stratégie n°4 : Jeu de rôle « Service Client » — Canevas rigide]
L'examen (épreuve orale ou écrite) peut demander de \textbf{reproduire} ou de \textbf{continuer} l'échange. Suivez ce canevas pour ne pas sortir du sujet.
\begin{enumerate}
    \item \textbf{Salutation pro :} \eng{Good morning/afternoon. [Company Name], [Your Name] speaking. How can I help you?}
    \item \textbf{Identification du besoin :} \eng{I'd like to subscribe to... / I want to buy a... / I need more data.}
    \item \textbf{Présentation de l'offre (Agent) :} \textbf{Minimum 2 options} avec \textbf{Prix, Volume, Validité}. Utilisez : \eng{We have... / For [price], you get... / It is valid for...}
    \item \textbf{Choix et confirmation (Client) :} \eng{I'll take the [Name] please. / The [Name] sounds good. My number is...}
    \item \textbf{Clôture technique (Agent) :} \eng{Done. Dial [Code] to activate. You will receive an SMS. Anything else?}
    \item \textbf{Politesse finale :} \eng{Thank you. Goodbye.}
\end{enumerate}
\end{methode}

\begin{exoresolu}[Exercice 4 : Production guidée – Écriture d'un dialogue complet]
\textbf{Consigne :} Vous êtes \textbf{Marie}, agente chez \eng{Orange Cameroon}. Un client (M. Ndi) appelle pour un forfait \textbf{hebdomadaire (Weekly)}. Il veut 3 Go pour 1 500 FCFA. Le code d'activation est \eng{*141*2\#}. Rédigez le dialogue complet (10-12 répliques). Utilisez le canevas de la méthode.

\tcblower
\textbf{Proposition de dialogue type (Modèle corrigé) :}

\begin{textebox}[Dialogue : Weekly Subscription]
\eng{Agent:} Good morning, Orange Cameroon, Marie speaking. How may I help you?\\
\eng{Client:} Good morning. My name is Jean Ndi. I would like to subscribe to a weekly internet bundle.\\
\eng{Agent:} Certainly, Mr. Ndi. We have a Weekly Bundle at 1,500 FCFA. It gives you 3 GB of data valid for 7 days.\\
\eng{Client:} That is exactly what I need. I will take that one.\\
\eng{Agent:} Great. May I have your phone number please?\\
\eng{Client:} Yes, it is 6 99 11 22 33.\\
\eng{Agent:} Thank you. Your subscription is now active. You will receive a confirmation message shortly.\\
\eng{Client:} Do I need to dial a code?\\
\eng{Agent:} Yes, please dial *141*2\# to start using your data immediately.\\
\eng{Client:} Alright. Thank you very much for your help.\\
\eng{Agent:} You are welcome, Mr. Ndi. Enjoy your connection. Goodbye!\\
\eng{Client:} Goodbye.
\end{textebox}

\textbf{Analyse linguistique du modèle :}
\begin{itemize}
    \item \textbf{Politesse professionnelle :} \eng{May I have...?} (plus formel que \eng{Can I have...?}).
    \item \textbf{Présentation de l'offre :} \eng{It gives you... valid for...} (Structure \textbf{Sujet + Verbe + Objet + Complément de durée}).
    \item \textbf{Choix :} \eng{I will take that one} (Décision spontanée $\to$ \textbf{Will}).
    \item \textbf{Confirmation :} \eng{Your subscription is now active} (Passif d'état / Adjectif \eng{active}).
    \item \textbf{Instruction technique :} \eng{Please dial... to start using...} (Infinitif de finalité + \textbf{-ing} après \eng{start}).
\end{itemize}
\end{exoresolu}

\begin{methode}[Stratégie n°5 : Transfert de compétences – Comprendre un SMS de confirmation / Notification USSD]
Au Bac, le support peut être un \textbf{texte court écrit} (SMS, écran USSD, email de confirmation) et non un audio. Les stratégies restent identiques : repérer \textbf{Expéditeur, Action, Montant, Solde, Code, Date limite}.
\begin{enumerate}
    \item \textbf{Identifiez la source :} \eng{From: CamTel Plus} / \eng{USSD Menu}.
    \item \textbf{Isolez la transaction :} \eng{You have subscribed to Monthly Max.}
    \item \textbf{Extrayez les données critiques :} \eng{Cost: 10,000 FCFA. Data: 25 GB. Expiry: 30/11/2025.}
    \item \textbf{Repérez l'action requise :} \eng{Dial *155\# to check balance.} / \eng{Reply STOP to unsubscribe.}
\end{enumerate}
\end{methode}

\begin{exoresolu}[Exercice 5 : Compréhension de notification écrite (SMS / USSD)]
\textbf{Consigne :} Lisez la notification USSD ci-dessous. Répondez aux questions en phrases complètes en anglais.

\begin{textebox}[USSD Notification Screen]
\eng{Welcome to CamTel Plus}\\
\eng{Subscription Confirmed}\\
\eng{Package: Monthly Max (25 GB)}\\
\eng{Amount Debited: 10,000 FCFA}\\
\eng{New Main Balance: 2,500 FCFA}\\
\eng{Data Balance: 25.00 GB}\\
\eng{Expiry Date: 15/12/2025}\\
\eng{Dial *155*0\# for Data Balance | Dial *155*9\# to Unsubscribe}\\
\eng{Thank you for choosing CamTel Plus!}
\end{textebox}

\textbf{Questions :}
\begin{enumerate}
    \item What package has the user just subscribed to?
    \item How much money was deducted from the account?
    \item How much credit remains on the main account?
    \item Until when is the data valid?
    \item What code must the user dial if he wants to cancel the subscription?
\end{enumerate}

\tcblower
\textbf{Solution détaillée :}
\begin{enumerate}
    \item \textbf{The user has subscribed to the « Monthly Max » package (25 GB).} (Ligne \eng{Package: Monthly Max (25 GB)}).
    \item \textbf{10,000 FCFA was debited.} (Ligne \eng{Amount Debited: 10,000 FCFA}). \textit{Note :} Passif \eng{was debited} car l'argent est retiré \emph{par} l'opérateur.
    \item \textbf{The remaining main balance is 2,500 FCFA.} (Ligne \eng{New Main Balance: 2,500 FCFA}).
    \item \textbf{The data is valid until December 15th, 2025.} (Ligne \eng{Expiry Date: 15/12/2025}). \textit{Format date UK :} 15/12 = 15 Décembre.
    \item \textbf{He must dial *155*9\# to unsubscribe.} (Ligne \eng{Dial *155*9\# to Unsubscribe}).
\end{enumerate}

\textbf{Point de grammaire (Passif / Causatif) :}
\eng{Amount Debited} = \eng{(Which was) Debited} (Participe passé adjectival).
\eng{You have subscribed} = \textbf{Present Perfect} (Action passée, conséquence présente).
\eng{Amount deducted} $\neq$ \eng{Amount deducted \textbf{by you}} (C'est l'opérateur qui déduit).
\end{exoresolu}

\courssub{Erreurs fréquentes et pièges du Bac}

\begin{attention}[Les 5 erreurs qui coûtent des points au Bac]
\begin{enumerate}
    \item \textbf{Faux ami / Orthographe : \eng{*Suscribe} au lieu de \eng{Subscribe}.} 
    \textit{Règle :} Double \textbf{B} après le \textbf{U} (\eng{sub-} + \eng{scribe}). Idem pour \eng{subscription} (pas *suscription).
    \item \textbf{Calque du français (Préposition) : \eng{*Subscribe a package} ou \eng{*Subscribe for a package}.} 
    \textit{Correct :} \eng{\textbf{Subscribe TO} a package / a plan / a service.} (Pensez à \eng{listen TO}).
    \item \textbf{Confusion \eng{Recharge} / \eng{Top up} (Verbe vs Nom) :} 
    \textit{Erreur :} \eng{« I want to make a recharge. »} (Anglais africain courant mais non standard). 
    \textit{Standard :} \eng{« I want to \textbf{top up} / \textbf{recharge} my account. »} (Verbe) ou \eng{« I want to buy a \textbf{top-up} / a \textbf{recharge card}. »} (Nom).
    \item \textbf{Ordre des mots dans la question indirecte / polie : \eng{*Can you tell me what is my balance?}} 
    \textit{Correct :} \eng{« Can you tell me \textbf{what my balance is}? »} (Sujet + Verbe, pas d'inversion, pas de \eng{do/does/did}).
    \item \textbf{Mauvaise lecture des chiffres / Dates (Piège USSD/SMS) :} 
    \textit{Erreur :} Lire \eng{10,000} « ten thousands » (pas de s à thousand) ou confondre \eng{15/12} (15 Dec UK) et \eng{12/15} (Dec 15 US). 
    \textit{Règle :} Au Cameroun (anglais britannique) : \textbf{JJ/MM/AAAA}. \eng{Expiry: 05/06/2025} = 5 Juin 2025 (pas le 6 Mai).
\end{enumerate}
\end{attention}

\begin{savaistu}[Le saviez-vous ? Le Cameroun et la révolution Mobile Money]
Le Cameroun est l'un des pays d'Afrique centrale où le \eng{Mobile Money} (Orange Money, MTN MoMo, Express Union Mobile) a transformé l'accès aux forfaits internet. La plupart des souscriptions (\eng{subscriptions}) se font désormais via USSD (\eng{*150\#, *123\#, *155\#}) sans passer par un agent. Le vocabulaire \eng{top up}, \eng{cash in}, \eng{cash out}, \eng{transfer}, \eng{buy bundle} est devenu du langage courant. Au Bac, un texte peut décrire une transaction 100\% Mobile Money : \eng{« Dial *150\#, select 'Buy Bundle', choose 'Weekly', enter PIN, confirm. »} Maîtriser cet anglais « procédural » (impératifs, séquencement : \eng{First, Then, Finally}) est un atout majeur.
\end{savaistu}

\newpage

\coursec{Exercices}

\coursec{Lesson 49 — Listening: Texts about Subscribing to Service Packages (Telephone/Internet Services)}

\courssub{Avant l'écoute : Préparation et anticipation}

\begin{methode}[Stratégies d'écoute : Se préparer efficacement]
\begin{enumerate}
\item \textbf{Analyser le contexte :} Lire le titre, regarder l'image ou le support visuel, identifier les interlocuteurs (agent/client, animateur radio) et le lieu (agence, radio, domicile).
\item \textbf{Anticiper le lexique :} Activer le vocabulaire thématique attendu : \cle{subscribe}, \cle{bundle/pack}, \cle{data}, \cle{validity}, \cle{USSD code}, \cle{dial}, \cle{activate}, \cle{balance}, \cle{top-up}, \cle{coverage}, \cle{unlimited}, \cle{renew}, \cle{throttling}, \cle{refund}.
\item \textbf{Préparer la grille de prise de notes :} Dessiner un tableau avec les colonnes essentielles : \emph{Offre / Prix / Volume data / Validité / Code USSD / Avantages / Conditions}.
\item \textbf{Écoute 1 (Globale) :} Identifier le thème principal, le ton, le nombre d'intervenants, l'intention communicative (informer, vendre, réclamer).
\item \textbf{Écoute 2 (Détaillée) :} Relever les chiffres (prix, Go, jours, minutes), les noms propres (MTN, Orange, Camtel, Yaoundé, Bafia), les mots-clés de liaison (\eng{however}, \eng{plus}, \cle{remember}, \cle{don't miss out}).
\item \textbf{Vérification :} Confirmer ou infirmer ses hypothèses, compléter le tableau.
\end{enumerate}
\end{methode}

\begin{experience}[Activité de mise en route : Brainstorming « Forfaits mobiles »]
En classe entière, au tableau, construisez un nuage de mots (mind map) autour de \cle{Mobile Plan / Bundle}. L'enseignant note les propositions des élèves en anglais (ex. \eng{data}, \eng{credit}, \eng{SMS}, \eng{voice call}, \eng{roaming}, \eng{4G/5G}, \eng{weekly}, \eng{monthly}). Puis classez-les en catégories : \textbf{Cost}, \textbf{Volume}, \textbf{Duration}, \textbf{Features}, \textbf{Actions}.
\end{experience}

\courssub{Le texte support : Scripts d'écoute}

\begin{textebox}[Script 1 : Radio Campus FM — Offre « Student Connect »]
Good morning, students of Yaoundé! This is Radio Campus FM. Are you tired of running out of data during your online research? MTN Cameroon launches the new “Student Connect” weekly bundle. For only 1,000 FCFA, you get 2 GB of high-speed internet valid for 7 days. To subscribe, dial *157*10\# and press 1. Activation is instant. You can check your remaining balance by dialing *157*0\#. Remember, this offer is exclusive to students with a valid student ID. Don't miss out! Stay connected, stay smart.
\end{textebox}

\begin{textebox}[Script 2 : Dialogue — Souscription forfait « Family Pack » (Agence Orange Bonanjo)]
\textbf{Agent:} Good morning, sir. Welcome to Orange Bonanjo. How can I help you?
\textbf{Client:} Good morning. I want to subscribe to the “Family Pack” for my household.
\textbf{Agent:} Excellent choice. The Family Pack costs 5,000 FCFA per month. It gives you 20 GB of data shared between up to 4 numbers, plus unlimited calls to Orange numbers.
\textbf{Client:} 20 GB for the whole family? Is that enough for video streaming?
\textbf{Agent:} For standard definition, yes. For HD streaming, you might need to top up. The pack is valid for 30 days. To activate, dial *141*5\# and select option 2.
\textbf{Client:} How long does activation take?
\textbf{Agent:} It takes a maximum of 5 minutes. You will receive a confirmation SMS.
\textbf{Client:} Great. I'll take it.
\end{textebox}

\begin{textebox}[Script 3 : Camtel Fibre Offres — Basic, Family, Premium]
Camtel presents three fibre optic packages for homes in Douala and Yaoundé. The “Basic” offer costs 15,000 FCFA monthly. It provides 10 Mbps speed and 100 GB data cap. The “Family” offer is 25,000 FCFA per month. Speed reaches 50 Mbps with 500 GB data. It includes a free landline. The “Premium” offer costs 45,000 FCFA monthly. You get 100 Mbps speed, unlimited data, a free landline and a free Wi-Fi 6 router. Installation is free for all packs if you sign a 12-month contract.
\end{textebox}

\courssub{Compréhension : Questions variées}

\begin{exercice}[Écoute : Annonce « Student Connect » — Vrai ou Faux justifié]
Écoutez l'annonce radiophonique (Script 1, lue deux fois par l'enseignant) et dites si les affirmations sont vraies (V) ou fausses (F). Justifiez chaque réponse par une phrase courte en anglais.
\begin{enumerate}
\item The “Student Connect” bundle costs 1,500 FCFA per week.
\item The data volume offered is 2 GB.
\item The validity period is 30 days.
\item The USSD code to subscribe is *157*10\#.
\item The offer is available to all MTN customers without restriction.
\end{enumerate}
\end{exercice}

\begin{exercice}[Dialogue client/agent — QCM]
Écoutez le dialogue (Script 2, lu deux fois). Choisissez la bonne réponse parmi les trois propositions.
\begin{enumerate}
\item What is the monthly cost of the Family Pack?
      \begin{itemize}
      \item[a)] 2,500 FCFA
      \item[b)] 5,000 FCFA
      \item[c)] 10,000 FCFA
      \end{itemize}
\item How much data does the pack provide?
      \begin{itemize}
      \item[a)] 10 GB
      \item[b)] 20 GB
      \item[c)] 50 GB
      \end{itemize}
\item How many numbers can share the data?
      \begin{itemize}
      \item[a)] 2
      \item[b)] 4
      \item[c)] 5
      \end{itemize}
\item What is the USSD code to activate the pack?
      \begin{itemize}
      \item[a)] *141*5\#
      \item[b)] *157*10\#
      \item[c)] *111*1\#
      \end{itemize}
\item What is the maximum activation delay mentioned by the agent?
      \begin{itemize}
      \item[a)] Instant
      \item[b)] 5 minutes
      \item[c)] 24 hours
      \end{itemize}
\end{enumerate}
\end{exercice}

\begin{exercice}[Complétion de tableau récapitulatif — Trois offres Camtel]
Écoutez le script (Script 3, lu deux fois) présentant trois offres de Camtel. Complétez le tableau ci-dessous en anglais.
\begin{center}
\begin{tabularx}{\linewidth}{|X|c|c|c|}\hline
\rowcolor{popblueL} \textbf{Offre} & \textbf{Prix mensuel (FCFA)} & \textbf{Vitesse / Données} & \textbf{Avantages inclus} \\ \hline
Basic  &  &  &  \\ \hline
Family &  &  &  \\ \hline
Premium &  &  &  \\ \hline
\end{tabularx}
\end{center}
\end{exercice}

\courssub{Lexique thématique : Télécommunications et forfaits}

\begin{aretenir}[Vocabulaire clé — Subscribing to Service Packages]
\begin{center}
\begin{tabularx}{\linewidth}{|X|X|X|}\hline
\rowcolor{popblueL} \textbf{English Term} & \textbf{Français} & \textbf{Exemple / Contexte} \\ \hline
\cle{Bundle} / \cle{Pack} & Forfait, pack & \eng{A weekly data bundle} \\ \hline
\cle{Subscribe (to)} & S'abonner (à) & \eng{Subscribe to the Family Pack} \\ \hline
\cle{Unsubscribe (from)} / \cle{Opt out} & Se désabonner (de) & \eng{Unsubscribe from the service} \\ \hline
\cle{Renew} & Renouveler & \eng{Renew your subscription before expiry} \\ \hline
\cle{Auto-renewal} & Renouvellement automatique & \eng{Disable auto-renewal in settings} \\ \hline
\cle{Validity} / \cle{Validity period} & Validité / Durée de validité & \eng{The validity period is 30 days} \\ \hline
\cle{Expire} & Expirer & \eng{Unused data will expire tomorrow} \\ \hline
\cle{Expiry date} & Date d'expiration & \eng{Check the expiry date on the SMS} \\ \hline
\cle{Data cap} / \cle{Cap} & Plafond de données & \eng{A 100 GB data cap} \\ \hline
\cle{Unlimited} & Illimité & \eng{Unlimited calls to same network} \\ \hline
\cle{Throttling} & Bridage (réduction débit) & \eng{Speed throttling after 50 GB} \\ \hline
\cle{Top-up} / \cle{Recharge} & Rechargement / Recharger & \eng{I need a 1,000 FCFA top-up} \\ \hline
\cle{Balance} & Solde (crédit/data) & \eng{Check your data balance} \\ \hline
\cle{USSD code} & Code USSD & \eng{Dial the USSD code *157\#} \\ \hline
\cle{Dial} & Composer (un numéro/code) & \eng{Dial *141*5\# to activate} \\ \hline
\cle{Activate} / \cle{Activation} & Activer / Activation & \eng{Activation is instant} \\ \hline
\cle{Confirmation SMS} & SMS de confirmation & \eng{You will receive a confirmation SMS} \\ \hline
\cle{Coverage} / \cle{Network coverage} & Couverture réseau & \eng{4G coverage in rural areas} \\ \hline
\cle{Landline} & Ligne fixe & \eng{Free landline included} \\ \hline
\cle{Broadband} / \cle{Fibre optic} & Haut débit / Fibre optique & \eng{Fibre optic backbone} \\ \hline
\cle{Subscriber} / \cle{User} & Abonné / Utilisateur & \eng{Rural subscribers face challenges} \\ \hline
\cle{Affordability} & Accessibilité financière & \eng{Affordability remains an issue} \\ \hline
\cle{Digital literacy} & Culture numérique / Illectronisme & \eng{Improve digital literacy for elders} \\ \hline
\end{tabularx}
\end{center}
\end{aretenir}

\begin{attention}[Faux amis et pièges fréquents]
\begin{itemize}
\item \textbf{Subscribe \emph{vs} Subscribe to :} On s'abonne \textbf{à} quelque chose (\eng{subscribe \textbf{to} a pack}). Ne pas oublier la préposition.
\item \textbf{Validity \emph{vs} Validity period :} \cle{Validity} est le nom abstrait (la validité), \cle{Validity period} est la durée. Ne pas dire \eng{*The validity is 30 days} mais \eng{The validity \textbf{period} is 30 days} ou \eng{It is valid \textbf{for} 30 days}.
\item \textbf{Expire (verbe) \emph{vs} Expiry (nom) :} \eng{The pack \textbf{expires} tomorrow} (verbe) / \eng{The \textbf{expiry date} is tomorrow} (nom).
\item \textbf{Recharge (nom/verbe) \emph{vs} Top-up (nom) :} \eng{Do a top-up} / \eng{Top up your account}. \cle{Recharge} s'utilise aussi comme verbe (\eng{Recharge your line}).
\item \textbf{Data :} Indénombrable. \eng{How \textbf{much} data?} (pas \eng{*How many data?}). Unités : \eng{MB, GB}.
\item \textbf{Coverage :} Indénombrable. \eng{Good \textbf{coverage}} (pas \eng{*a good coverage}).
\item \textbf{USSD :} \emph{Unstructured Supplementary Service Data}. Se prononce « you-ess-es-di ».
\end{itemize}
\end{attention}

\courssub{Méthode d'écoute : Stratégies et exemple traité}

\begin{methode}[Exemple traité : Repérer les chiffres et noms propres (Script 1)]
\textbf{Consigne :} « Combien coûte le forfait ? Quelle est la durée de validité ? Quel code composer ? »
\begin{enumerate}
\item \textbf{Écoute 1 :} J'entends « MTN », « students », « Yaoundé », « 1,000 », « 2 GB », « 7 days », « *157*10\# ». Je note : \emph{Opérateur=MTN, Cible=Étudiants, Prix=1000, Data=2GB, Durée=7j, Code=*157*10\#}.
\item \textbf{Écoute 2 :} Je vérifie : « 1,000 FCFA » (pas 1,500). « 7 days » (pas 30). « *157*10\# » (confirmé). « Exclusive to students » (restriction).
\item \textbf{Rédaction des réponses :} J'écris des phrases complètes : \eng{The bundle costs 1,000 FCFA. It is valid for 7 days. The USSD code is *157*10\#.}
\end{enumerate}
\end{methode}

\courssub{Grammaire en contexte : Temps verbaux et formes}

\begin{propriete}[Rappel : Temps verbaux fréquents dans les procédures télécoms]
\begin{center}
\begin{tabularx}{\linewidth}{|X|X|X|}\hline
\rowcolor{popblueL} \textbf{Temps / Forme} & \textbf{Emploi (Contexte Telecom)} & \textbf{Exemples} \\ \hline
\textbf{Present Simple} & Habitudes, faits généraux, procédures fixes & \eng{Customers \textbf{subscribe} every Monday.} \\ \hline
\textbf{Present Perfect} & Action passée avec résultat présent (\cle{just}, \cle{already}, \cle{yet}) & \eng{The agent \textbf{has just activated} your line.} \\ \hline
\textbf{Future Simple (\cle{will})} & Prédiction, promesse, décision instantanée & \eng{The company \textbf{will launch} 5G soon.} \\ \hline
\textbf{Imperatif} & Instructions, mode d'emploi, consignes & \eng{\textbf{Dial} *157\#. \textbf{Do not forget} to renew.} \\ \hline
\textbf{Type 1 Conditional (\cle{If} + Present, \cle{will} + BV)} & Condition réelle future (procédure USSD) & \eng{If you \textbf{dial} *157\#, you \textbf{will see} the menu.} \\ \hline
\textbf{Present Perfect (\cle{for/since} / \cle{How long})} & Durée depuis le passé jusqu'à maintenant & \eng{How long \textbf{have you had} this number?} \\ \hline
\end{tabularx}
\end{center}
\end{propriete}

\begin{exercice}[Conjugaison et formes verbales (Contexte télécoms)]
Complétez les phrases suivantes en conjuguant le verbe entre parenthèses au temps approprié (Present Simple, Present Perfect, Future Simple, Impératif, Conditionnel Type 1).
\begin{enumerate}
\item Customers \textbf{\_\_\_\_\_\_\_\_\_\_} (subscribe) to the weekly bundle every Monday.
\item The agent \textbf{\_\_\_\_\_\_\_\_\_\_} (just / activate) your line. You can start browsing now.
\item If you \textbf{\_\_\_\_\_\_\_\_\_\_} (dial) *157\#, you \textbf{\_\_\_\_\_\_\_\_\_\_} (see) the main menu.
\item \textbf{\_\_\_\_\_\_\_\_\_\_} (not / forget) to renew your pack before it expires tomorrow.
\item By the end of the year, the company \textbf{\_\_\_\_\_\_\_\_\_\_} (launch) 5G in Buea and Bamenda.
\item How long \textbf{\_\_\_\_\_\_\_\_\_\_} (you / have) this phone number?
\item The data balance \textbf{\_\_\_\_\_\_\_\_\_\_} (expire) in two days if you don't use it.
\item Please \textbf{\_\_\_\_\_\_\_\_\_\_} (enter) your PIN to confirm the transaction.
\end{enumerate}
\end{exercice}

\courssub{Production guidée : Réemploi du lexique}

\begin{exercice}[Traduction : Phrases thématiques (Français $\rightarrow$ Anglais)]
Traduisez en anglais en réemployant le vocabulaire de la leçon (\cle{subscribe}, \cle{bundle}, \cle{top up}, \cle{balance}, \cle{coverage}, \cle{unlimited}, \cle{validity}, \cle{renew}).
\begin{enumerate}
\item Je veux m'abonner au forfait mensuel illimité.
\item Combien coûte le rechargement de 5 Go ?
\item La validité de ce forfait est de 30 jours.
\item Vérifiez votre solde avant de regarder une vidéo en streaming.
\item La couverture réseau est excellente dans le centre-ville de Yaoundé.
\item Vous devez renouveler votre abonnement avant la date d'expiration.
\item Ce forfait famille offre des appels illimités vers le même opérateur.
\item Quel est le code USSD pour souscrire à cette offre ?
\end{enumerate}
\end{exercice}

\begin{exercice}[Transformation de phrases]
Transformez les phrases selon les indications entre parenthèses.
\begin{enumerate}
\item ``Dial *157\# to check your balance.'' (Passif : \emph{Your balance \_\_\_\_\_\_\_\_\_\_})
\item ``The activation is instant.'' (Négation)
\item ``You get 10 GB for 2,000 FCFA.'' (Question sur le prix : \emph{How much \_\_\_\_\_\_\_\_\_\_?})
\item ``He subscribed to the Premium pack last week.'' (Present Perfect avec \emph{just})
\item ``If you top up today, you get a bonus.'' (Inversion conditionnelle Type 1 : \emph{Should you \_\_\_\_\_\_\_\_\_\_})
\item ``The agent explained the procedure to me.'' (Passif : \emph{I \_\_\_\_\_\_\_\_\_\_})
\item ``This offer expires tomorrow.'' (Future Simple avec \emph{will})
\item ``Do you have a student ID?'' (Discours indirect : \emph{He asked if \_\_\_\_\_\_\_\_\_\_})
\end{enumerate}
\end{exercice}

\begin{exercice}[Texte à trous : Scénario de souscription]
Complétez le texte suivant avec les mots de la liste. Il y a deux mots en trop.
\emph{\{ activate, coverage, dial, expire, unlimited, bundle, top up, balance, subscribe, validity, router, confirmation \}}
\begin{textebox}[Complétion]
To \textbf{(1)} \_\_\_\_\_\_\_\_\_\_ to the new ``Youth Boost'' offer, customers must \textbf{(2)} \_\_\_\_\_\_\_\_\_\_ *145*20\# on their phone. The \textbf{(3)} \_\_\_\_\_\_\_\_\_\_ costs 1,500 FCFA and includes 5 GB of data with a \textbf{(4)} \_\_\_\_\_\_\_\_\_\_ of 14 days. Calls to the same network are \textbf{(5)} \_\_\_\_\_\_\_\_\_\_. Once you subscribe, you will receive a \textbf{(6)} \_\_\_\_\_\_\_\_\_\_ SMS. To \textbf{(7)} \_\_\_\_\_\_\_\_\_\_ the service, simply restart your phone. You can check your data \textbf{(8)} \_\_\_\_\_\_\_\_\_\_ at any time by dialing *145*0\#. Remember that unused data will \textbf{(9)} \_\_\_\_\_\_\_\_\_\_ after the validity period. If you travel to the Far North region, check the network \textbf{(10)} \_\_\_\_\_\_\_\_\_\_ map on the operator's website.
\end{textebox}
\end{exercice}

\begin{exercice}[Appariements : Termes et définitions / Étapes USSD]
Associez chaque terme de la colonne A à sa définition ou action correspondante dans la colonne B.
\begin{center}
\begin{tabularx}{\linewidth}{|X|X|}\hline
\rowcolor{popblueL} \textbf{Colonne A} & \textbf{Colonne B} \\ \hline
1. USSD Code & A. Amount of money added to the account to buy data or credit. \\ \hline
2. Top-up / Recharge & B. Short code (e.g. *157\#) used to access mobile services. \\ \hline
3. Data Cap / Cap & C. The date when the package stops working. \\ \hline
4. Expiry Date & D. Maximum amount of data allowed in a package. \\ \hline
5. Bundle / Pack & E. A package combining data, voice and/or SMS for a fixed price. \\ \hline
6. Coverage & F. The geographical area where the network signal is available. \\ \hline
7. Balance & G. The remaining amount of data or credit. \\ \hline
8. Auto-renewal & H. Automatic subscription renewal at the end of the validity period. \\ \hline
\end{tabularx}
\end{center}
\textit{Écrivez les paires correctes (ex: 1-B, 2-A, ...).}
\end{exercice}

\courssub{Je me prépare au Bac}

\begin{exercice}[Compréhension de texte long + Questions de langue — Type Baccalauréat]
Lisez attentivement le texte suivant et répondez aux questions qui suivent.
\begin{textebox}[Article : « Bridging the Digital Divide: Cameroon's Rural Connectivity Challenge »]
In the bustling streets of Douala, 4G coverage is almost ubiquitous. Street vendors, students, and taxi drivers seamlessly stream videos, make WhatsApp calls, and mobile money transactions. However, drive 50 kilometres east towards the rural subdivisions of the Centre Region, and the digital landscape changes drastically. Here, network coverage is patchy, and internet speed often crawls at 2G levels.

This disparity is the focus of the government's “National Broadband Plan 2025”. The plan aims to extend fibre optic backbone to 360 subdivision headquarters. Private operators like MTN, Orange, and Camtel are incentivised to deploy base stations in underserved zones through tax breaks and universal service fund subsidies.

For the average rural subscriber, the challenge is not only coverage but affordability. A 1 GB daily bundle costs 500 FCFA — a significant portion of a farmer's daily income. Operators have introduced “social bundles”: weekly packages of 500 MB for 200 FCFA, accessible via USSD codes like *157*99\#. Yet, digital literacy remains a barrier. Many elders do not know how to \cle{subscribe}, \cle{renew}, or \cle{check their balance} without assistance.

In Bafia, a community digital centre funded by a partnership between the Ministry of Posts and Telecommunications and a local NGO offers free Wi-Fi and training. “Before, I had to travel to Yaoundé to send an email,” says Marie, a 42-year-old trader. “Now, I top up my phone here and sell my cocoa to buyers in Douala via WhatsApp. My income has doubled.”

Experts argue that infrastructure alone is insufficient. Content in local languages (Bassa, Ewondo, Fulfulde) and relevant e-government services (birth certificates, tax payment) are crucial to drive adoption. As Cameroon races towards its 2035 emergence vision, connecting the “last mile” is not just a technical issue — it is an economic imperative.
\end{textebox}

\textbf{A. Comprehension Questions (Answer in English, using your own words as much as possible).}
\begin{enumerate}
\item Contrast the internet experience in Douala with that in rural Centre Region. (2 pts)
\item What is the goal of the “National Broadband Plan 2025” regarding fibre optic deployment? (1 pt)
\item Name two incentives offered to private operators to deploy in rural areas. (2 pts)
\item Why is the 1 GB daily bundle considered unaffordable for a rural farmer? (1 pt)
\item What solution do operators offer for low-income users? Give the price and volume. (2 pts)
\item According to the text, what prevents elders from using mobile internet independently? (1 pt)
\item How has the community digital centre in Bafia changed Marie's business? (2 pts)
\item Identify two non-technical factors experts say are necessary to drive internet adoption in rural areas. (2 pts)
\end{enumerate}

\textbf{B. Language Work}
\begin{enumerate}
\item Find in the text words or phrases that mean the same as: (2 pts)
      \begin{itemize}
      \item[a)] existing everywhere (par. 1)
      \item[b)] financial help from the government (par. 2)
      \item[c)] obstacle / barrier (par. 4)
      \item[d)] essential / vital (par. 6)
      \end{itemize}
\item Rewrite the following sentences as indicated: (3 pts)
      \begin{itemize}
      \item[a)] ``Operators have introduced social bundles.'' (Passive voice)
      \item[b)] ``Many elders do not know how to subscribe.'' (Use: \emph{Many elders are unable \_\_\_\_\_\_\_\_\_\_})
      \item[c)] ``If the government builds more towers, coverage will improve.'' (Begin with: \emph{Should \_\_\_\_\_\_\_\_\_\_})
      \end{itemize}
\item Put the verbs in brackets into the correct tense/form: (3 pts)
      By the time the plan \textbf{(achieve)} its goals in 2025, thousands of kilometres of fibre \textbf{(lay)}. Currently, the Universal Service Fund \textbf{(finance)} several pilot projects. If the projects \textbf{(be)} successful, rural GDP \textbf{(rise)} significantly.
\end{enumerate}
\end{exercice}

\begin{exercice}[Traduction : Paragraphe Anglais $\rightarrow$ Français]
Traduisez le paragraphe suivant en français correct et naturel.
\begin{textebox}[Extrait à traduire]
``For the average rural subscriber, the challenge is not only coverage but affordability. A 1 GB daily bundle costs 500 FCFA — a significant portion of a farmer's daily income. Operators have introduced ``social bundles'': weekly packages of 500 MB for 200 FCFA, accessible via USSD codes like *157*99\#. Yet, digital literacy remains a barrier. Many elders do not know how to subscribe, renew, or check their balance without assistance.''
\end{textebox}
\end{exercice}

\begin{exercice}[Expression écrite : Sujet de composition — Type Bac]
Choisissez \textbf{UN} des deux sujets ci-dessous. Rédigez une réponse de \textbf{150 à 180 mots} environ. Votre production sera évaluée sur la pertinence, la cohérence, la correction linguistique (grammaire, vocabulaire, orthographe) et la richesse du vocabulaire thématique (télécommunications, forfaits, USSD, etc.).

\textbf{Sujet 1 : Article d'opinion pour le journal de l'école}
``Mobile internet packages are too expensive for Cameroonian students. Write an article for your school magazine arguing for or against this statement. Propose concrete solutions (student bundles, campus Wi-Fi, partnerships with operators).''

\textbf{Sujet 2 : Lettre formelle / Email de réclamation}
``You subscribed to a ``Monthly Unlimited'' bundle (10,000 FCFA) with your operator. After two weeks, the speed dropped drastically (throttling), and customer service has not resolved the issue after three complaints. Write a formal letter of complaint to the General Manager of the operator requesting a refund or an immediate technical fix. Use the standard layout of a formal letter.''
\end{exercice}

\begin{methode}[Méthode : Rédaction d'une lettre formelle de réclamation (Sujet 2)]
\textbf{Structure type :}
\begin{enumerate}
\item \textbf{Expéditeur} (Adresse, Email, Téléphone) — en haut à gauche.
\item \textbf{Destinataire} (Nom / Titre, Société, Adresse) — en haut à droite, sous l'expéditeur.
\item \textbf{Lieu et Date} (ex. \eng{Yaoundé, 24 October 2025}) — à droite.
\item \textbf{Objet} (Objet : \cle{Complaint} / \cle{Refund Request} — \cle{Monthly Unlimited Bundle}) — gras, souligné.
\item \textbf{Formule d'appel} (\eng{Dear Sir/Madam,} ou \eng{Dear General Manager,}).
\item \textbf{Corps du texte} (3 paragraphes) :
      \begin{itemize}
      \item \textbf{Paragraphe 1 :} Rappel du contrat (numéro d'abonné, date souscription, offre \cle{Monthly Unlimited 10,000 FCFA}).
      \item \textbf{Paragraphe 2 :} Exposé du problème (chute de débit / \cle{throttling} après 2 semaines, 3 réclamations infructueuses au service client, numéros de ticket si possible).
      \item \textbf{Paragraphe 3 :} Demande précise (\cle{full refund} ou \cle{immediate technical fix} / restauration du débit) + délai raisonnable (ex. \eng{within 7 working days}).
      \end{itemize}
\item \textbf{Formule de politesse} (\eng{Yours faithfully,} si \cle{Dear Sir/Madam} / \eng{Yours sincerely,} si nom connu).
\item \textbf{Signature} + Nom imprimé.
\end{enumerate}
\textbf{Lexique utile :} \eng{I am writing to complain about...}, \eng{Despite repeated complaints...}, \eng{I am entitled to...}, \eng{I look forward to your prompt response.}
\end{methode}

\begin{methode}[Méthode : Article d'opinion (Sujet 1)]
\textbf{Structure :}
\begin{enumerate}
\item \textbf{Titre accrocheur} (ex. \eng{Data Costs: A Barrier to Student Success?}).
\item \textbf{Introduction :} Accroche (statistique, anecdote), énoncé du sujet, thèse (Pour/Contre).
\item \textbf{Développement (2-3 arguments) :} 1 argument par paragraphe. Connecteurs (\eng{Firstly, Moreover, However, For instance}). Exemples camerounais (forfaits étudiants MTN/Orange, Wi-Fi campus, partenariats MINPOSTEL/Universités).
\item \textbf{Contre-argument (optionnel) :} \eng{Some argue that... but...}
\item \textbf{Conclusion :} Résumé + Appel à l'action / Proposition concrète (\eng{Partnerships with operators to zero-rate educational sites}).
\item \textbf{Signature} (Nom, Classe).
\end{enumerate}
\end{methode}

\begin{savaistu}[Le saviez-vous ? — Le « Last Mile » au Cameroun]
L'expression \cle{« last mile »} (dernier kilomètre) désigne la portion finale du réseau qui relie l'infrastructure principale (fibre optique, antenne relais) au domicile de l'abonné. Au Cameroun, bien que le backbone national (réseau dorsal) soit relativement dense grâce aux projets \cle{Camtel} et \cle{ACE} (Africa Coast to Europe), le « last mile » reste le goulet d'étranglement majeur en zone rurale. Le \cle{Universal Service Fund (FSU)}, financé par une taxe sur le chiffre d'affaires des opérateurs, vise précisément à subventionner ce « last mile » dans les zones non rentables commercialement (zones blanches). La \cle{Numérisation des services publics} (actes de naissance, impôts, inscription concours) via des \cle{USDD codes} accessibles sur téléphone basique (feature phone) est le levier principal pour justifier l'investissement auprès des populations rurales.
\end{savaistu}

\newpage

\coursec{Corrigés des exercices}

\courssub{Corrigé de l'exercice 1}

\begin{corrige}[Exercice 1 — Vrai ou Faux justifié (annonce « Student Connect »)]
\begin{enumerate}
\item \textbf{F.} The bundle costs 1,000 FCFA, not 1,500 FCFA. Justification : \eng{“For only 1,000 FCFA, you get 2 GB of high-speed internet.”}
\item \textbf{V.} The script clearly states : \eng{“you get 2 GB of high-speed internet.”}
\item \textbf{F.} The validity period is 7 days (\eng{“valid for 7 days”}), not 30 days.
\item \textbf{V.} The script states : \eng{“To subscribe, dial *157*10\# and press 1.”}
\item \textbf{F.} The offer is restricted : \eng{“this offer is exclusive to students with a valid student ID.”}
\end{enumerate}
\textbf{Points de vigilance :} en anglais, on écrit les prix avec une virgule pour les milliers (\eng{1,000 FCFA}) ; ne pas confondre \cle{valid for 7 days} (durée) et \cle{valid until} (date limite). La justification doit être une phrase complète, pas un simple mot relevé.
\end{corrige}

\courssub{Corrigé de l'exercice 2}

\begin{corrige}[Exercice 2 — QCM (dialogue « Family Pack »)]
\begin{enumerate}
\item \textbf{b)} 5,000 FCFA. L'agent dit : \eng{“The Family Pack costs 5,000 FCFA per month.”}
\item \textbf{b)} 20 GB. \eng{“It gives you 20 GB of data shared between up to 4 numbers.”}
\item \textbf{b)} 4. \eng{“shared between up to 4 numbers.”} Attention : \cle{up to} signifie « jusqu'à » (maximum), pas « exactement ».
\item \textbf{a)} *141*5\#. \eng{“To activate, dial *141*5\# and select option 2.”} Le code *157*10\# appartient au Script 1 (piège classique d'interférence entre les deux écoutes).
\item \textbf{b)} 5 minutes. \eng{“It takes a maximum of 5 minutes.”} Ne pas confondre avec le Script 1 où l'activation est \eng{instant}.
\end{enumerate}
\textbf{Points de vigilance :} dans un QCM d'écoute, les distracteurs reprennent souvent des chiffres entendus ailleurs dans le script ou dans un autre script ; il faut associer chaque chiffre à sa « fonction » (prix, volume, durée, délai).
\end{corrige}

\courssub{Corrigé de l'exercice 3}

\begin{corrige}[Exercice 3 — Tableau récapitulatif des offres Camtel]
\begin{center}
\begin{tabularx}{\linewidth}{|X|c|c|c|}\hline
\rowcolor{popblueL} \textbf{Offre} & \textbf{Prix mensuel (FCFA)} & \textbf{Vitesse / Données} & \textbf{Avantages inclus} \\ \hline
Basic  & 15,000 & 10 Mbps / 100 GB data cap & Free installation (with a 12-month contract) \\ \hline
Family & 25,000 & 50 Mbps / 500 GB & Free landline + free installation (12-month contract) \\ \hline
Premium & 45,000 & 100 Mbps / unlimited data & Free landline, free Wi-Fi 6 router, free installation (12-month contract) \\ \hline
\end{tabularx}
\end{center}
\textbf{Justifications relevées dans le script :} \eng{“The ‘Basic' offer costs 15,000 FCFA monthly. It provides 10 Mbps speed and 100 GB data cap.”} ; \eng{“The ‘Family' offer is 25,000 FCFA per month. Speed reaches 50 Mbps with 500 GB data. It includes a free landline.”} ; \eng{“The ‘Premium' offer costs 45,000 FCFA monthly. You get 100 Mbps speed, unlimited data, a free landline and a free Wi-Fi 6 router.”} ; \eng{“Installation is free for all packs if you sign a 12-month contract.”}
\textbf{Points de vigilance :} \cle{Mbps} se prononce « mégabits par seconde » et ne prend pas de -s variable en français technique courant ; \cle{data cap} = plafond de données (à ne pas traduire par « casquette » !) ; \cle{unlimited} est un adjectif invariable placé avant le nom : \eng{unlimited data}.
\end{corrige}

\courssub{Corrigé de l'exercice 4}

\begin{corrige}[Exercice 4 — Conjugaison et formes verbales]
\begin{enumerate}
\item \eng{Customers \textbf{subscribe} to the weekly bundle every Monday.} — Present Simple : habitude signalée par \cle{every Monday}.
\item \eng{The agent \textbf{has just activated} your line.} — Present Perfect : \cle{just} + résultat présent (\eng{You can start browsing now}).
\item \eng{If you \textbf{dial} *157\#, you \textbf{will see} the main menu.} — Conditionnel Type 1 : \cle{If} + Present Simple, puis \cle{will} + base verbale. Jamais \eng{*If you will dial...}.
\item \eng{\textbf{Do not forget} / \textbf{Don't forget} to renew your pack before it expires.} — Impératif négatif : \cle{Do not} + base verbale.
\item \eng{By the end of the year, the company \textbf{will launch} 5G in Buea and Bamenda.} — Future Simple : prédiction. (Le Future Perfect \eng{will have launched} serait possible mais n'était pas exigé.)
\item \eng{How long \textbf{have you had} this phone number?} — Present Perfect avec \cle{How long} : durée commencée dans le passé et qui continue. Attention : \cle{have} (possession durable), pas \eng{*have you bought} (action ponctuelle incompatible avec une durée).
\item \eng{The data balance \textbf{expires} in two days...} — Present Simple pour une échéance programmée ; \eng{\textbf{will expire}} (prédiction) est également accepté.
\item \eng{Please \textbf{enter} your PIN to confirm the transaction.} — Impératif affirmatif : base verbale seule.
\end{enumerate}
\textbf{Points de vigilance :} ne jamais mettre \cle{will} dans la subordonnée en \cle{if} ; l'impératif n'a pas de sujet exprimé ; \cle{forget} est suivi de \cle{to} + infinitif quand l'action n'est pas encore faite (\eng{Don't forget to renew}).
\end{corrige}

\courssub{Corrigé de l'exercice 5}

\begin{corrige}[Exercice 5 — Traduction (Français $\rightarrow$ Anglais)]
\begin{enumerate}
\item \eng{I want to \textbf{subscribe to} the monthly \textbf{unlimited bundle}.} — Ne pas oublier la préposition \cle{to} après \cle{subscribe}.
\item \eng{How much does the \textbf{5 GB top-up} cost?} — \cle{How much} car le prix est indénombrable ; on accepte aussi \eng{How much is a 5 GB recharge?}
\item \eng{The \textbf{validity} of this \textbf{bundle} is 30 days.} ou mieux : \eng{This bundle is \textbf{valid for} 30 days.}
\item \eng{\textbf{Check} your \textbf{balance} before \textbf{streaming} a video.} — Impératif + \cle{before} + verbe en -ing.
\item \eng{Network \textbf{coverage} is excellent in \textbf{downtown Yaoundé}.} — \cle{coverage} est indénombrable : pas d'article \eng{*a coverage}.
\item \eng{You must \textbf{renew} your \textbf{subscription} before the \textbf{expiry date}.} — \cle{must} + base verbale ; \cle{expiry date} (nom + nom), pas \eng{*expiration date of validity}.
\item \eng{This \textbf{family pack} offers \textbf{unlimited calls} to the same \textbf{network}.}
\item \eng{What is the \textbf{USSD code} to \textbf{subscribe to} this offer?} — Question directe : inversion \eng{What is...?}.
\end{enumerate}
\textbf{Points de vigilance :} « s'abonner à » = \cle{subscribe to} ; « illimité » se place avant le nom ; « Go » se note \cle{GB} en anglais ; « Mo » = \cle{MB}.
\end{corrige}

\courssub{Corrigé de l'exercice 6}

\begin{corrige}[Exercice 6 — Transformation de phrases]
\begin{enumerate}
\item \eng{Your balance \textbf{can be checked by dialing} *157\#.} — Passif avec modal : \cle{can be} + participe passé ; \cle{by} + -ing.
\item \eng{The activation \textbf{is not instant}.} — Négation du verbe \cle{be} : \cle{be} + \cle{not}.
\item \eng{\textbf{How much does 10 GB cost?}} — Question sur le prix : \cle{How much} + auxiliaire \cle{does} + base verbale.
\item \eng{He \textbf{has just subscribed} to the Premium pack.} — Present Perfect : \cle{has} + participe passé ; \cle{just} se place entre l'auxiliaire et le participe.
\item \eng{\textbf{Should you top up} today, you will get a bonus.} — Inversion formelle du Type 1 : \cle{Should} + sujet + base verbale (= \eng{If you should top up...}).
\item \eng{I \textbf{was told} the procedure by the agent.} / \eng{The procedure \textbf{was explained to me} by the agent.} — Attention : \cle{explain} ne prend pas de double complément ; la forme \eng{*I was explained the procedure} est à éviter en production soignée : préférer \cle{The procedure was explained to me}.
\item \eng{This offer \textbf{will expire} tomorrow.} — Future Simple : \cle{will} + base verbale.
\item \eng{He asked \textbf{if I had} a student ID.} — Discours indirect : question fermée introduite par \cle{if}, retour au présent de l'indicatif supprimé, concordance des temps (\cle{have} $\rightarrow$ \cle{had}), ordre sujet-verbe rétabli (pas d'inversion).
\end{enumerate}
\textbf{Points de vigilance :} au discours indirect, on ne garde ni l'auxiliaire \cle{do} ni l'inversion ; au passif, seul le verbe \cle{be} porte le temps, le participe passé est invariable.
\end{corrige}

\courssub{Corrigé de l'exercice 7}

\begin{corrige}[Exercice 7 — Texte à trous (scénario de souscription)]
\begin{enumerate}
\item \textbf{subscribe} — \eng{To subscribe to the new “Youth Boost” offer...} (infinitif de but, suivi de \cle{to}).
\item \textbf{dial} — \eng{customers must dial *145*20\#} (base verbale après \cle{must}).
\item \textbf{bundle} — \eng{The bundle costs 1,500 FCFA...}
\item \textbf{validity} — \eng{with a validity of 14 days.}
\item \textbf{unlimited} — \eng{Calls to the same network are unlimited.}
\item \textbf{confirmation} — \eng{you will receive a confirmation SMS.}
\item \textbf{activate} — \eng{To activate the service, simply restart your phone.}
\item \textbf{balance} — \eng{check your data balance at any time.}
\item \textbf{expire} — \eng{unused data will expire after the validity period} (base verbale après \cle{will}).
\item \textbf{coverage} — \eng{check the network coverage map.}
\end{enumerate}
Mots en trop : \cle{top up} et \cle{router}.
\textbf{Points de vigilance :} après un modal (\cle{must}, \cle{will}), toujours la base verbale ; distinguer \cle{activate} (verbe) et \cle{activation} (nom) — ici la structure \eng{To + verbe} impose le verbe.
\end{corrige}

\courssub{Corrigé de l'exercice 8}

\begin{corrige}[Exercice 8 — Appariements termes / définitions]
\begin{center}
\begin{tabularx}{\linewidth}{|c|X|}\hline
\rowcolor{popblueL} \textbf{Paire} & \textbf{Justification} \\ \hline
1 — B & A USSD code is a short code (e.g. *157\#) used to access mobile services. \\ \hline
2 — A & A top-up is an amount of money added to the account to buy data or credit. \\ \hline
3 — D & A data cap is the maximum amount of data allowed in a package. \\ \hline
4 — C & The expiry date is the date when the package stops working. \\ \hline
5 — E & A bundle is a package combining data, voice and/or SMS for a fixed price. \\ \hline
6 — F & Coverage is the geographical area where the network signal is available. \\ \hline
7 — G & The balance is the remaining amount of data or credit. \\ \hline
8 — H & Auto-renewal is the automatic renewal of the subscription at the end of the validity period. \\ \hline
\end{tabularx}
\end{center}
\textbf{Points de vigilance :} ne pas confondre \cle{balance} (solde restant) et \cle{top-up} (montant ajouté) ; \cle{expiry date} (nom) correspond au verbe \cle{expire}.
\end{corrige}

\courssub{Corrigé de l'exercice 9}

\begin{corrige}[Exercice 9 — Bac Blanc : compréhension et travail de langue]
\textbf{A. Comprehension Questions} (barème indicatif : 13 pts)
\begin{enumerate}
\item In Douala, 4G coverage is almost ubiquitous, so people can stream videos and make mobile money transactions seamlessly, whereas in the rural Centre Region network coverage is patchy and internet speed often crawls at 2G levels. (2 pts : 1 pt par idée contrastée)
\item The plan aims to extend the fibre optic backbone to 360 subdivision headquarters. (1 pt)
\item Private operators are offered tax breaks and universal service fund subsidies. (2 pts : 1 pt par incitation)
\item Because 500 FCFA represents a significant portion of a farmer's daily income. (1 pt)
\item Operators have introduced “social bundles”: weekly packages of 500 MB for 200 FCFA, accessible via USSD codes. (2 pts : 1 pt solution, 1 pt prix + volume)
\item Digital literacy remains a barrier: many elders do not know how to subscribe, renew or check their balance without assistance. (1 pt)
\item She no longer has to travel to Yaoundé: she tops up her phone at the centre and sells her cocoa to buyers in Douala via WhatsApp, so her income has doubled. (2 pts)
\item Content in local languages (Bassa, Ewondo, Fulfulde) and relevant e-government services (birth certificates, tax payment). (2 pts)
\end{enumerate}
\textbf{B. Language Work} (barème indicatif : 8 pts)
\begin{enumerate}
\item a) \textbf{ubiquitous} ; b) \textbf{subsidies} (ou \cle{tax breaks}) ; c) \textbf{barrier} ; d) \textbf{imperative} (ou \cle{crucial}). (2 pts : 0,5 pt par mot)
\item a) \eng{Social bundles \textbf{have been introduced} by operators.} (passif au Present Perfect) ; b) \eng{Many elders are unable \textbf{to subscribe}.} ; c) \eng{\textbf{Should the government build} more towers, coverage will improve.} (3 pts : 1 pt par transformation)
\item \eng{By the time the plan \textbf{achieves} its goals in 2025, thousands of kilometres of fibre \textbf{will have been laid}. Currently, the Universal Service Fund \textbf{finances} (ou \textbf{is financing}) several pilot projects. If the projects \textbf{are} successful, rural GDP \textbf{will rise} significantly.} (3 pts : 0,6 pt par verbe)
\end{enumerate}
\textbf{Points de vigilance :} après \cle{By the time} suivi d'une date future, on utilise le Present Simple dans la subordonnée et le Future Perfect dans la principale ; le passif de \cle{lay} exige le participe \cle{laid} (verbe irrégulier : lay — laid — laid) ; répondre par des phrases complètes et, autant que possible, avec ses propres mots (\emph{in your own words}).
\end{corrige}

\courssub{Corrigé de l'exercice 10}

\begin{corrige}[Exercice 10 — Traduction (Anglais $\rightarrow$ Français)]
« Pour l'abonné rural moyen, le défi ne tient pas seulement à la couverture, mais aussi au coût. Un forfait journalier de 1 Go coûte 500 FCFA — une part importante du revenu quotidien d'un cultivateur. Les opérateurs ont lancé des « forfaits sociaux » : des packs hebdomadaires de 500 Mo à 200 FCFA, accessibles au moyen de codes USSD comme *157*99\#. Pourtant, la culture numérique reste un obstacle. Beaucoup de personnes âgées ne savent pas comment s'abonner, renouveler leur forfait ou vérifier leur solde sans aide. »
\textbf{Points de vigilance :} \cle{affordability} se rend par « coût abordable / accessibilité financière », jamais par « abordabilité » en registre soigné ; \cle{digital literacy} = « culture numérique » (l'« illectronisme » désigne plutôt l'absence de cette culture) ; \cle{a farmer} = « un cultivateur / un paysan » selon le contexte ; conserver les sigles (FCFA, USSD, Go/Mo adaptés en Mo/Go français).
\end{corrige}

\courssub{Corrigé de l'exercice 11}

\begin{corrige}[Exercice 11 — Expression écrite (type Bac)]
\textbf{Barème indicatif (sur 20) :} pertinence et respect du sujet (4 pts) ; organisation et cohérence (4 pts) ; correction grammaticale (4 pts) ; richesse et exactitude du vocabulaire thématique (4 pts) ; orthographe et ponctuation (2 pts) ; respect du nombre de mots et de la mise en forme (2 pts).

\textbf{Proposition de rédaction modèle — Sujet 1 (article d'opinion) :}
\begin{textebox}[Model Answer — Opinion Article]
\textbf{Data Costs: A Barrier to Student Success?}

Every evening, thousands of Cameroonian students stop their online research because their data bundle has expired. In my view, mobile internet packages are indeed too expensive for students, and urgent solutions are needed.

Firstly, a standard 2 GB weekly bundle costs about 1,000 FCFA, which is a heavy burden for a student with no income. Moreover, online classes, digital libraries and examination registration now require a stable connection. For instance, during the GCE registration period, many candidates in Bamenda had to top up several times just to upload their documents.

Some argue that operators already offer student bundles. However, these offers are limited in volume and validity, and campus Wi-Fi remains rare.

To bridge this gap, the government and operators should partner to zero-rate educational websites and provide free Wi-Fi in every university. Affordable data is not a luxury: it is the key to equal opportunities for all students.

\emph{By A. Ngo Bell, Tle SH.}
\end{textebox}

\textbf{Proposition de rédaction modèle — Sujet 2 (lettre de réclamation) :}
\begin{textebox}[Model Answer — Formal Letter of Complaint]
P.O. Box 1234, Yaoundé\\
Tel: 6XX XX XX XX\\[0.3em]
The General Manager\\
{}[Operator] Cameroon\\
Douala\\[0.3em]
Yaoundé, 24 October 2025\\[0.3em]
\textbf{Subject: Complaint and Refund Request — “Monthly Unlimited” Bundle}\\[0.3em]
Dear Sir/Madam,\\[0.3em]
I am writing to complain about the “Monthly Unlimited” bundle I subscribed to on 1 October 2025 for 10,000 FCFA on the number 6XX XX XX XX.\\[0.3em]
After only two weeks, my connection speed dropped drastically, although the offer was advertised as unlimited. I have contacted your customer service three times, but the issue has not been resolved.\\[0.3em]
I therefore request a full refund or an immediate technical fix within seven working days. I look forward to your prompt response.\\[0.3em]
Yours faithfully,\\[0.3em]
A. Ngo Bell
\end{textebox}
\textbf{Points de vigilance :} respecter la mise en page formelle (adresses, date, objet, formules d'appel et de politesse appariées : \eng{Dear Sir/Madam} $\rightarrow$ \eng{Yours faithfully}) ; employer le Present Perfect pour les réclamations déjà faites (\eng{I have contacted...}) ; réemployer le lexique de la leçon (\cle{subscribe}, \cle{bundle}, \cle{unlimited}, \cle{refund}, \cle{top up}) ; éviter les calques du français (\eng{*I am writing you} $\rightarrow$ \eng{I am writing to you}).
\end{corrige}

\newpage

\coursec{Fiche bilan}

\begin{popfigure}
\begin{tikzpicture}[
  mindmap,
  grow cyclic,
  every node/.style={concept, execute at begin node=\hskip0pt},
  concept color=popblue,
  level 1/.append style={level distance=4.5cm, sibling angle=360/7},
  level 2/.append style={level distance=3cm, sibling angle=360/4},
  telecom/.style={concept color=poporange},
  strategy/.style={concept color=popgreen},
  vocab/.style={concept color=poppurple},
  grammar/.style={concept color=poppink},
  questions/.style={concept color=popgold},
  production/.style={concept color=popdark},
  context/.style={concept color=popmuted}
]
\node[concept, text width=2.8cm, align=center, font=\bfseries\large] {Lesson 49\\Subscribing to\\Service Packages}
  child[telecom] { node[concept, text width=2.2cm, align=center, font=\bfseries] {Theme:\\Telecom \& Internet\\Packages} }
  child[strategy] { node[concept, text width=2.2cm, align=center, font=\bfseries] {Listening\\Strategies\\Keywords, Numbers,\\Proper Nouns} }
  child[vocab] { node[concept, text width=2.2cm, align=center, font=\bfseries] {Key Vocabulary\\Plan, Data, Bundle,\\Subscription, Fee} }
  child[grammar] { node[concept, text width=2.2cm, align=center, font=\bfseries] {Grammar Tools\\Present Perfect,\\Modals, Conditionals} }
  child[questions] { node[concept, text width=2.2cm, align=center, font=\bfseries] {Comprehension\\Tasks\\True/False, MCQ,\\Short Answers} }
  child[production] { node[concept, text width=2.2cm, align=center, font=\bfseries] {Guided Production\\Role-play, Summary,\\Email writing} }
  child[context] { node[concept, text width=2.2cm, align=center, font=\bfseries] {Cameroon Context\\MTN, Orange,\\Camtel, Bendskin\\WiFi} };
\end{tikzpicture}
\legende{Carte mentale : Lesson 49 — Subscribing to Service Packages (Tle SH)}
\end{popfigure}

\begin{bilan}

\end{bilan}

\courssub{Lexique clé — Telecommunications \& Service Packages}

\begin{center}
\begin{tabularx}{\linewidth}{|X|X|X|}
\hline
\rowcolor{popblueL}
\textbf{English Term} & \textbf{Français} & \textbf{Context / Collocation} \\
\hline
\cle{A subscription plan} / \cle{A package} & Un forfait / Un abonnement & \eng{Choose a monthly package.} (« Choisis un forfait mensuel. ») \\
\cle{Prepaid} / \cle{Postpaid} & Prépayé / Postpayé (facturé) & \eng{Prepaid SIM cards are popular.} \\
\cle{Data allowance} / \cle{Data cap} & Volume de données (plafond) & \eng{Unlimited data allowance.} \\
\cle{To subscribe / To sign up for} & S'abonner / Souscrire & \eng{I want to subscribe to the 4G plan.} \\
\cle{A bundle} & Un pack groupé (voix + data + SMS) & \eng{A family bundle saves money.} \\
\cle{Rollover data} & Report de données (non utilisées) & \eng{Unused data rolls over to next month.} \\
\cle{Roaming charges} & Frais d'itinérance & \eng{Avoid roaming charges abroad.} \\
\cle{A contract} / \cle{Terms and conditions} & Un contrat / Conditions générales & \eng{Read the terms before signing.} \\
\cle{Activation fee} / \cle{Set-up fee} & Frais d'activation / d'installation & \eng{There is a one-time activation fee.} \\
\cle{Customer care} / \cle{Helpline} & Service client / Numéro vert & \eng{Call customer care for help.} \\
\cle{Network coverage} & Couverture réseau & \eng{Check coverage in your area.} \\
\cle{To top up} / \cle{To recharge} & Recharger (crédit) & \eng{I need to top up my credit.} \\
\hline
\end{tabularx}
\end{center}

\courssub{Points de grammaire indispensables (Outils pour l'écoute et la production)}

\begin{center}
\begin{tabularx}{\linewidth}{|l|X|X|}
\hline
\rowcolor{popblueL}
\textbf{Notion} & \textbf{Règle / Forme clé} & \textbf{Exemple (Écoute / Production)} \\
\hline
\cle{Present Perfect} & \eng{have/has + V-ed/Participe passé}\newline Expérience, changement d'état, résultat présent. & \eng{“I \textbf{have subscribed} to the new plan.”} (« Je me suis abonné… ») \\
\cle{Modals : Politeness \& Obligation} & \eng{Could/Would} (politesse), \eng{Must/Have to} (obligation), \eng{Should} (conseil). & \eng{“\textbf{Could} you explain the fees?” / “You \textbf{must} show your ID.”} \\
\cle{First Conditional (Real Future)} & \eng{If + Present Simple, ... will + BV}\newline Offres, promesses, conséquences réelles. & \eng{“If you \textbf{subscribe} today, you \textbf{will get} a discount.”} \\
\cle{Comparatives / Superlatives} & \eng{-er/more ... than}, \eng{the -est/most}\newline Comparer les forfaits. & \eng{“This plan is \textbf{cheaper than} the other one.” / “The \textbf{best} value.”} \\
\cle{Passive Voice (Present/Past)} & \eng{is/are + V-ed}, \eng{was/were + V-ed}\newline Focus sur l'action/service, pas l'agent. & \eng{“The SIM card \textbf{is activated} instantly.” / “Fees \textbf{were charged}.”} \\
\hline
\end{tabularx}
\end{center}

\courssub{Méthode express — Réussir l'écoute « Service Packages »}

\begin{enumerate}
\item \textbf{Anticiper (30 sec) :} Lire le titre, les questions, les images. Activer le lexique \cle{plan, fee, data, contract, coverage}.
\item \textbf{Écoute 1 (Global) :} Identifier le \textbf{thème} (quel opérateur ? quel type de forfait ?), les \textbf{interlocuteurs} (client / agent), le \textbf{but} (plainte, information, souscription).
\item \textbf{Écoute 2 (Détails/Chiffres) :} Noter \textbf{chiffres} (prix, Go, minutes, durée contrat), \textbf{noms propres} (noms de forfaits : “MaxiPlus”, “Youth Pack”), \textbf{mots-clés de liaison} (\eng{however, therefore, included, excluded}).
\item \textbf{Écoute 3 (Vérification) :} Confirmer Vrai/Faux, choisir QCM, compléter phrases. Attention aux \textbf{négations} (\eng{not, never, without, unless}) et aux \textbf{faux amis} (\eng{actually = en fait, pas actuellement}).
\item \textbf{Transfert :} Réutiliser 5--7 mots notés pour résumer oralement ou écrire un SMS/email de confirmation.
\end{enumerate}



\begin{aretenir}[Checklist avant le Bac — Lesson 49]
$\square$ Je sais épeler et prononcer les termes clés : \eng{subscription, allowance, bundle, roaming, coverage, top-up}.\par
$\square$ Je distingue \eng{prepaid} (prépayé) de \eng{postpaid} (postpayé) et je sais expliquer la différence.\par
$\square$ Je repère les \textbf{chiffres} (prix, Go, minutes) et les \textbf{noms propres} (noms d'offres) à l'écoute.\par
$\square$ J'utilise le \textbf{Present Perfect} pour annoncer un changement récent : \eng{I've just subscribed...}\par
$\square$ J'utilise les \textbf{modaux} pour la politesse (\eng{Could/Would}) et l'obligation (\eng{Must/Have to}).\par
$\square$ Je construis une phrase au \textbf{First Conditional} pour une offre : \eng{If you sign up, you will save...}\par
$\square$ Je sais comparer deux forfaits avec les \textbf{comparatifs/superlatifs}.\par
$\square$ Je reconnais la \textbf{voix passive} pour décrire les services : \eng{The service is included.}\par
$\square$ Je sais rédiger un \textbf{email/SMS formel} de confirmation ou de réclamation (structure : objet, faits, demande).\par
$\square$ Je connais les 3 grands opérateurs camerounais (\eng{MTN, Orange, Camtel}) et le vocabulaire local (\eng{credit, bundle, data}).\par
$\square$ Je gère mon temps : lecture questions (1 min), écoute 1 (global), écoute 2 (détails), relecture (30 s).\par
$\square$ Je ne panique pas si je rate un mot : j'infère le sens par le contexte (mots transparents, ton, logique).
\end{aretenir}

\findecours

\end{document}
